\documentclass[11pt]{article}

\usepackage[margin=1.1in]{geometry}

\usepackage{amssymb,mathtools,natbib}

% Anchored formal environments for causalsmith papers.
% First mandatory argument = obj_id (machine anchor joining the paper to the
% Lean crosswalk; invisible in the rendered PDF). Optional argument = name.
% Each environment also defines \label{obj:<obj_id>} so prose can use
% \cref{obj:T-1}; cleveref derives the printed kind and number from the target.
\usepackage{amsmath}
\usepackage{mathrsfs}
\usepackage{amsthm}
\usepackage{booktabs}
% Long displays may break across pages; let TeX stretch a little harder before an
% overfull \hbox so wide formulas/tables are less likely to run off the page.
\allowdisplaybreaks
\setlength{\emergencystretch}{3em}
\newtheorem{theorem}{Theorem}
\newtheorem{lemma}{Lemma}
\newtheorem{assumption}{Assumption}
\newtheorem{definition}{Definition}
\newtheorem{proposition}{Proposition}
\newtheorem{remark}{Remark}
\newtheorem{citedresult}{Cited result}
% The amsthm counter is `algorithmbox` (not `algorithm`) so a later `\usepackage{algorithm}` cannot clash.
\newcounter{algorithmbox}
\renewcommand{\thealgorithmbox}{\arabic{algorithmbox}}
\NewDocumentEnvironment{theoremv}{m o s}{\begin{theorem}\label{obj:#1}{\normalfont[\texttt{\detokenize{#1}}]\IfValueT{#2}{\ (#2)}\IfBooleanT{#3}{\verificationfootnotemark}\ }}{\end{theorem}}
\NewDocumentEnvironment{lemmav}{m o s}{\begin{lemma}\label{obj:#1}{\normalfont[\texttt{\detokenize{#1}}]\IfValueT{#2}{\ (#2)}\IfBooleanT{#3}{\verificationfootnotemark}\ }}{\end{lemma}}
\NewDocumentEnvironment{assumptionv}{m o s}{\begin{assumption}\label{obj:#1}{\normalfont[\texttt{\detokenize{#1}}]\IfValueT{#2}{\ (#2)}\IfBooleanT{#3}{\verificationfootnotemark}\ }}{\end{assumption}}
\NewDocumentEnvironment{definitionv}{m o s}{\begin{definition}\label{obj:#1}{\normalfont[\texttt{\detokenize{#1}}]\IfValueT{#2}{\ (#2)}\IfBooleanT{#3}{\verificationfootnotemark}\ }}{\end{definition}}
% A `propositionv` is a proved result tiered as a Proposition (theorem-grade, machine-verified).
\NewDocumentEnvironment{propositionv}{m o s}{\begin{proposition}\label{obj:#1}{\normalfont[\texttt{\detokenize{#1}}]\IfValueT{#2}{\ (#2)}\IfBooleanT{#3}{\verificationfootnotemark}\ }}{\end{proposition}}
% A `remarkv` holds NON-load-bearing interpretive / open-question content (e.g. a causalsmith `oeq:`
% open question). It is neither proved nor assumed; no theorem may depend on it.
\NewDocumentEnvironment{remarkv}{m o s}{\begin{remark}\label{obj:#1}{\normalfont[\texttt{\detokenize{#1}}]\IfValueT{#2}{\ (#2)}\IfBooleanT{#3}{\verificationfootnotemark}\ }}{\end{remark}}
% An `algorithmv` is a DEFINITION-kind object presented as a boxed, numbered-step procedure (an
% estimator / selection rule built in stages). Same anchor contract as `definitionv`; its body is
% ordinary LaTeX whose steps are `\begin{enumerate}` items, so tex2html needs no extra machinery.
% Presented in the CONVENTIONAL algorithm style readers expect from `algorithm`+`algorithmic`:
% a rule above, an "Algorithm N (Title)" caption line, a rule under the caption, the numbered
% steps, and a closing rule. It is NOT a float — the anchor contract requires the object to stay
% where the outline places it, and floats drift. The obj-id tag is suppressed here (it is
% bookkeeping, and a caption line carrying it reads as debug output in a finished paper).
\NewDocumentEnvironment{algorithmv}{m o s}%
  {\par\addvspace{\medskipamount}\noindent\rule{\linewidth}{0.8pt}\par\nobreak\vspace{2pt}%
   \refstepcounter{algorithmbox}\label{obj:#1}%
   \noindent\textbf{Algorithm \thealgorithmbox}\IfValueT{#2}{ \textnormal{(#2)}}%
   \IfBooleanT{#3}{\verificationfootnotemark}\par\nobreak\vspace{2pt}%
   \noindent\rule{\linewidth}{0.4pt}\par\nobreak\vspace{2pt}\normalfont}%
  {\par\nobreak\vspace{2pt}\noindent\rule{\linewidth}{0.8pt}\par\addvspace{\medskipamount}}
% A `citedv` is an IMPORTED external result the paper relies on but does NOT prove
% (a causalsmith `gate_class:"cited"` node). It is machine-checked against the cited
% source at F2.5, never proved here; the fixed provenance note makes that explicit
% so the environment can never be read as a contribution of this paper. The \citep
% attribution is supplied by the surrounding prose (P2), keyed to references.bib.
\NewDocumentEnvironment{citedv}{m o s}{\begin{citedresult}\label{obj:#1}{\normalfont[\texttt{\detokenize{#1}}]\IfValueT{#2}{\ (#2)}\IfBooleanT{#3}{\verificationfootnotemark}\ \textit{(imported result; machine-checked against the cited source, not proved here.)}\ }}{\end{citedresult}}
% \leanref{obj_id}{display text}: an inline first-mention link to a from-note object's
% verified Lean code. In the PDF it renders the display text plainly (the Lean link is a
% web-only affordance); tex2html turns it into a drawer-opening span.
\newcommand{\leanref}[2]{#2}
% For a theorem/lemma whose Lean declaration accepts a source-matched published proposition as an
% input, the starred anchored environment puts the mark in its heading; the generated text remains
% outside the frozen mathematical environment. Keep the combined macro for older paper bundles.
\newcommand{\verificationfootnotemark}{\footnotemark}
\newcommand{\verificationfootnotetext}[1]{\footnotetext{#1}}
\newcommand{\verificationfootnote}[1]{\verificationfootnotemark\verificationfootnotetext{#1}}
% xcolor before hyperref (hyperref would otherwise pull in plain `color` for colorlinks, and a
% later xcolor load clashes). The page-1 identity stamp at the end of this file needs a grey.
\usepackage{xcolor}
\usepackage{graphicx}
\usepackage{eso-pic}
% hyperref follows natbib and the environment macros; cleveref must follow hyperref.
\usepackage[colorlinks=true,linkcolor=blue,citecolor=blue,urlcolor=blue]{hyperref}
\usepackage[capitalise,nameinlink,noabbrev]{cleveref}
% Pin the names explicitly: labels are placed inside the underlying amsthm environments, and these
% declarations make the PDF contract independent of cleveref's theorem-name inference. House style:
% reference names always print with a capital first letter ("Theorem 1", "Definition 2"), in any
% sentence position — hence `capitalise` above and capitalized \crefname forms below.
\crefname{theorem}{Theorem}{Theorems}
\Crefname{theorem}{Theorem}{Theorems}
\crefname{lemma}{Lemma}{Lemmas}
\Crefname{lemma}{Lemma}{Lemmas}
\crefname{assumption}{Assumption}{Assumptions}
\Crefname{assumption}{Assumption}{Assumptions}
\crefname{definition}{Definition}{Definitions}
\Crefname{definition}{Definition}{Definitions}
\crefname{proposition}{Proposition}{Propositions}
\Crefname{proposition}{Proposition}{Propositions}
\crefname{remark}{Remark}{Remarks}
\Crefname{remark}{Remark}{Remarks}
\crefname{citedresult}{Cited result}{Cited results}
\Crefname{citedresult}{Cited result}{Cited results}
\crefname{algorithmbox}{Algorithm}{Algorithms}
\Crefname{algorithmbox}{Algorithm}{Algorithms}
% ---------------------------------------------------------------------------
% Page-1 identity stamp (arXiv style) and the pinned title-block date.
%
% Split of responsibility: THIS file owns the FORMAT (placement, font, colour, punctuation),
% the generated `paper_stamp.tex` owns only the VALUES (number+version, area, revised date,
% DOI, link target). P4 refreshes this template on every emit, so a layout fix reaches every
% bundle from a recompile; and because `paper_stamp.tex` is not one of the P2 assembly
% digest's authored inputs, a DOI that arrives after publication needs only that one small
% file rewritten plus a recompile — never a redraft, never a reassembly.
%
% Defaults first, so a bundle WITHOUT a stamp file compiles exactly as it always did:
% `\cspaperdate` falls back to `\today` and no stamp is drawn.
\providecommand*{\cspaperdate}{\today}  % the \date{} target
\providecommand*{\csstampid}{}          % e.g. CSWP-2026-008v3 (empty ⇒ no stamp at all)
\providecommand*{\csstamparea}{}        % e.g. Stat
\providecommand*{\csstampdate}{}        % e.g. 27 Aug 2026
\providecommand*{\csstampdoi}{}         % e.g. 10.5281/zenodo.123 (empty ⇒ DOI part omitted)
\providecommand*{\csstamplink}{}        % href target (DOI url, else the paper's page; may be empty)
\IfFileExists{paper_stamp.tex}{% GENERATED by CausalSmith P4 — paper identity stamp values. Do not hand-edit.
%
% Field order on the stamp (owned by paper_macros.tex):
%   <number>v<version> · doi:<doi> · [<Area>] <date>   — the DOI is never the last token, so a
%   text extractor cannot run it into whatever follows. Without a DOI that part is omitted.
% Format lives in paper_macros.tex; only the values are here, and this file is NOT one of
% the P2 assembly digest's authored inputs. A DOI arriving after publication therefore needs
% this file rewritten plus a `latexmk` recompile — never a redraft or a reassembly.
\renewcommand*{\csstampid}{CSWP-2026-040v3}
\renewcommand*{\csstamparea}{Experimentation}
\renewcommand*{\csstampdate}{4 Oct 2026}
\renewcommand*{\csstampdoi}{}
\renewcommand*{\csstamplink}{https://causalsmith.org/papers/exp_bivariate_sobolev_design_capacity_v1}
\renewcommand*{\cspaperdate}{4 October 2026}
}{}
\makeatletter
% Field order is deliberate: the DOI is NEVER the last token on the line. A trailing identifier
% runs into whatever a text extractor emits next, which makes it unsafe to match mechanically
% (the Zenodo stamp matcher reads exactly this line); the date is a safe terminator.
%   <number>v<version> · doi:<doi> · [<Area>] <date>      — with a DOI
%   <number>v<version> · [<Area>] <date>                  — without
\newcommand{\cs@stampsep}{\,\textperiodcentered\,}
\newcommand{\cs@stamptext}{%
  \csstampid
  \ifx\csstampdoi\@empty\else\cs@stampsep doi:\csstampdoi\fi
  \ifx\csstamparea\@empty\else\cs@stampsep[\csstamparea]\fi
  \ifx\csstampdate\@empty\else\quad\csstampdate\fi
}
\ifx\csstampid\@empty\else
  \definecolor{csstampgrey}{gray}{0.45}
  % Background shipout material: it occupies no space, so the text block and the \thanks
  % footnote are byte-identical to an unstamped compile. The starred form applies to the
  % NEXT page shipped only — i.e. page 1. Placed in the left margin (the text block starts
  % at 1.1in), reading bottom-to-top, in a Times-like face at 8pt.
  \AddToShipoutPictureBG*{%
    \AtPageLowerLeft{%
      \put(\LenToUnit{0.42in},\LenToUnit{1.1in}){%
        \rotatebox{90}{%
          \fontfamily{ptm}\fontsize{8}{10}\selectfont
          \ifx\csstamplink\@empty
            \textcolor{csstampgrey}{\cs@stamptext}%
          \else
            \expandafter\href\expandafter{\csstamplink}{\textcolor{csstampgrey}{\cs@stamptext}}%
          \fi
        }%
      }%
    }%
  }
\fi
\makeatother


\title{Minimax excess variance with smooth pairwise interactions in randomized experiments}

\author{CausalSmith\thanks{Machine-generated by the CausalSmith research pipeline.}}

\date{\cspaperdate}

\begin{document}

\maketitle

\begin{abstract}
We characterize the finite-sample minimax excess variance of Horvitz--Thompson estimation under covariate-dependent experimental assignment. Covariates are independent uniform draws on the unit cube, and the conditional potential-outcome midpoint mean is centered, has an exact main-and-pair-effect decomposition, and satisfies a pooled nonperiodic Sobolev restriction-norm budget. For every sample size \(n\geq2\), covariate dimension \(d\geq2\), and smoothness \(0<s\leq1\), the minimax scaled excess variance is comparable to \(\min\{1,(B/n)^s\}\), where \(B=\binom d2\) is the number of coordinate pairs, with positive comparison constants uniform over this parameter range. A single fair, outcome-independent assignment law, independent of smoothness, attains this scale on all original units. An independent coefficient-sign prior supplies the matching lower bound against every admissible design. The characterization is logarithm-free, includes the endpoint \(s=1\), and yields vanishing minimax scaled excess exactly when \(d=o(\sqrt n)\) for each fixed smoothness. The causal interpretation applies to arbitrary square-integrable potential outcomes with the specified uniform-covariate prognostic class.
\end{abstract}

\section{Introduction}

This paper establishes a finite-sample capacity characterization for covariate balancing in randomized experiments. Under independent uniform covariates, we consider centered prognostic functions consisting exactly of canonical main effects and pairwise interactions, with a common squared nonperiodic Sobolev restriction-norm budget pooled across components. Writing \(n\) for the number of experimental units, \(d\) for covariate dimension, and \(s\) for component smoothness, the best achievable scaled excess variance of the unit-weight Horvitz--Thompson estimator for this class is of order
\[
\min\left\{1,\left(\frac{\binom d2}{n}\right)^s\right\},
\qquad n,d\geq2,\quad 0<s\leq1.
\]
One fair, outcome-independent assignment procedure attains this scale simultaneously across the stated smoothness range. The matching converse applies to every admissible covariate-dependent assignment law.

The statistical question is how much precision assignment can deliver when the outcome regression has many potentially relevant interactions. Pairwise structure associates each component with at most two covariates, while allowing the number of components to grow quadratically with ambient dimension. A pooled regularity budget bounds their total squared component norm and permits its allocation across coordinates and pairs to vary. Our characterization identifies \(B=\binom d2\), the coordinate-pair count, relative to sample size as the quantity governing worst-case excess variance. Below saturation, \(B<n\), the rate is \((B/n)^s\); at and above saturation, it is constant order. For every fixed smoothness, vanishing minimax scaled excess is equivalent to \(d=o(\sqrt n)\), including along arbitrary integer dimension sequences.

The design problem fixes the estimator and specifies both the available information and the population comparison. An admissible design is a Borel assignment kernel using the complete original covariate array, with conditional treatment probability one-half for each unit and assignment independent of potential outcomes given covariates. Every original unit receives a treatment sign and enters the Horvitz--Thompson estimator at its original weight. Performance is evaluated by expected squared prognostic imbalance, averaging over covariate sampling and assignment. The supremum over prognostic functions lies outside that expectation, so each function is fixed across sample realizations. These choices define the minimax signing risk \(R_s(n,d)\), the scaled population imbalance criterion in \cref{obj:def:criterion}, over the function and design classes in \cref{obj:def:sobolev-class,obj:def:design-class}.

The causal interpretation follows from the excess-variance identity of \citet[Proposition 2.3, equation (2.4), under Assumption 2.1 and Definition 2.2]{Cytrynbaum2026Limits}. For independent original units, admissible assignment, and finite potential-outcome second moments, scaled Horvitz--Thompson variance above the law-specific benchmark equals scaled expected squared imbalance of the conditional potential-outcome midpoint mean. Every prognostic function in our class is realized by a Gaussian outcome completion. Consequently, the minimax excess over all uniform-covariate outcome laws with square-integrable potential outcomes and the specified prognosis equals \(R_s(n,d)\), as established in \cref{obj:lem:published-prognostic-capacity,obj:thm:capacity-answer}. The Gaussian subclass realizes the same capacity and provides a concrete calibration with scaled variance benchmark \(4.01\).

The closest rate comparison is \citet[Theorem B.8(ii), Lemma B.9, and Theorem 4.9]{Cytrynbaum2026Limits}. His correctly specified pairwise model admits a structured Gram--Schmidt Walk upper bound of order \((d^2\log n/n)^{s\wedge1}\), while the full-dimensional smooth-class lower bound specializes to exponent \(s\) in two dimensions. For our independent uniform sample and centered pooled order-two Sobolev class, \cref{obj:thm:capacity-answer,obj:thm:log-free-bivariate-capacity} matches upper and lower bounds with explicit interaction-count dependence and a logarithm-free rate. The comparison includes the frequency-rich regime \(B/n\to0\) and the endpoint \(s=1\). An existing random-matrix route supplies one-frequency saturation through \citet[Theorem 1.2(2a)]{KoganNandyHuang2024Extremal}; our characterization covers both that regime and the subcritical frequency-rich comparison at finite sample sizes.

The construction connects minimax experimental design \citep[Sections 2.2--2.4]{Kallus2018} to randomized balancing. In the subcritical regime, independent coordinate reflections and a common random shift supply ordered real spectral features. A finite boundary-rounding procedure protects an initial segment of these features while preserving coordinate means; the protected segment contracts as assignment coordinates reach their boundaries. The feature ordering is independent of smoothness. Combining its row-imbalance guarantees with the reflected nonperiodic Sobolev budget yields simultaneous attainment. In the saturated regime, the procedure uses independent fair signs. The complete assignment law and its guarantees are specified in \cref{obj:def:capacity-handle,obj:def:ordered-boundary-rounding,obj:thm:log-free-bivariate-capacity} under exact-real arithmetic conventions.

The converse uses a finite pair-cosine family with independent coefficient signs drawn before the covariate sample. Every realization satisfies the pooled Sobolev budget. A sample-size-dependent frequency cutoff supplies enough feature directions for a signing lower bound, including directions from pairs sharing an original coordinate. Averaging over this legal prior preserves the fixed-function population comparison and yields the full-design lower bound in \cref{obj:thm:full-design-lower}. Together, the two comparisons have constants bounded uniformly over \(0<s\leq1\). The uniform subclass also transfers the converse and the necessary dimension condition to the larger bounded-density outcome classes specified in \cref{obj:lem:published-prognostic-capacity}.

The main capacity results are stated in \cref{obj:thm:capacity-answer,obj:thm:log-free-bivariate-capacity}.

The paper proceeds through related work, setup and assumptions, main results, the smoothness-independent assignment procedure, and discussion and extensions. The appendices develop causal variance transfer, spectral and rounding bounds, the independent-prior converse, and proof assembly, followed by proofs of the main results.


\section{Related work}

The closest econometric comparison is \citet[Proposition 2.3, equation (2.4), Theorem B.8(ii), Lemma B.9, and Theorem 4.9]{Cytrynbaum2026Limits}. His excess-variance identity expresses the design-dependent part of Horvitz--Thompson variance as scaled expected squared prognostic imbalance, the criterion used in \cref{obj:def:criterion}. Writing \leanref{sym:n}{\ensuremath{n}} for sample size, \leanref{sym:d}{\ensuremath{d}} for covariate dimension, and \leanref{sym:s}{\ensuremath{s}} for smoothness, his correctly specified pairwise model admits a structured Gram--Schmidt Walk (GSW) upper bound of order \((d^2\log n/n)^{s\wedge1}\), supported by the pooled Mat\'ern approximation calculation. His lower bound \(n^{-2s/d}\) concerns a full-dimensional smooth class; its bivariate specialization has exponent \(s\). For independent uniform covariates and the centered order-two pooled Sobolev class in \cref{obj:def:sobolev-class}, the present characterization matches upper and lower bounds across sample size and dimension at order \(\min\{1,(B/n)^s\}\), where \leanref{sym:B}{\ensuremath{B=\binom{d}{2}}} is the number of coordinate pairs, throughout the \leanref{S-1}{parameter range \(n,d\geq2\) and \(0<s\leq1\)}. The comparison therefore concerns the logarithmic factor and the explicit interaction-count dependence for this specified class; \cref{obj:prop:published-class-inclusion} supplies the class-transfer relation.

A precise comparison with random-matrix theory identifies an existing route to saturation. Applied to independent, centered, unit-variance cosine coordinates under uniform sampling, \citet[Theorem 1.2(2a)]{KoganNandyHuang2024Extremal} implies the following spectral limit. Let \(G\) denote the sample Gram matrix of the pair-product features using one cosine frequency per coordinate, normalized by \(B\). As \(n,d\) tend to infinity with \(n/B\to\rho\in(0,1)\), where \(\rho\) is the limiting sample-to-feature ratio, its smallest eigenvalue converges almost surely to \((1-\sqrt{\rho})^2\). This positive spectral edge yields a constant-order signing lower bound for the normalized one-frequency submodel. The additional scope of \cref{obj:thm:capacity-answer,obj:thm:log-free-bivariate-capacity} is the matched finite-sample characterization, including the frequency-rich regime \(B/n\to0\), together with attainment by the same assignment procedure for every smoothness in the stated range. The independent coefficient prior in \cref{obj:def:cosine-prior} supplies the growing frequency family for that subcritical comparison.

The criterion also belongs to the minimax experimental-design tradition. \citet[Sections 2.2--2.4]{Kallus2018} connects restrictions on conditional outcome functions to worst-case imbalance and optimal randomized designs, while \citet[Sections 3--5]{Kallus2020} studies mixed assignment distributions and randomization constraints. Here, \cref{obj:def:design-class,obj:def:criterion} specifies fair, outcome-independent assignment laws and places the worst-function comparison outside the expectation over covariates and assignment. The lower bound in \cref{obj:thm:full-design-lower} applies to this entire admissible design class. Its role is to characterize the best achievable excess variance for the fixed Horvitz--Thompson criterion under the pooled smoothness restriction.

The attaining procedure connects this statistical objective to randomized balancing. The GSW design of \citet[Lemma 6.1 and Theorem 6.3]{HarshawSavjeSpielmanZhang2024GSW} combines assignment fairness with covariance control for augmented covariate features. The Edge-Walk method of \citet[Section 2.1]{LovettMeka2015EdgeWalk} develops constrained random walks for partial coloring. The procedure in \cref{obj:def:ordered-boundary-rounding} uses finite randomized boundary moves with an ordered family of balancing constraints. Its statistical application in \cref{obj:def:capacity-handle} uses a feature ordering independent of smoothness, allowing one assignment law to attain the comparison scale simultaneously across the stated Sobolev classes.

\citet[Lemma 4.2 and Proposition 5.1]{Chen2026} studies nonlinear feature balance through a normalized feature Gram matrix \(K\). For \(\phi\in(0,1)\), the design has assignment covariance bounded by \(\{\phi I+(1-\phi)K\}^{-1}\), and a fixed-schedule ridge approximation bound controls the actual Horvitz--Thompson variance. That comparison is conditional on the realized schedule and a chosen nonlinear feature map. The present result instead gives a population minimax capacity matched over an explicit pooled Sobolev class, with a single smoothness-independent construction and a converse for every admissible design.

Discrepancy-based integration provides a related use of randomized balancing with distinct sampling budgets and error criteria. \citet[Theorem 1.1 and Section 3.1]{BansalJiang2024QMC} starts from \(n^2\) independent uniform points and partitions them into \(n\) sets of \(n\) points. For each retained set, the mean squared integration error is bounded by \(\widetilde O_d(\sigma_{\mathsf{SO}}(f)^2/n^2)\), where \(f\) is the integrand and \(\sigma_{\mathsf{SO}}(f)\) is its smoothed-out variation; the displayed order suppresses logarithmic factors. \citet[Theorem 1.1, Section 2.2, and Theorem 3.3]{ChenJiangKirk2025HighDimensionalQMC} retains this quadratic input pool and develops bounds for weighted Fourier classes and low superposition dimension, with explicit coordinate-weight and dimension dependence. \citet[Theorem 1.2]{EzeunalaJhaJiang2026QMCII} obtains the smoothed-out-variation mean squared error bound while retaining \(n\) points from \((1+\varepsilon)n\) independent uniform inputs with high probability, for fixed \(\varepsilon\in(0,1)\); its constants and logarithmic factors depend on dimension and \(\varepsilon\).

Other thinning results emphasize different forms of approximation. The online procedure of \citet[Theorems 1--2]{Dwivedi2019} retains at least one point from every two consecutive uniform inputs and controls unnormalized counting discrepancy over axis-parallel rectangles. Kernel thinning in \citet[Theorem 1 and Corollary 1]{Dwivedi2024} compresses an input approximation into a smaller subset and bounds maximum mean discrepancy, equivalently worst-case integration error over the associated reproducing kernel Hilbert-space unit ball, under the stated regularity and input-radius conditions. In the present experiment, all \(n\) original units receive treatment signs, and performance is measured by expected squared prognostic imbalance over a pooled nonperiodic Sobolev class. These sampling, function-class, and loss specifications make the comparison with integration results precise.

Finally, \citet[Theorems 2.2--2.3 and 3.1]{SudijonoDobribanTchetgen2026Minimax} studies joint minimax choice of design and estimator for finite-population treatment-effect estimation with binary or bounded potential outcomes. For binary outcomes, Bernoulli randomization paired with a nonlinear estimator attains their minimax risk, for which they derive a second-order expansion. The present optimization fixes the Horvitz--Thompson estimator and varies covariate-dependent assignment under a smooth prognostic class. It characterizes the excess-variance gains available through assignment for that estimator and model.


\section{Setup and assumptions}

Throughout, \(n,d\geq2\) are integers and \(0<s\leq1\). Unit indices \leanref{sym:i}{\ensuremath{i}} range from \(1\) to \(n\), and coordinate indices \leanref{sym:j}{\ensuremath{j}} and \leanref{sym:ell}{\ensuremath{\ell}} range from \(1\) to \(d\), with \(j<\ell\) for coordinate pairs. A cube point is denoted by \leanref{sym:x}{\ensuremath{x}}. Integrals over cubes use Lebesgue measure, and torus integrals use normalized Haar probability. The notation \(\|\cdot\|_2\) denotes the Euclidean norm for vectors and the \(L^2\) norm for functions, as determined by the argument. Generic positive constants may change between occurrences; a subscript \(s\) records permitted dependence on smoothness. For nonnegative quantities, \(a\lesssim_s b\) means \(a\leq C(s)b\), and \(a\asymp_s b\) means that both comparison inequalities hold.

\subsection{Original-unit sampling}

The full covariate sample \leanref{sym:X}{\ensuremath{X=(X_1,\ldots,X_n)}} records the covariates of all original experimental units. Assignment takes values in the sign space \leanref{sym:Zspace}{\ensuremath{\mathcal Z_n=\{-1,1\}^n}}. An assignment vector \leanref{sym:Z}{\ensuremath{Z=(Z_1,\ldots,Z_n)}} selects treatment through the indicators \leanref{sym:A}{\ensuremath{A_i=(Z_i+1)/2}}. A covariate-dependent assignment law \leanref{sym:pi}{\ensuremath{\pi}} is a Borel probability kernel from \([0,1]^{nd}\) to \(\mathcal P(\mathcal Z_n)\), the probability measures on the finite sign space. Thus the design can use the entire original covariate array when assigning every unit.

The sampling distribution fixes the population measure against which component centering, regularity, and expected imbalance are evaluated.

\begin{assumptionv}{ass:uniform-draw}[Independent uniform covariate sample]
The covariate sample \(X=(X_1,\ldots,X_n)\) satisfies
\[
\operatorname{Law}(X)=\operatorname{Unif}([0,1]^d)^{\otimes n}.
\]
\end{assumptionv}

Independent uniform sampling is the exact uniform special case of the independent sampling and bounded-density condition in \citep[Assumption 2.1]{Cytrynbaum2026Limits}. It supplies a common population distribution for the design criterion and the canonical component decomposition. Independence also separates the randomness of the original sample from the dependence among treatment assignments that the design creates.

\subsection{Canonical components and regularity}

The prognostic function \leanref{sym:m}{\ensuremath{m}} describes the conditional mean of the potential-outcome midpoint. We work with its centered version, whose integral over the cube is zero. Under the product sampling distribution, the canonical main effects \leanref{sym:g1}{\ensuremath{g_j(m)}} and canonical pair effects \leanref{sym:g2}{\ensuremath{g_{j\ell}(m)}} are obtained by marginal integration. The following definition fixes their normalization and the interpretation of almost-everywhere equalities.

\begin{definitionv}{synth_1}[Canonical main and pair effects]
Let \(d\geq2\) and let \(m:[0,1]^d\to\mathbb R\) be a finite-valued Borel function in \(L^2([0,1]^d)\) with Lebesgue integral zero. Its canonical main and pair effects are
\[
g_j(m)(x_j)
=\int_{[0,1]^{d-1}}m(x_j,u_{-j})\,du_{-j},
\qquad 1\leq j\leq d,
\]
and
\[
g_{j\ell}(m)(x_j,x_\ell)
=\int_{[0,1]^{d-2}}
m(x_j,x_\ell,u_{-\{j,\ell\}})\,du_{-\{j,\ell\}}
-g_j(m)(x_j)-g_\ell(m)(x_\ell),
\qquad j<\ell.
\]
The integration arguments fill the remaining coordinates in their original order; a zero-dimensional integral is evaluation. These formulas define equivalence classes almost everywhere. Their centering convention is
\[
\int_0^1g_j(m)(v)\,dv=0,
\qquad
\int_0^1g_{j\ell}(m)(v,x_\ell)\,dv=0,
\qquad
\int_0^1g_{j\ell}(m)(x_j,v)\,dv=0,
\]
where the last two equalities hold for almost every remaining argument. Whenever \(m\) admits an exact order-two decomposition into square-integrable main and pair effects, these centering conditions determine its components uniquely up to Lebesgue null sets.
\end{definitionv}

Coordinatewise centering assigns each pair effect to variation jointly associated with its two coordinates. Together with the zero integral of \(m\), it fixes the intercept and identifies the components in the canonical ANOVA decomposition \citep{Hoeffding1948,Sobol1993}. The specification restricts the prognostic function to the sum of these main and pair effects.

\begin{assumptionv}{ass:order-two}[Exact order-two decomposition]
The prognostic function \(m\) has the decomposition
\[
m(x)=\sum_{j=1}^d g_j(m)(x_j)
+\sum_{j<\ell}g_{j\ell}(m)(x_j,x_\ell)
\]
for Lebesgue-almost every \(x\in[0,1]^d\), where \(g_j(m)\) and \(g_{j\ell}(m)\) are its canonical main and pair effects, respectively.
\end{assumptionv}

The exact order-two canonical ANOVA condition with zero intercept is the specification used here from \citep{Cytrynbaum2026Limits}. It describes a population model in which main effects and pairwise interactions account for the prognostic function exactly. Canonical centering makes the allocation of regularity across components intrinsic to \(m\).

Regularity is measured on the original cube by extension to Euclidean space. For a component dimension \leanref{sym:p}{\ensuremath{p\in\{1,2\}}}, the squared nonperiodic Sobolev restriction norm \leanref{sym:N}{\ensuremath{N_{p,s}(g)^2}} is
\[
N_{p,s}(g)^2
=
\inf_{\substack{G\in L^1(\mathbb R^p)\cap L^2(\mathbb R^p)\\
G|_{[0,1]^p}=g\ {\rm a.e.}}}
\int_{\mathbb R^p}
\left(1+\frac{\|\omega\|_2^2}{p}\right)^s
|\widehat G(\omega)|^2\,d\omega.
\]
Here the Euclidean Fourier transform \leanref{sym:Fourier}{\ensuremath{\widehat G}}, at angular frequency \leanref{sym:omega}{\ensuremath{\omega\in\mathbb R^p}}, uses the convention
\[
\widehat G(\omega)
=(2\pi)^{-p/2}\int_{\mathbb R^p}
e^{-\mathrm i\omega^{\mathsf T}u}G(u)\,du,
\]
and the dimension-normalized weight is exactly \((1+\|\omega\|_2^2/p)^s\). The infimum runs over integrable, square-integrable extensions agreeing almost everywhere on the cube. We impose one common budget across all canonical components.

\begin{assumptionv}{ass:pooled-budget}[Pooled component Sobolev budget]
The canonical main effects \(g_j(m)\) and pair effects \(g_{j\ell}(m)\) satisfy
\[
\sum_{j=1}^d N_{1,s}(g_j(m))^2
+\sum_{j<\ell}N_{2,s}(g_{j\ell}(m))^2
\leq 1,
\]
where \(N_{1,s}\) and \(N_{2,s}\) are the nonperiodic Sobolev restriction norms for one- and two-dimensional components, respectively.
\end{assumptionv}

The pooled component Sobolev budget fixes the dimension-normalized restriction-norm convention used throughout this paper. Pooling bounds the total squared component norm while allowing its distribution across coordinates and pairs to vary. Consequently, the dimension dependence of the criterion reflects the number of available interaction directions under a common regularity budget.

These component conditions define the population class over which the design will be evaluated. Write \leanref{sym:H}{\ensuremath{\mathcal H_{d,s}}} for this centered pooled Sobolev class.

\begin{assumptionv}{ass:fair}[Conditional fairness of assignment]
Under the assignment law \(\pi\), each assignment sign \(Z_i\) has conditional mean zero given the covariate sample \(X\):
\[
\mathbb E_\pi[Z_i\mid X]=0
\qquad\text{for every }i=1,\ldots,n.
\]
\end{assumptionv}

The finite-valued Borel representatives permit evaluation at sampled covariates. Uniform sampling makes changes on Lebesgue null sets immaterial to the expected loss, so the class retains the usual almost-everywhere interpretation of Sobolev functions.

\subsection{Assignment and signing loss}

The assignment family combines covariate-dependent balancing with equal marginal treatment probabilities. Its first requirement concerns each unit's assignment conditional on the complete covariate sample.

\begin{assumptionv}{ass:assignment-independence}[Conditional assignment independence]
A Borel assignment kernel \(\pi\) is outcome-independent if the following holds uniformly over pre-assignment data-generating laws. For every probability law \(\mu_X\) of the full covariate sample and every Markov kernel \(\kappa\) giving the pre-assignment potential-outcome array \(Y\) conditionally on \(X\), the experiment first draws \((X,Y)\) from \(\mu_X\otimes\kappa\) and then draws the assignment signs \(Z\) from \(\pi(X)\). Under this experiment law, \(Z\) and \(Y\) are conditionally independent given \(X\).
\end{assumptionv}

Conditional half-probability marginals are the fairness condition in \citep[Definition 2.2]{Cytrynbaum2026Limits}. Since signs take values in \(\{-1,1\}\), each unit receives treatment with conditional probability \(1/2\). The joint assignment distribution can nevertheless depend on all covariates and can correlate assignments across units to control prognostic imbalance.

The second requirement specifies the information available to the assignment law. Potential outcomes are generated before assignment, and the design draws signs using the covariate array.

\begin{definitionv}{def:sobolev-class}[Pooled Sobolev class \(\mathcal H_{d,s}\)]
The class \(\mathcal H_{d,s}\) consists of functions \(m:[0,1]^d\to\mathbb R\) with the following properties:
\begin{itemize}
\item \textbf{(Measurability.)} \(m\) is finite-valued and Borel measurable.
\item \textbf{(Centering.)} \(m\) belongs to \(L^2([0,1]^d)\) and has Lebesgue integral zero.
\item \textbf{(Component conditions.)} \(m\) satisfies \cref{obj:ass:order-two,obj:ass:pooled-budget} at \((d,s)\).
\end{itemize}
Canonical component equalities and centering are interpreted almost everywhere; arbitrary representatives on null sets give the same risk.
\end{definitionv}

Outcome-independent assignment is the information restriction in \citep[Definition 2.2]{Cytrynbaum2026Limits}. The sequential formulation expresses conditional independence for every pre-assignment outcome law: the design's randomization uses covariates and its own randomness. In conjunction with conditional fairness, it supplies the assignment conditions needed for the causal variance identity below.

Let \leanref{sym:D}{\ensuremath{\mathcal D_{n,d,s}}} denote the resulting admissible design class.

\begin{definitionv}{def:design-class}[Admissible design class $\mathcal D_{n,d,s}$]
The admissible design class is
\[
\mathcal D_{n,d,s}
=
\left\{
\pi:[0,1]^{nd}\to\mathcal P(\mathcal Z_n):
\pi \text{ is a Borel probability kernel satisfying }
\cref{obj:ass:fair,obj:ass:assignment-independence}
\text{ at }(n,d,s)
\right\}.
\]
Here \(\mathcal Z_n\) is the space of length-\(n\) assignment sign vectors, and \(\mathcal P(\mathcal Z_n)\) is its space of probability measures. Each kernel uses the covariates of all original units as its input, with assignment independent of function values, pilot outcomes, and the unknown outcome model.
\end{definitionv}

For a fixed admissible design and a fixed prognostic function, define the scaled signing loss \leanref{sym:loss}{\ensuremath{\mathcal L_{n,d,s}(\pi,m)}} by
\[
\mathcal L_{n,d,s}(\pi,m)
=
\frac4n
\mathbb E_{X,Z}
\left[
\left(\sum_{i=1}^n Z_i m(X_i)\right)^2
\right].
\]
The expectation integrates both the independent uniform sample and the conditional assignment law \(\pi(X)\). The minimax signing risk \leanref{sym:R}{\ensuremath{R_s(n,d)}} optimizes this population loss over admissible designs.

\begin{definitionv}{def:criterion}[Minimax signing risk $R_s(n,d)$]
The minimax scaled signing risk is
\[
R_s(n,d)
=
\inf_{\pi\in\mathcal D_{n,d,s}}
\sup_{m\in\mathcal H_{d,s}}
\frac{4}{n}
\mathbb E_{X,Z}
\left[
\left(\sum_{i=1}^{n}Z_i m(X_i)\right)^2
\right].
\]
Here \(\mathcal D_{n,d,s}\) is the admissible design class, \(\mathcal H_{d,s}\) is the centered prognostic class with pooled order-two Sobolev budget, \(X\) is the independent uniform covariate sample, and \(Z\) is the assignment sign vector drawn under \(\pi(X)\).
\end{definitionv}

The order of optimization is substantive. A design specifies its assignment rule as a function of the covariate array; a fixed population function is then evaluated under that rule, with sampling and assignment averaged inside the loss. The supremum therefore ranges over functions fixed across sample realizations. This order measures how well a single covariate-based design controls population prognostic imbalance uniformly over \(\mathcal H_{d,s}\).

To connect signing loss to treatment-effect estimation, fix the unit-weight Horvitz--Thompson estimator \citep{Horvitz1952}. In the next definition, \(A(a)\) denotes the estimator as a function of a treatment vector \(a\); the subscripted \(A_i\) denotes unit \(i\)'s treatment indicator.

\begin{definitionv}{synth_5}[Unit-weight Horvitz--Thompson estimator]
The unit-weight Horvitz--Thompson estimator \(A\) is defined, for any nonnegative integer \(n\), real potential-outcome schedule \(Y=((Y_i(0),Y_i(1)))_{i=1}^n\), and assignment signs \(Z\in\{-1,1\}^n\), by
\[
A\!\left(\frac{Z+1}{2}\right)
=
\frac{2}{n}\sum_{i=1}^n Z_i
Y_i\!\left(\frac{Z_i+1}{2}\right).
\]
Here \((Z_i+1)/2\) selects treatment when \(Z_i=1\) and control when \(Z_i=-1\). For \(n=0\), the estimator is defined to be zero.
\end{definitionv}

Conditional fairness gives each observed outcome its inverse assignment-probability weight. This estimator remains fixed throughout the design optimization in \cref{obj:def:criterion}.

\subsection{Spectral assignment features}

The assignment construction uses a \leanref{S-4}{reflected and commonly shifted representation} of the original covariates. For an integer \(p\), the period-two torus \leanref{sym:Torus}{\ensuremath{\mathbb T_2^p=[0,2)^p}} has addition modulo two. Its coordinatewise fold \leanref{sym:fold}{\ensuremath{b_p}} is
\[
b_p(y)_h=\min\{y_h,2-y_h\},
\qquad h=1,\ldots,p.
\]
The reflected prognostic function \leanref{sym:F}{\ensuremath{F=m\circ b_d}} is consequently a function on \(\mathbb T_2^d\).

An independent fair bit array \leanref{sym:reflection}{\ensuremath{(b_{ij})}}, with \(b_{ij}\in\{0,1\}\), selects the lifted covariates \leanref{sym:lift}{\ensuremath{\widetilde X_i}} coordinatewise:
\[
\widetilde X_{ij}
=
\begin{cases}
X_{ij},&b_{ij}=0,\\
2-X_{ij},&b_{ij}=1,
\end{cases}
\quad\text{interpreted modulo two}.
\]
Let \leanref{sym:T}{\ensuremath{T}} be a common Haar-uniform shift, independent of the covariates, reflections, and subsequent rounding randomness. The transformed covariate for original unit \(i\) is \leanref{sym:W}{\ensuremath{W_i=(\widetilde X_i+T)\bmod 2}}. Each row used by the assignment procedure still has one entry for each original unit.

For an integer frequency \leanref{sym:k}{\ensuremath{k\in\mathbb Z^d}}, the torus character \leanref{sym:ek}{\ensuremath{e_k(y)=\exp(\mathrm i\pi k^{\mathsf T}y)}} gives the normalized Fourier coefficient \leanref{sym:Fhat}{\ensuremath{\widehat F(k)}} through
\[
\widehat F(k)
=
\int_{\mathbb T_2^d}F(y)\overline{e_k(y)}\,dH(y),
\]
where \(H\) is normalized Haar probability. The frequency complexity \leanref{sym:t}{\ensuremath{t(k)}} ranks frequencies by squared Euclidean magnitude divided by the number of active coordinates. The ordered real features \leanref{sym:features}{\ensuremath{(\phi_h)_{h\geq1}}} convert this ranking into the prescribed rows for assignment.

\begin{definitionv}{synth_3}[Ordered real spectral features and rounding rows]
For \(d\geq2\), write \(\mathbb T_2^d=[0,2)^d\) with addition modulo two. For each \(k\in\mathbb Z^d\) with
\[
1\leq|\operatorname{supp}(k)|\leq2,
\qquad
\operatorname{supp}(k)=\{j:k_j\ne0\},
\]
choose the representative of \(\{k,-k\}\) whose first nonzero coordinate is positive. Order these representatives by increasing
\[
t(k)=\frac{\|k\|_2^2}{|\operatorname{supp}(k)|},
\]
breaking ties lexicographically. For each representative, list the two real features
\[
\sqrt2\cos(\pi k^{\mathsf T}y),
\qquad
\sqrt2\sin(\pi k^{\mathsf T}y),
\qquad y\in\mathbb T_2^d,
\]
consecutively in that order. This defines the ordered sequence \((\phi_h)_{h\geq1}\).

For a reflected and shifted sample \(W_1,\ldots,W_n\), the prescribed real spectral rows are
\[
a_h=\bigl(\phi_h(W_1),\ldots,\phi_h(W_n)\bigr)\in\mathbb R^n,
\qquad h\geq1.
\]
With \(K=\lfloor n/4\rfloor\), the input matrix for ordered boundary rounding is precisely
\[
a=\bigl(\phi_h(W_i)\bigr)_{\substack{1\leq h\leq K\\1\leq i\leq n}}
\in\mathbb R^{K\times n}.
\]
The feature ordering and the choice of rows depend on \(d\) and \(n\) and are independent of the smoothness parameter \(s\).
\end{definitionv}

The support restriction matches the main and pair structure of the prognostic class. Ordering the real features before selecting rows gives a single prescribed balancing matrix for each sample size and dimension. Its independence from \(s\) supplies the common feature input for the smoothness-independent assignment procedure.

\subsection{Outcome laws and causal excess variance}

The causal interpretation uses a single-unit law \(P\) of covariates \(U\) and potential outcomes \(O(0),O(1)\). Define its average treatment effect and conditional prognostic function by
\[
\vartheta(P)=\mathbb E_P[O(1)-O(0)],
\qquad
m_P(u)=
\mathbb E_P\!\left[
\frac{O(1)+O(0)}2\;\middle|\;U=u
\right].
\]
Let \leanref{sym:Q_s}{\ensuremath{\mathcal Q_{d,s}}} consist of these laws with uniform covariates, finite second moments for both potential outcomes, and \(m_P=m\) almost everywhere for some \(m\in\mathcal H_{d,s}\). Thus this outcome-law class fixes the prognostic restriction while allowing square-integrable conditional outcome distributions.

From \(n\) independent copies of \(P\), followed by assignment under \(\pi(U_{1:n})\), the actual estimator is
\[
\widehat T_P
=
\frac2n\sum_{i=1}^n Z_i
O_i\!\left(\frac{Z_i+1}{2}\right).
\]
Under the admissibility conditions, its expectation is \(\vartheta(P)\). The law-specific variance benchmark is
\[
\mathsf V(P)
=
\operatorname{Var}_P\!\left(
\mathbb E_P[O(1)-O(0)\mid U]
\right)
+
2\mathbb E_P\operatorname{Var}_P(O(1)\mid U)
+
2\mathbb E_P\operatorname{Var}_P(O(0)\mid U).
\]
This definition supplies the benchmark directly in terms of the outcome law. It combines variation in the conditional treatment effect with the two conditional outcome variances.

Define the causal excess-variance criterion by
\[
\mathcal V_n(\pi,P)
=
n\operatorname{Var}_{P,\pi}(\widehat T_P)-\mathsf V(P).
\]
The published excess-variance identity gives
\[
\mathcal V_n(\pi,P)
=
\frac4n\mathbb E_{P,\pi}
\left[
\left(\sum_{i=1}^n Z_i m_P(U_i)\right)^2
\right]
=
\mathcal L_{n,d,s}(\pi,m_P)
\]
for \(P\in\mathcal Q_{d,s}\) and \(\pi\in\mathcal D_{n,d,s}\)
\citep[Proposition 2.3, equation (2.4), under Assumption 2.1 and Definition 2.2]{Cytrynbaum2026Limits}.
Hence the signing loss is the scaled actual Horvitz--Thompson variance above a benchmark specific to the outcome law.

A \leanref{S-3}{Gaussian causal completion} supplies a concrete outcome law for every member of the prognostic class. Fix the mean-effect constant \leanref{sym:tau0}{\ensuremath{\tau_0=0.2}} and heterogeneity amplitude \leanref{sym:delta}{\ensuremath{\delta=0.1}}, and define the conditional treatment effect \leanref{sym:tau}{\ensuremath{\tau(x)}} by
\[
\tau(x)=\tau_0+\delta\sqrt2\sin(2\pi x_1).
\]
The Gaussian noise array \leanref{sym:epsilon}{\ensuremath{\epsilon}} has law \(\mathcal N(0,I_{2n})\), independently of \(X\). The pre-assignment potential-outcome schedule \leanref{sym:Y}{\ensuremath{Y}} is generated by
\[
Y_i(0)=m(X_i)-\frac{\tau(X_i)}2+\epsilon_i(0),
\qquad
Y_i(1)=m(X_i)+\frac{\tau(X_i)}2+\epsilon_i(1).
\]
Its conditional mean midpoint is \(m(X_i)\), so its single-unit law belongs to \(\mathcal Q_{d,s}\). The actual estimator \leanref{sym:theta_hat}{\ensuremath{\widehat\theta}}, realized finite-population target \leanref{sym:theta_n}{\ensuremath{\theta_n}}, and sampling-law target \leanref{sym:theta}{\ensuremath{\theta}} are
\[
\widehat\theta
=
\frac2n\sum_{i=1}^n Z_iY_i(A_i),
\qquad
\theta_n
=
\frac1n\sum_{i=1}^n\bigl(Y_i(1)-Y_i(0)\bigr),
\qquad
\theta=\mathbb E[\theta_n]=\tau_0.
\]
For this family, the law-specific benchmark is the constant \leanref{sym:Vstar}{\ensuremath{V^*=\delta^2+4=4.01}}. The published identity therefore expresses \(n\operatorname{Var}(\widehat\theta)-V^*\) as the signing loss of \(m\).

Because the Gaussian completion realizes every \(m\in\mathcal H_{d,s}\), the prognostic image of \(\mathcal Q_{d,s}\) is exactly \(\mathcal H_{d,s}\). It follows from the same identity that
\[
\inf_{\pi\in\mathcal D_{n,d,s}}
\sup_{P\in\mathcal Q_{d,s}}
\mathcal V_n(\pi,P)
=
R_s(n,d).
\]
This equality makes the minimax signing criterion a causal excess-variance criterion over the full uniform-covariate outcome-law class. The following result records the relation between that class and its Gaussian subclass, the corresponding published admissibility conditions, and elementary lower and upper guarantees.

\begin{propositionv}{prop:published-class-inclusion}[Published class inclusion and guarantees]*
Let \(0<s\leq1\), and use the prognostic class, design class, and minimax criterion of \cref{obj:def:sobolev-class,obj:def:design-class,obj:def:criterion}. For each integer \(d\geq2\), let \(\mathcal Q_{d,s}\) be the class of single-unit laws \(P\) of \((U,O(0),O(1))\) such that \(U\) is uniform on \([0,1]^d\), both potential outcomes have finite second moments, and, for some \(m\in\mathcal H_{d,s}\),
\[
m_P(u):=\mathbb E_P\!\left[\frac{O(1)+O(0)}2\;\middle|\;U=u\right]
=m(u)
\quad\text{for almost every }u\in[0,1]^d.
\]
Then the following statements hold.
\begin{itemize}
\item The laws obtained from \(m\in\mathcal H_{d,s}\) by the stated Gaussian causal completion form a proper subset of \(\mathcal Q_{d,s}\), for every integer \(d\geq2\).

\item For every pair of integers \(n,d\geq2\) and every Borel assignment probability kernel \(\pi\) based on the original covariate array, membership in \(\mathcal D_{n,d,s}\) is equivalent to the following published admissibility conditions: under independent uniform covariates \(X\),
\[
\mathbb E_\pi[Z_i\mid X]=0
\quad\text{almost surely, for every }i=1,\ldots,n,
\]
and, for every single-unit law \(P\) whose covariate density is bounded above and below by finite strictly positive constants and whose potential outcomes have finite second moments, the assignment signs are conditionally independent of the full potential-outcome array given the original covariate array in the experiment formed from \(n\) independent copies of \(P\) and the kernel \(\pi\). These are the admissibility conditions of \citet[Assumption 2.1 and Definition 2.2]{Cytrynbaum2026Limits}, specialized to the stated covariate distribution and available information.

\item For every integer \(n\geq2\), the two-dimensional minimax criterion satisfies
\[
R_s(n,2)\geq
\frac{2}{(48+32\pi^2)\,16384^s}\,n^{-s}.
\]

\item For every pair of integers \(n,d\geq2\) and every \(m\in\mathcal H_{d,s}\), independent fair assignment signs, independent of the original covariate sample, give
\[
\frac4n\,\mathbb E\!\left[
\left(\sum_{i=1}^n Z_i m(X_i)\right)^2
\right]
=
4\int_{[0,1]^d}m(x)^2\,dx
=
4\mathbb E[m(X_1)^2]
\leq4,
\]
where \(X_1,\ldots,X_n\) are independent uniform covariate vectors.
\end{itemize}
\end{propositionv}
% CAUSALSMITH-CITED-SCOPE-BEGIN prop:published-class-inclusion
\verificationfootnotetext{\textbf{Formalization scope.} This result uses the cited conclusion \citep{Cytrynbaum2026Limits} (Proposition 2.3 (Excess Variance), equation (2.4), with Assumption 2.1 and Definition 2.2) and \citep{BartlettMendelson2002Complexities} (Theorem 12(4), with Rademacher complexity as in Definition 2). The cited source proof is not formalized here: Lean verifies its use and all remaining steps. If that cited conclusion is false, the portions of this result that depend on it are not certified.}
% CAUSALSMITH-CITED-SCOPE-END prop:published-class-inclusion

\Cref{obj:prop:published-class-inclusion} distinguishes the outcome distributions from the prognostic functions that determine excess variance. The Gaussian family is a proper subclass of the admissible outcome laws, while realizing every prognostic function in the pooled Sobolev class. The admissibility equivalence aligns the original-unit assignment family with the published causal framework. Finally, the two-dimensional lower bound supplies an \(n^{-s}\) benchmark, and independent fair assignment supplies the universal upper bound of \(4\); these comparisons locate the design gains characterized by the main results.


\section{Main results}

The minimax criterion in \cref{obj:def:criterion} admits a two-sided characterization in terms of the number of coordinate pairs relative to the sample size. For presentation, define the interaction count \(B\) and the comparison scale \leanref{sym:a_s}{\ensuremath{a_s(n,d)}} by
\[
B=\binom d2,
\qquad
a_s(n,d)=\min\{1,(B/n)^s\},
\qquad n,d\geq2,\quad 0<s\leq1.
\]
The following result establishes the log-free capacity bound and the exact dimension condition for vanishing excess. The reflection lift and Gaussian completion appearing in the statement are specified in \cref{obj:synth_2,obj:synth_4}.

\begin{definitionv}{synth_4}[Pre-assignment Gaussian completion]
The Gaussian completion is the conditional law of the pre-assignment potential-outcome array \(Y\) given \(X\), generated independently across units by
\[
Y_i(0)=m(X_i)-\frac{0.2+0.1\sqrt2\sin(2\pi X_{i1})}{2}+\varepsilon_{i0},
\qquad
Y_i(1)=m(X_i)+\frac{0.2+0.1\sqrt2\sin(2\pi X_{i1})}{2}+\varepsilon_{i1}.
\]
The variables \(\varepsilon_{ia}\) are mutually independent standard Gaussians, independent of the covariates, and \(m\) is the centered measurable square-integrable prognostic function supplied to the completion.
\end{definitionv}

The comparison scale separates two regimes. In the subcritical regime \(B<n\), minimax excess is of order \((B/n)^s\): the interaction count enters through its ratio to the number of experimental units. When \(B\geq n\), the scale saturates at one, and the displayed lower and upper bounds are constants independent of \(n\) and \(d\). The assignment procedure \leanref{sym:pi_star}{\ensuremath{\pi^\star_{n,d}}} specified in \cref{obj:thm:log-free-bivariate-capacity} attains this scale in both regimes, using independent fair assignment in the saturated regime and spectral rounding in the subcritical regime.

At the endpoint \(s=1\), the subcritical rate is \(B/n\), with the same upper constant and a strictly positive lower constant. More generally, the lower constant \leanref{sym:c_s}{\ensuremath{c_s}} is bounded below by \(c_1>0\) throughout \(0<s\leq1\). The two-sided comparison therefore describes the endpoint together with the fractional-smoothness cases, with a logarithm-free rate. For each fixed smoothness value, it also gives an exact criterion for vanishing minimax scaled excess along arbitrary integer dimension sequences: the number of coordinate pairs must be negligible relative to sample size, equivalently \(d_n=o(\sqrt n)\). This equivalence accommodates irregular as well as monotone dimension sequences.

\subsection{Simultaneous attainment}

For the spectral branch of the assignment procedure, the auxiliary randomization transforms the covariates of the original units as follows.

\begin{definitionv}{synth_2}[Randomized reflection lift]
Let \(\mathbb T_2^d=[0,2)^d\) with coordinatewise addition modulo two. For \(x\in[0,1]^d\) and \(\eta\in\{0,1\}^d\), define the coordinatewise reflection map \(\operatorname{lift}:[0,1]^d\times\{0,1\}^d\to\mathbb T_2^d\) by
\[
\operatorname{lift}(x,\eta)_j
=
\begin{cases}
x_j \pmod 2,&\eta_j=0,\\
2-x_j \pmod 2,&\eta_j=1.
\end{cases}
\]
For the original sample \(X=(X_1,\ldots,X_n)\), take independent fair coins \(\eta_{ij}\), independent of \(X\), the common shift, and all rounding seeds, and put
\[
\widetilde X_i=\operatorname{lift}(X_i,\eta_i),
\qquad
W_i=\widetilde X_i+T\pmod2,
\]
where \(\eta_i=(\eta_{i1},\ldots,\eta_{id})\) and \(T\) is an independent common Haar-uniform shift on \(\mathbb T_2^d\).
\end{definitionv}

The reflections act separately across units and coordinates, whereas the Haar shift is shared by all units. The transformed covariates in \cref{obj:synth_2} supply the spectral rows used to generate assignment signs for the original sample.

The simultaneous guarantee separates the choice of assignment from the smoothness value used to assess precision. The procedure in \cref{obj:thm:log-free-bivariate-capacity} uses the same spectral rows and randomization for every \(0<s\leq1\), while the resulting bound varies with the exponent \(s\). Its finite completion property applies to every original sample and every allowed realization of the auxiliary seeds. Thus the displayed risk bounds describe a fully specified assignment law on the original units.

\subsection{Log-free capacity and dimension frontier}

The lower bound has a separate statistical role: it applies to every admissible assignment law through a single prior chosen independently of the covariate sample. The prior specification and its class-membership properties are developed further in \cref{obj:def:cosine-prior,obj:lem:cosine-prior}.

The normalized prognostic function \leanref{sym:m_xi}{\ensuremath{m_\xi}} in \cref{obj:def:cosine-prior} ranges over a finite family of pair-cosine functions. Its law, \leanref{sym:nu_star}{\ensuremath{\nu^\star_{n,d,s}}}, depends on sample size, dimension, and smoothness, while remaining independent of both the covariate sample and the chosen assignment law. Every realization belongs to the centered pooled Sobolev class by \cref{obj:thm:log-free-bivariate-capacity}. The contraction inequality used in the supporting converse analysis is \citet[Theorem~12(4), with Rademacher complexity as in Definition~2]{BartlettMendelson2002Complexities}; its applications here use the absolute-value map with finite classes and finite-dimensional Euclidean unit-ball linear classes. The following result states the resulting prior-average bound over the full design class.

\begin{theoremv}{thm:log-free-bivariate-capacity}[Log-free capacity and dimension frontier]*
For every pair of integers \(n,d\geq2\), write
\[
B=\binom d2.
\]
There is a single fair, outcome-independent Borel original-sample assignment kernel \(\pi^\star_{n,d}\), independent of the smoothness parameter, belonging to the design class of \cref{obj:def:design-class}. It assigns independent fair signs when \(n\leq B\). When \(B<n\), it is the law of the signs produced by spectral rounding with independent reflection coins, an independent common Haar shift, and independent uniform rounding seeds.

More precisely, the reflected and shifted sample is
\[
W_i=\operatorname{lift}(X_i,\text{reflection coins})+T,
\]
where \(T\) is the common Haar shift on the period-two torus. The input rows are the first \(\lfloor n/4\rfloor\) prescribed real features evaluated at \(W_i\). For every choice of reflection coins and common shift, and every rounding-seed array whose entries belong to \([0,1]\), the fractional coordinate after \(n\) rounding iterations equals the corresponding output sign, for every original sample and every unit.

For every real \(s\) with \(0<s\leq1\), let
\[
a_s(n,d)=b_s(n,d)=\min\{1,(B/n)^s\},
\qquad
c_s=\frac{2}{(48+32\pi^2)16384^s},
\qquad
c_1=\frac{2}{(48+32\pi^2)16384}.
\]
Thus \(a_s(n,d)=1\) when \(n<B\), and \(a_s(n,d)=(B/n)^s\) when \(B\leq n\). The constants satisfy
\[
0<c_1\leq c_s<3072.
\]
With \(\mathcal H_{d,s}\) as in \cref{obj:def:sobolev-class}, and \(R_s(n,d)\) the minimax criterion of \cref{obj:def:criterion}, define the scaled signing loss by
\[
\mathcal L_{n,d,s}(\pi,m)
=\frac4n\,\mathbb E_{\pi}
 \left(\sum_{i=1}^n Z_i m(X_i)\right)^2.
\]
The same kernel satisfies, simultaneously for all \(0<s\leq1\),
\[
c_1a_s(n,d)
\leq c_sa_s(n,d)
\leq R_s(n,d)
\leq \sup_{m\in\mathcal H_{d,s}}
       \mathcal L_{n,d,s}(\pi^\star_{n,d},m)
\leq3072a_s(n,d).
\]

For each such \(s\), the finite coefficient prior \(\nu^\star_{n,d,s}\), with normalization
\[
C_0=\sqrt{48+32\pi^2},
\]
has every coefficient-sign realization \(m_\xi\) centered and in \(\mathcal H_{d,s}\):
\[
\int_{[0,1]^d}m_\xi(x)\,dx=0,
\qquad
m_\xi\in\mathcal H_{d,s}.
\]
Drawn independently of the covariate sample, it certifies, for every admissible design \(\pi\),
\[
c_sa_s(n,d)
\leq
\mathbb E_{m\sim\nu^\star_{n,d,s}}
       \mathcal L_{n,d,s}(\pi,m).
\]

Let \(\mathcal Q_{d,s}\) be the frozen uniform-covariate, zero-intercept, correctly specified outcome-law class with finite second moments and centered prognostic function in \(\mathcal H_{d,s}\). For a single-unit law \(P\), its prognostic function, actual HT estimator, law-specific benchmark, and scaled excess are
\[
m_P(u)=\mathbb E_P[(O(1)+O(0))/2\mid U=u],
\qquad
\widehat T_P=\frac2n\sum_{i=1}^n
 Z_iO_i((Z_i+1)/2),
\]
\[
\mathsf V(P)
=\operatorname{Var}_P\!\left(
 \mathbb E_P[O(1)-O(0)\mid U]\right)
+2\mathbb E_P\operatorname{Var}_P(O(1)\mid U)
+2\mathbb E_P\operatorname{Var}_P(O(0)\mid U),
\]
\[
\mathcal V_n(\pi,P)
=n\operatorname{Var}_{P,\pi}(\widehat T_P)-\mathsf V(P).
\]
Here the experiment uses independent original-unit copies, and conditional means are interpreted modulo covariate null sets. Then
\[
\inf_{\pi\in\mathcal D_{n,d,s}}
 \sup_{P\in\mathcal Q_{d,s}}\mathcal V_n(\pi,P)
=R_s(n,d),
\]
\[
\sup_{P\in\mathcal Q_{d,s}}
 \mathcal V_n(\pi^\star_{n,d},P)
=
\sup_{m\in\mathcal H_{d,s}}
 \mathcal L_{n,d,s}(\pi^\star_{n,d},m).
\]
The Gaussian completions of functions in \(\mathcal H_{d,s}\) form a proper subclass of \(\mathcal Q_{d,s}\). Restricting the law supremum to these Gaussian completions gives the same minimax capacity \(R_s(n,d)\) and the same worst-case loss for \(\pi^\star_{n,d}\).

For every pair of real density bounds satisfying
\[
0<\mathrm{lower}\leq1\leq\mathrm{upper},
\]
let \(\mathcal A_{d,s}\) be the bounded-density class defined in \cref{obj:lem:published-prognostic-capacity}: its laws have finite potential-outcome second moments, satisfy
\[
\int m_P(u)^2\,dP_U(u)\leq1,
\]
and have \(m_P\) equal almost everywhere to a permitted intercept plus a centered function in \(\mathcal H_{d,s}\), with covariate density between the displayed bounds. Its minimax excess satisfies
\[
R_s(n,d)
\leq
\inf_{\pi\in\mathcal D_{n,d,s}}
 \sup_{P\in\mathcal A_{d,s}}\mathcal V_n(\pi,P).
\]

For every admissible \(\pi\) and every \(m\in\mathcal H_{d,s}\), the actual HT estimator \(\widehat\theta\) under the Gaussian completion of \(m\), with benchmark \(V^*=4.01\), satisfies
\[
n\operatorname{Var}(\widehat\theta)-4.01
=\mathcal L_{n,d,s}(\pi,m)<\infty.
\]
Consequently, for every \(\varepsilon>0\),
\[
n\geq B\max\!\left\{
1,\left(\frac{3072}{4.01\varepsilon}\right)^{1/s}
\right\}
\quad\Longrightarrow\quad
\sup_{m\in\mathcal H_{d,s}}
 \mathcal L_{n,d,s}(\pi^\star_{n,d},m)
\leq4.01\varepsilon.
\]
Conversely, for every \(0<\varepsilon<c_s/4.01\),
\[
R_s(n,d)\leq4.01\varepsilon
\quad\Longrightarrow\quad
n\geq B\left(\frac{c_s}{4.01\varepsilon}\right)^{1/s}.
\]

Finally, for every fixed \(s\in(0,1]\) and every integer sequence \(d_n\geq2\),
\[
R_s(n,d_n)\longrightarrow0
\quad\Longleftrightarrow\quad
\frac{\binom{d_n}{2}}n\longrightarrow0
\quad\Longleftrightarrow\quad
d_n=o(\sqrt n).
\]
For every fixed pair \(0<\mathrm{lower}\leq1\leq\mathrm{upper}\), vanishing minimax excess over the corresponding bounded-density classes implies the same dimension frontier:
\[
\inf_{\pi\in\mathcal D_{n,d_n,s}}
 \sup_{P\in\mathcal A_{d_n,s}}\mathcal V_n(\pi,P)
\longrightarrow0
\quad\Longrightarrow\quad
d_n=o(\sqrt n).
\]
\end{theoremv}
% CAUSALSMITH-CITED-SCOPE-BEGIN thm:log-free-bivariate-capacity
\verificationfootnotetext{\textbf{Formalization scope.} This result uses the cited conclusion \citep{BartlettMendelson2002Complexities} (Theorem 12(4), with Rademacher complexity as in Definition 2) and \citep{Cytrynbaum2026Limits} (Proposition 2.3 (Excess Variance), equation (2.4), with Assumption 2.1 and Definition 2.2). The cited source proof is not formalized here: Lean verifies its use and all remaining steps. If that cited conclusion is false, the portions of this result that depend on it are not certified.}
% CAUSALSMITH-CITED-SCOPE-END thm:log-free-bivariate-capacity

\subsection{An independent-prior converse}

The converse in \cref{obj:thm:full-design-lower} certifies the lower comparison scale \leanref{sym:b_s}{\ensuremath{b_s(n,d)}} for every admissible assignment law, including laws that depend on the entire observed covariate sample. The frequency cutoff lets the prior supply the subcritical lower bound as well as the saturated bound. Because the guarantee concerns an average over legal prognostic functions drawn independently of the sample, each admissible design has worst-case loss at least as large as that average.

The next theorem isolates the prior-average lower bound that drives the converse for the entire admissible design class.

\begin{theoremv}{thm:full-design-lower}[Full-design capacity lower bound]*
For every \(s\in(0,1]\), define the lower constant \(c_s\) and the comparison scale \(b_s(n,d)\) by
\[
c_s=\frac{2}{(48+32\pi^2)\,16384^s},
\qquad
b_s(n,d)=\min\left\{1,\left(\frac{\binom d2}{n}\right)^s\right\}.
\]
Then \(c_s>0\), independently of \(n,d\). For all integers \(n,d\geq2\), every design \(\pi\in\mathcal D_{n,d,s}\) of \cref{obj:def:design-class} satisfies
\[
\mathbb E_\xi\,\mathcal L_{n,d,s}(\pi,m_\xi)
\geq c_s b_s(n,d),
\]
where the expectation is over the independent coefficient-sign cosine prior and \(m_\xi\) is its associated prognostic function; the coefficient signs are drawn independently of the covariate sample \(X\). Consequently, the minimax criterion of \cref{obj:def:criterion} satisfies
\[
R_s(n,d)\geq c_s b_s(n,d).
\]
\end{theoremv}
% CAUSALSMITH-CITED-SCOPE-BEGIN thm:full-design-lower
\verificationfootnotetext{\textbf{Formalization scope.} This result uses the cited conclusion \citep{BartlettMendelson2002Complexities} (Theorem 12(4), with Rademacher complexity as in Definition 2). The cited source proof is not formalized here: Lean verifies its use and all remaining steps. If that cited conclusion is false, the portions of this result that depend on it are not certified.}
% CAUSALSMITH-CITED-SCOPE-END thm:full-design-lower

\subsection{Causal transfer and prospective precision}

The completion in \cref{obj:synth_4} retains \(m\) as the conditional midpoint mean while fixing the treatment-effect function and the conditional noise variances. Its scaled variance benchmark is \(4.01\), so a signing-loss bound translates directly into an actual excess-variance bound. This translation uses the excess-variance identity of \citet[Proposition~2.3 (Excess Variance), equation~(2.4), with Assumption~2.1 and Definition~2.2]{Cytrynbaum2026Limits}, under conditional fairness, outcome-independent assignment, and finite potential-outcome second moments.

The exact minimax equality has a causal interpretation. The published excess-variance identity makes the conditional midpoint mean the function that determines excess above the law-specific benchmark; the Gaussian completion realizes every function in the prognostic class. Accordingly, \cref{obj:thm:log-free-bivariate-capacity} equates the prognostic minimax criterion with the outcome-law minimax criterion over \(\mathcal Q_{d,s}\), and gives the corresponding equality for the worst-case loss of the specified procedure. The same capacity is realized within the Gaussian subclass, where the benchmark is \(4.01\).

For prospective precision calculations under the Gaussian completion, the sufficient sample-size bound guarantees scaled variance at most \(4.01(1+\varepsilon)\) uniformly over the prognostic class. The necessary bound applies when \(0<\varepsilon<c_s/4.01\). At \(s=1\), \(c_1\approx3.36\times10^{-7}\), so this displayed tolerance range is approximately \(\varepsilon<8.4\times10^{-8}\); the ratio \(3072/c_1\) exceeds nine billion. These constants are conservative finite-sample certificates of the interaction-count scaling rather than a claim of sharp numerical calibration.

The larger bounded-density class in \cref{obj:lem:published-prognostic-capacity} inherits the converse because its minimax scaled excess is at least \(R_s(n,d)\). For fixed density bounds \(0<\mathrm{lower}\leq1\leq\mathrm{upper}\) and fixed smoothness, vanishing minimax scaled excess in that class therefore requires \(d_n=o(\sqrt n)\). This necessary dimension condition applies with the permitted intercept, the stated prognostic second-moment bound, and finite potential-outcome second moments.


\section{A smoothness-independent assignment procedure}

The simultaneous attainment in \cref{obj:thm:log-free-bivariate-capacity} is achieved by an assignment rule that prioritizes low-frequency covariate balance. The rule uses sample size, dimension, observed covariates, and auxiliary random seeds. The same rule serves every smoothness level \(0<s\leq1\): smoothness enters the evaluation of its performance through the spectral budget.

We first describe the spectral construction used to supply the balancing rows. Work on the period-two torus \(\mathbb T_2^d=[0,2)^d\), with addition modulo two and normalized Haar probability. Its coordinatewise fold onto the cube is
\[
b_d(y)_j=\min\{y_j,2-y_j\}.
\]
For each original unit and coordinate, independently choose the lift \(X_{ij}\) or \(2-X_{ij}\), each with probability \(1/2\). Add a common Haar-uniform shift \(T\in\mathbb T_2^d\), independent of the sample, the reflection choices, and the rounding seeds. The resulting transformed covariates are
\[
W_i=\bigl(\text{lift of }X_i+T\bigr)\bmod 2.
\]
These transformations retain the correspondence between each transformed covariate vector and its original experimental unit.

The reflected prognostic function is \(F=m\circ b_d\). Writing
\[
e_k(y)=\exp(\mathrm i\pi k^{\mathsf T}y),
\qquad
\widehat F(k)=\int_{\mathbb T_2^d}F(y)\overline{e_k(y)}\,dy,
\]
where integration uses normalized Haar probability, gives the Fourier representation underlying the spectral budget. Reflection and translation preserve the original prognostic values through
\[
F(W_i-T)=m(X_i).
\]
The order-two structure in \cref{obj:ass:order-two} directs attention to nonzero integer frequencies supported on one or two coordinates. Their frequency complexity is
\[
t(k)=\frac{\|k\|_2^2}{|\operatorname{supp}(k)|},
\qquad
\operatorname{supp}(k)=\{j:k_j\ne0\}.
\]
Take one representative of each pair \(\{k,-k\}\), choosing the representative whose first nonzero entry is positive. Order these representatives by increasing complexity, with lexicographic tie breaking, and place the cosine row before the sine row for each representative. The corresponding real features are
\[
\sqrt2\cos(\pi k^{\mathsf T}W_i),
\qquad
\sqrt2\sin(\pi k^{\mathsf T}W_i).
\]
The first \(\lfloor n/4\rfloor\) rows of this ordered family form the input matrix for rounding. Thus reflection, the common shift, frequency ordering, and real features constitute a single construction shared across smoothness levels.

The assignment procedure protects an initial segment of these rows while enough fractional coordinates remain available. As coordinates reach their assignment boundaries, the protected segment becomes shorter. Low-frequency rows consequently receive protection for more of the procedure. Randomization between the two boundary directions preserves the conditional mean of each assignment coordinate, linking covariate balance to fair marginal assignment probabilities.

The independent prior used in the converse places random signs on a finite family of pair-cosine features. Its sample-dependent cutoff and smoothness-dependent normalization are specified next.

\begin{definitionv}{def:cosine-prior}[Pair-cosine prior \(m_\xi\)]
Let
\[
B=\binom d2,\qquad
L=\max\left\{1,\left\lceil\sqrt{4096n/B}\right\rceil\right\},\qquad
M=BL^2,\qquad
C_0=\sqrt{48+32\pi^2}.
\]
Before observing \(X\), draw every coefficient \(\xi_\alpha\) independently and uniformly from \(\{-1,+1\}\). For \(\alpha=(j,h,k,\ell)\), where \(j<h\) and \(1\leq k,\ell\leq L\), put
\[
V_\alpha(x)=2\cos(\pi kx_j)\cos(\pi\ell x_h),
\qquad
m_\xi(x)=\frac{1}{C_0L^s\sqrt M}\sum_\alpha\xi_\alpha V_\alpha(x).
\]
The prior \(\nu^\star_{n,d,s}\) is the law of this normalized pair-cosine function.
\end{definitionv}

The pair-cosine features \leanref{sym:V}{\ensuremath{V_\alpha(x)}} in \cref{obj:def:cosine-prior} are indexed by \leanref{sym:alpha}{\ensuremath{\alpha}}; the frequency cutoff \leanref{sym:L}{\ensuremath{L}} gives the feature count \leanref{sym:M}{\ensuremath{M}}. The coefficient signs \leanref{sym:xi}{\ensuremath{\xi_\alpha}} are sampled independently of \(X\), while the normalization constant \leanref{sym:C0}{\ensuremath{C_0}} and the factor \(L^s\sqrt M\) determine the amplitude of the prior functions.

The following construction combines this assignment rule with the comparison scale and the independent prior used in the converse.

\begin{algorithmv}{def:capacity-handle}[Capacity construction \((a_s,\pi^\star,\nu^\star)\)]
For \(n,d\) and \(s\), construct the capacity handle as follows.
\begin{enumerate}
\item Set
\[
a_s(n,d)=\min\{1,(\tbinom d2/n)^s\}.
\]
\item If \(n\leq\binom d2\), let \(\pi^\star_{n,d}\) assign independent fair signs. If \(\binom d2<n\), let \(\pi^\star_{n,d}\) use the reflected, commonly shifted spectral features and ordered-boundary rounding on the original sample.
\item Let \(\nu^\star_{n,d,s}\) be the independent uniform coefficient-sign prior defining the pair-cosine functions at the cutoff in \cref{obj:def:cosine-prior}.
\item Return the triple
\[
(a_s(n,d),\pi^\star_{n,d},\nu^\star_{n,d,s}),
\]
\end{enumerate}
\end{algorithmv}

At saturation, \(n\leq\binom d2\), the comparison scale in \cref{obj:def:capacity-handle} equals one, and independent fair signs attain the upper comparison established in \cref{obj:thm:log-free-bivariate-capacity}. When the sample exceeds the coordinate-pair count, the same construction instead balances the ordered spectral features. Its assignment component \(\pi^\star_{n,d}\), defined in \cref{obj:def:capacity-handle}, is independent of \(s\), whereas the comparison scale and the coefficient prior retain their smoothness dependence.

To implement the balancing branch, the fractional assignment starts at zero and moves within the active-coordinate subspace subject to the current row constraints. The boxed procedure specifies the ordering, boundary probabilities, and tie conventions that determine each move.

\begin{algorithmv}{def:ordered-boundary-rounding}[Ordered boundary rounding $Z$]
The input is a real matrix \(a\in\mathbb R^{K\times n}\), with
\[
K=\lfloor n/4\rfloor.
\]
\begin{enumerate}
\item Initialize the fractional assignment state by
\[
u=0\in[-1,1]^n.
\]
A coordinate \(i\) is active exactly when \(|u_i|<1\). Inactive coordinates remain fixed.

\item At the start of each phase, set
\[
r=\#\{i:|u_i|<1\}.
\]
If \(r\leq3\), round the remaining active coordinates independently, in increasing index order, according to
\[
\mathbb P(u_i\leftarrow+1)=\frac{1+u_i}{2},
\qquad
\mathbb P(u_i\leftarrow-1)=\frac{1-u_i}{2},
\]
and proceed to the output step. If \(r\geq4\), set
\[
q_1=\lfloor r/4\rfloor
\]
and retain this value throughout the phase, until the active count is at most \(\lfloor r/2\rfloor\).

\item At each move, form the orthogonal projection \(P\) onto the vectors supported on the current active coordinates and orthogonal to the first \(q_1\) input rows:
\[
P=
\operatorname{proj}_{\left\{
v\in\mathbb R^n:
v_i=0\ \text{if }|u_i|=1,\ 
\sum_{i=1}^n a_{\ell i}v_i=0\ (1\leq\ell\leq q_1)
\right\}}.
\]
Compute \(P\) by orthonormalizing the constraint rows restricted to the active coordinates in row order, discarding zero Gram--Schmidt residuals and normalizing residuals with positive norm. Subtract their rank-one projections from the identity on the active coordinate subspace.

\item Process \(Pe_1,\ldots,Pe_n\) in index order, where \(e_1,\ldots,e_n\) are the standard coordinate vectors. Discard zero Gram--Schmidt residuals and normalize residuals with positive norm to obtain the ordered orthonormal basis \(v_1,\ldots,v_h\), characterized by
\[
h=\dim(\operatorname{im}P),
\qquad
\operatorname{span}\{v_1,\ldots,v_h\}=\operatorname{im}P,
\qquad
v_b^{\mathsf T}v_c=
\begin{cases}
1,&b=c,\\
0,&b\ne c.
\end{cases}
\]

\item For each basis vector, compute the boundary step lengths and their product:
\[
\delta_b^+
=
\min\left(
\left\{\frac{1-u_i}{(v_b)_i}:(v_b)_i>0\right\}
\cup
\left\{\frac{1+u_i}{-(v_b)_i}:(v_b)_i<0\right\}
\right),
\]
\[
\delta_b^-
=
\min\left(
\left\{\frac{1+u_i}{(v_b)_i}:(v_b)_i>0\right\}
\cup
\left\{\frac{1-u_i}{-(v_b)_i}:(v_b)_i<0\right\}
\right),
\qquad
t_b=\delta_b^+\delta_b^-.
\]
These sets contain active nonzero entries. Set the normalizer to
\[
S=\sum_{b=1}^h t_b^{-1}.
\]

\item Choose a basis index according to
\[
\mathbb P(\text{chosen index}=b)=\frac{t_b^{-1}}{S}.
\]
Conditional on that choice, update the state according to
\[
u\leftarrow
\begin{cases}
u+\delta_b^+v_b,
&\text{with probability }\displaystyle\frac{\delta_b^-}{\delta_b^++\delta_b^-},\\[6pt]
u-\delta_b^-v_b,
&\text{with probability }\displaystyle\frac{\delta_b^+}{\delta_b^++\delta_b^-}.
\end{cases}
\]
Repeat the moves until the active count is at most \(\lfloor r/2\rfloor\), then start a new phase. Every random choice uses an independent uniform seed. Inverse distribution functions use half-open intervals, assigning the value one to the last interval. Comparisons with zero and boundary comparisons are exact; ties in finite minima use the least index.

\item Output the terminal assignment sign vector
\[
Z=u\in\mathcal Z_n.
\]
\end{enumerate}
\end{algorithmv}

The protected rows in \cref{obj:def:ordered-boundary-rounding} retain their current imbalances throughout each phase. Starting from zero therefore gives the earliest rows the strongest initial protection. The asymmetric boundary probabilities and the terminal coordinate probabilities preserve conditional means, yielding a fair sign for every original unit. Each boundary move fixes at least one previously active coordinate, so the construction requires at most \(n\) boundary moves, followed by independent rounding of at most three remaining coordinates. This is a finite guarantee under the specified exact-real arithmetic conventions.

Every output coordinate remains attached to its original unit. Reflection and shifting provide balancing features, while all \(n\) units enter the original Horvitz--Thompson estimator with their original weight \(2/n\). Consequently, the assignment improvement applies directly to the original-sample criterion in \cref{obj:def:criterion}.

The quantitative link between ordered protection and prognostic imbalance is supplied by \cref{obj:lem:ordered-prefix-rounding}. Its row-imbalance bound combines with the reflected spectral budget and transfer in \cref{obj:lem:exact-reflection-budget,obj:lem:reflection-transfer}. Frequency ordering allocates the strongest balance to the rows whose coefficients receive the greatest weight under smoothness. These bounds yield the simultaneous smoothness attainment stated in \cref{obj:thm:log-free-bivariate-capacity}. The supporting bounds and measurable realization arguments accompany those auxiliary results in the appendix.


\section{Discussion and extensions}

The characterization in \cref{obj:thm:log-free-bivariate-capacity} gives an interaction-count benchmark for experimental sample size. Under the independent uniform sampling law of \cref{obj:ass:uniform-draw}, the relevant comparison is between \(B=\binom d2\), the number of coordinate pairs, and \(n\), the number of experimental units. For the centered prognostic class with the pooled Sobolev budget in \cref{obj:def:sobolev-class}, minimax scaled excess variance is comparable to \(\min\{1,(B/n)^s\}\). Thus the interaction count summarizes the dimension dependence of the worst-case precision requirement within this class.

\subsection{Causal calibration and relative excess}

To interpret the assignment guarantee in terms of causal precision, we first specify the potential outcomes, treatment, estimator, and targets used in the Gaussian calibration.

\begin{definitionv}{def:causal-completion}[Causal completion and estimator \(\widehat\theta\)]
The actual estimator \(\widehat\theta\) uses \(Y\), the pre-assignment potential-outcome array, and the treatment indicator
\[
A=(Z+1)/2,
\]
where \(Z\) is the vector of assignment signs. Its conditional estimand is \(\theta_n\), the average realized potential-outcome contrast, conditional on the realized full schedule. Its sampling-law estimand is \(\theta\), the sampling-law mean of that contrast. The setting and constructions are specified in \cref{obj:ass:order-two,obj:ass:pooled-budget,obj:ass:uniform-draw,obj:ass:fair,obj:ass:assignment-independence,obj:def:sobolev-class,obj:def:design-class,obj:def:criterion,obj:def:cosine-prior,obj:def:capacity-handle,obj:def:ordered-boundary-rounding,obj:lem:cosine-prior,obj:thm:full-design-lower,obj:lem:reflection-transfer,obj:lem:ht-transfer,obj:prop:published-class-inclusion,obj:thm:log-free-bivariate-capacity,obj:lem:isotropic-row-signing,obj:lem:ordered-prefix-rounding,obj:lem:published-prognostic-capacity,obj:thm:capacity-answer,obj:lem:exact-reflection-budget,obj:synth_1,obj:synth_2,obj:synth_3,obj:synth_4,obj:synth_5}.
\end{definitionv}

Concretely, for each \(m\in\mathcal H_{d,s}\), the completion in \cref{obj:synth_4} fixes
\[
\tau(x)=0.2+0.1\sqrt2\sin(2\pi x_1)
\]
and generates the pre-assignment schedule as
\[
Y_i(0)=m(X_i)-\frac{\tau(X_i)}2+\epsilon_i(0),
\qquad
Y_i(1)=m(X_i)+\frac{\tau(X_i)}2+\epsilon_i(1).
\]
The \(2n\) noise variables are independent standard Gaussians, independent of the covariates. After this schedule is generated, the admissible assignment law draws \(Z\) using the covariates and satisfies \cref{obj:ass:fair,obj:ass:assignment-independence}. With \(A_i=(Z_i+1)/2\), the observed outcome is \(Y_i(A_i)\), and the estimator and its two targets are
\[
\widehat\theta=\frac2n\sum_{i=1}^n Z_iY_i(A_i),
\qquad
\theta_n=\frac1n\sum_{i=1}^n\bigl(Y_i(1)-Y_i(0)\bigr),
\qquad
\theta=\mathbb E[\theta_n]=0.2.
\]
Fairness and outcome-independent assignment make the estimator conditionally unbiased for the realized schedule average \(\theta_n\), and its sampling-law expectation is \(\theta\).

The benchmark for this family is \(V^*=4.01\): treatment-effect heterogeneity contributes \(0.01\), and the two unit conditional outcome variances contribute \(4\). Here variance is taken over covariate sampling, potential-outcome noise, and assignment. The exact identity in \cref{obj:lem:ht-transfer} is
\[
n\operatorname{Var}(\widehat\theta)-V^*
=
\mathcal L_{n,d,s}(\pi,m).
\]
Consequently, the signing criterion measures actual Horvitz--Thompson excess variance above \(V^*/n\) for this completion. The Gaussian family calibrates the abstract imbalance loss while realizing every prognostic function in the pooled class.

A relative-excess tolerance \(\varepsilon>0\) requires
\[
n\operatorname{Var}(\widehat\theta)
\leq V^*(1+\varepsilon)
\]
uniformly over the Gaussian completions indexed by \(m\in\mathcal H_{d,s}\). For the assignment procedure in \cref{obj:def:capacity-handle}, the sufficient sample-size implication of \cref{obj:thm:log-free-bivariate-capacity} is
\[
n\geq B\max\left\{
1,\left(\frac{3072}{4.01\varepsilon}\right)^{1/s}
\right\}
\quad\Longrightarrow\quad
\sup_{m\in\mathcal H_{d,s}}
\left\{
\frac{n\operatorname{Var}(\widehat\theta)}{4.01}-1
\right\}
\leq\varepsilon.
\]
The same assignment procedure supplies this guarantee for every \(0<s\leq1\).

Conversely, any admissible design meeting that uniform relative-excess requirement must have \(R_s(n,d)\leq4.01\varepsilon\). For \(0<\varepsilon<c_s/4.01\), where \(c_s\) is the lower comparison constant in \cref{obj:thm:log-free-bivariate-capacity}, this entails
\[
n\geq B\left(\frac{c_s}{4.01\varepsilon}\right)^{1/s}.
\]
The tolerance restriction places the necessary implication below the saturated lower bound. At \(s=1\), \(c_1\approx3.36\times10^{-7}\), hence \(c_1/4.01\approx8.4\times10^{-8}\), while \(3072/c_1>9\times10^9\). The constants therefore provide conservative finite-sample certificates of the \(B\varepsilon^{-1/s}\) scaling; they should not be read as a sharp numerical calibration.

\subsection{Precision over bounded-density classes}

The broader comparison in \cref{obj:thm:log-free-bivariate-capacity} gives a necessary precision condition beyond uniform sampling. The class \(\mathcal A_{d,s}\) defined in \cref{obj:lem:published-prognostic-capacity} allows covariate densities between fixed positive finite bounds containing one. Its potential outcomes have finite second moments, and its conditional midpoint mean is a permitted intercept plus a centered member of \(\mathcal H_{d,s}\), with squared midpoint mean integrating to at most one under the covariate law.

Because this class contains the uniform Gaussian completions, its minimax scaled excess satisfies
\[
R_s(n,d)
\leq
\inf_{\pi\in\mathcal D_{n,d,s}}
\sup_{P\in\mathcal A_{d,s}}\mathcal V_n(\pi,P),
\]
where \(\mathcal V_n(\pi,P)\) is excess above the law-specific benchmark specified in \cref{obj:thm:log-free-bivariate-capacity}. Thus a uniform precision requirement over the broader class must accommodate the uniform-sampling lower bound. In particular, for fixed \(s\) and fixed density bounds, vanishing minimax scaled excess over these classes requires \(d_n=o(\sqrt n)\). This is the necessary interaction-count condition inherited from the uniform subclass.

\subsection{Limitations and future work}

Attainment under nonuniform covariate densities is a natural next question. The two-sided characterization and the specified assignment guarantee apply to independent uniform sampling, while the bounded-density comparison supplies a converse. Extending attainment would require assignment and approximation bounds that account for the covariate density.

A second direction concerns prognostic functions beyond the exact centered order-two decomposition and pooled budget in \cref{obj:ass:order-two,obj:ass:pooled-budget}. Higher-order interactions, approximation error in the decomposition, and unknown prognostic intercepts call for explicit control of their contributions to imbalance. Such extensions could clarify how sample-size requirements change when pairwise structure is an approximation to the outcome regression.

Finally, the finite construction in \cref{obj:def:ordered-boundary-rounding} uses exact-real arithmetic conventions. Finite-precision certification and bit complexity remain implementation questions. Quantifying how numerical error affects boundary decisions, fairness, and row imbalance would connect the mathematical assignment guarantee to a certified numerical procedure.


\appendix

\section*{Appendices}

\section{Causal variance identities and class transfer}

The connection between prognostic balance and causal precision begins with the realized full potential-outcome schedule. Conditional fairness and assignment independence, as specified in \cref{obj:ass:fair,obj:ass:assignment-independence}, make the Horvitz--Thompson estimator conditionally unbiased for the finite-population contrast \citep{Horvitz1952}. Its conditional variance is determined by the potential-outcome midpoints and the covariance of the assignment signs. The following statement also gives the sampling-law variance identity for the Gaussian completion in \cref{obj:def:causal-completion}.

\begin{lemmav}{lem:ht-transfer}[Conditional HT variance transfer]*
Let the following conditions hold:
\begin{itemize}
\item \textbf{(Dimensions and smoothness.)} \(n,d\) are integers with \(n\geq2\), \(d\geq2\), and \(0<s\leq1\).
\item \textbf{(Prognosis.)} \(m\in\mathcal H_{d,s}\), as defined in \cref{obj:def:sobolev-class}.
\item \textbf{(Assignment.)} \(\pi\) is a covariate-dependent Borel assignment probability kernel satisfying \cref{obj:ass:fair,obj:ass:assignment-independence}.
\end{itemize}
Let \(X\), the full covariate sample, consist of \(n\) independent uniform draws from \([0,1]^d\), and let \(Z\in\{-1,1\}^n\), the assignment sign vector, have conditional law \(\pi(X)\). For a real potential-outcome schedule \(Y=(Y_i(0),Y_i(1))_{i=1}^n\), define the HT estimator, finite-population contrast, midpoint vector, and conditional assignment covariance by
\[
\widehat\theta
 =\frac2n\sum_{i=1}^n Z_iY_i\!\left(\frac{Z_i+1}{2}\right),
\qquad
\theta_n=\frac1n\sum_{i=1}^n\bigl(Y_i(1)-Y_i(0)\bigr),
\qquad
q_i=\frac{Y_i(1)+Y_i(0)}2,
\qquad
[C_\pi(X)]_{ij}=\operatorname{Cov}_\pi(Z_i,Z_j\mid X).
\]
For almost every \(X\), simultaneously for every such schedule \(Y\),
\[
\mathbb E_\pi[\widehat\theta\mid X,Y]=\theta_n,
\qquad
n\operatorname{Var}_\pi(\widehat\theta\mid X,Y)
 =\frac4n q^{\mathsf T}C_\pi(X)q.
\]
Define the scaled expected squared prognostic imbalance by
\[
\mathcal L_{n,d,s}(\pi,m)
 =\frac4n\,\mathbb E_{X,Z}
   \left[\left(\sum_{i=1}^n Z_i m(X_i)\right)^2\right].
\]
Under the frozen Gaussian completion associated with \(m\), with assignment law \(\pi(X)\), the sampling-law causal target and scaled variance benchmark satisfy
\[
\theta=\mathbb E_{X,Y}[\theta_n]=0.2,
\qquad
V^*=4.01.
\]
Moreover,
\[
\mathcal L_{n,d,s}(\pi,m)<\infty,
\qquad
n\operatorname{Var}_{X,Y,Z}(\widehat\theta)-V^*
 =\mathcal L_{n,d,s}(\pi,m).
\]
\end{lemmav}
% CAUSALSMITH-CITED-SCOPE-BEGIN lem:ht-transfer
\verificationfootnotetext{\textbf{Formalization scope.} This result uses the cited conclusion \citep{BartlettMendelson2002Complexities} (Theorem 12(4), with Rademacher complexity as in Definition 2). The cited source proof is not formalized here: Lean verifies its use and all remaining steps. If that cited conclusion is false, the portions of this result that depend on it are not certified.}
% CAUSALSMITH-CITED-SCOPE-END lem:ht-transfer

\begin{proof}[Proof of \cref{obj:lem:ht-transfer}]
For the cited expected absolute-supremum contraction conclusion, we invoke \citep[Theorem 12(4), with the convention of Definition 2]{BartlettMendelson2002Complexities}.
The calculation below uses the measurable, centered, square-integrable prognosis supplied by \cref{obj:def:sobolev-class}. Since \(n,d\geq2\), all divisions by \(n\) and first-coordinate evaluations are well defined.

\begin{enumerate}
\item
Choose \(X\) in the probability-one set on which every fairness equality in \cref{obj:ass:fair} holds. This set is chosen independently of the schedule \(Y\). For either value of \(Z_i\),
\[
2Z_iY_i\!\left(\frac{Z_i+1}{2}\right)
=
Y_i(1)-Y_i(0)+2Z_iq_i.
\]
Summing gives the pointwise identity
\[
\widehat\theta
=
\theta_n+\frac2n\sum_{i=1}^n q_iZ_i.
\]
Under \cref{obj:ass:assignment-independence}, assignment conditional on \((X,Y)\) uses the kernel \(\pi(X)\). Consequently,
\[
\mathbb E_\pi[\widehat\theta\mid X,Y]
=
\theta_n+\frac2n\sum_{i=1}^n
q_i\mathbb E_\pi[Z_i\mid X]
=
\theta_n.
\]
Expanding the centered square and using fairness in the definition of the assignment covariance yields
\[
\begin{aligned}
\operatorname{Var}_\pi(\widehat\theta\mid X,Y)
&=
\frac4{n^2}
\mathbb E_\pi\!\left[
\left(\sum_{i=1}^n q_iZ_i\right)^2
\,\middle|\,X\right]\\
&=
\frac4{n^2}\sum_{i=1}^n\sum_{j=1}^n
q_i\mathbb E_\pi[Z_iZ_j\mid X]q_j\\
&=
\frac4{n^2}q^{\mathsf T}C_\pi(X)q.
\end{aligned}
\]
These calculations apply to every real schedule on the same probability-one set, proving both conditional assertions. % lean: CausalSmith.Experimentation.BivariateSobolevDesignCapacity.ht_conditional_transfer

\item
Define the Gaussian completion's treatment-effect function and its sample average by
\[
\begin{aligned}
\tau_{\mathrm G}(x)
&=
0.2+0.1\sqrt2\sin(2\pi x_1),
&&x\in[0,1]^d,\\
\overline\tau_{\mathrm G}(X)
&=
\frac1n\sum_{i=1}^n\tau_{\mathrm G}(X_i),
&&X\in([0,1]^d)^n.
\end{aligned}
\]
The completion in \cref{obj:synth_4} gives
\[
\theta_n
=
\frac1n\sum_{i=1}^n
\bigl(\tau_{\mathrm G}(X_i)+\varepsilon_{i1}-\varepsilon_{i0}\bigr),
\qquad
\mathbb E_Y[\theta_n\mid X]
=
\overline\tau_{\mathrm G}(X).
\]
Each first coordinate is uniform on \([0,1]\), and integration over a full sine period gives
\[
\int_0^1\sin(2\pi v_{\mathrm{int}})\,dv_{\mathrm{int}}
=
\left[-\frac{\cos(2\pi v_{\mathrm{int}})}{2\pi}\right]_0^1
=
0.
\]
Thus
\[
\mathbb E_X[\tau_{\mathrm G}(X_i)]=0.2,
\qquad
\theta
=
\mathbb E_{X,Y}[\theta_n]
=
\mathbb E_X[\overline\tau_{\mathrm G}(X)]
=
0.2.
\]
The frozen benchmark evaluates to
\[
V^*=(0.1)^2+4=4.01.
\]
% lean: CausalSmith.Experimentation.BivariateSobolevDesignCapacity.gaussianTheta_eq_tau0

\item
For every covariate array and sign vector, define
\[
S_m(X,Z)=\sum_{i=1}^n Z_i m(X_i),
\qquad
(X,Z)\in([0,1]^d)^n\times\{-1,1\}^n.
\]
Cauchy--Schwarz and \(Z_i^2=1\) give the pointwise bound
\[
S_m(X,Z)^2
\leq
n\sum_{i=1}^n m(X_i)^2.
\]
Because \(\pi(X)\) is a probability law,
\[
0\leq
\mathbb E_\pi[S_m(X,Z)^2\mid X]
\leq
n\sum_{i=1}^n m(X_i)^2.
\]
Integrating against the uniform sample law therefore gives
\[
\mathbb E_{X,Z}[S_m(X,Z)^2]
\leq
n^2\int_{[0,1]^d}m(x)^2\,dx,
\]
and hence
\[
0\leq\mathcal L_{n,d,s}(\pi,m)
\leq
4n\int_{[0,1]^d}m(x)^2\,dx
<\infty.
\]
The same bound ensures integrability of the conditional imbalance second moment. % lean: CausalSmith.Experimentation.BivariateSobolevDesignCapacity.loss_lt_top_of_memLp

\item
Write \(\varepsilon=(\varepsilon_{i0},\varepsilon_{i1})_{i=1}^n\) for the Gaussian noise array in \cref{obj:synth_4}, and define the selected-noise contribution on its full domain by
\[
N_{\mathrm{sel}}(Z,\varepsilon)
=
\frac2n\sum_{i=1}^n
Z_i\varepsilon_{i,(Z_i+1)/2},
\qquad
(Z,\varepsilon)\in\{-1,1\}^n\times\mathbb R^{n\times2}.
\]
Substituting the completion, for either assignment sign, gives
\[
2Z_iY_i\!\left(\frac{Z_i+1}{2}\right)
=
\tau_{\mathrm G}(X_i)
+2Z_i m(X_i)
+2Z_i\varepsilon_{i,(Z_i+1)/2}.
\]
Consequently,
\[
\widehat\theta
=
\overline\tau_{\mathrm G}(X)
+\frac2n S_m(X,Z)
+N_{\mathrm{sel}}(Z,\varepsilon).
\]
For fixed \(Z\), the signed selected errors are independent across units, each with mean zero and variance one. With \(\mathbb E_\varepsilon\) denoting expectation over the Gaussian noise at fixed \(X,Z\), it follows that
\[
\begin{aligned}
\mathbb E_\varepsilon[N_{\mathrm{sel}}(Z,\varepsilon)]
&=
\frac2n\sum_{i=1}^n
\mathbb E_\varepsilon[
Z_i\varepsilon_{i,(Z_i+1)/2}]
=0,\\
\operatorname{Var}_\varepsilon(N_{\mathrm{sel}}(Z,\varepsilon))
&=
\frac4{n^2}\sum_{i=1}^n1
=\frac4n,\\
\mathbb E_\varepsilon[N_{\mathrm{sel}}(Z,\varepsilon)^2]
&=\frac4n.
\end{aligned}
\]
Thus
\[
\mathbb E_\varepsilon[\widehat\theta]
=
\overline\tau_{\mathrm G}(X)+\frac2n S_m(X,Z),
\]
and expansion of the square gives
\[
\begin{aligned}
\mathbb E_\varepsilon[(\widehat\theta-0.2)^2]
&=
\left(
\overline\tau_{\mathrm G}(X)-0.2+\frac2n S_m(X,Z)
\right)^2\\
&\quad+
2\left(
\overline\tau_{\mathrm G}(X)-0.2+\frac2n S_m(X,Z)
\right)
\mathbb E_\varepsilon[N_{\mathrm{sel}}(Z,\varepsilon)]\\
&\quad+
\mathbb E_\varepsilon[N_{\mathrm{sel}}(Z,\varepsilon)^2]\\
&=
\left(
\overline\tau_{\mathrm G}(X)-0.2+\frac2n S_m(X,Z)
\right)^2+\frac4n.
\end{aligned}
\]
% lean: CausalSmith.Experimentation.BivariateSobolevDesignCapacity.ht_completion_noise_centered_second_moment

\item
For almost every \(X\), fairness gives
\[
\mathbb E_\pi[S_m(X,Z)\mid X]
=
\sum_{i=1}^n m(X_i)\mathbb E_\pi[Z_i\mid X]
=
0.
\]
The finite assignment space permits exchanging the assignment average and the Gaussian-noise integral. Averaging the preceding noise mean therefore yields
\[
\mathbb E_{Y,Z}[\widehat\theta\mid X]
=
\overline\tau_{\mathrm G}(X).
\]
For the centered second moment, expand
\[
\begin{aligned}
\left(
\overline\tau_{\mathrm G}(X)-0.2+\frac2n S_m(X,Z)
\right)^2
&=
(\overline\tau_{\mathrm G}(X)-0.2)^2\\
&\quad+
\frac4n(\overline\tau_{\mathrm G}(X)-0.2)S_m(X,Z)
+\frac4{n^2}S_m(X,Z)^2.
\end{aligned}
\]
The middle term has assignment mean zero, so
\[
\mathbb E_{Y,Z}[(\widehat\theta-0.2)^2\mid X]
=
(\overline\tau_{\mathrm G}(X)-0.2)^2
+\frac4{n^2}\mathbb E_\pi[S_m(X,Z)^2\mid X]
+\frac4n.
\]
The bounded treatment-effect function, the integrable imbalance bound, and Gaussian second moments ensure that this expression is integrable over \(X\), and that the actual estimator has a finite centered second moment. % lean: CausalSmith.Experimentation.BivariateSobolevDesignCapacity.ht_gaussian_kernel_centered_second_moment

\item
Integrating the conditional mean over the covariates gives
\[
\mathbb E_{X,Y,Z}[\widehat\theta]
=
\mathbb E_X[\overline\tau_{\mathrm G}(X)]
=
0.2.
\]
Square integrability and the iterated construction of the experiment law now imply
\[
\operatorname{Var}_{X,Y,Z}(\widehat\theta)
=
\mathbb E_{X,Y,Z}[(\widehat\theta-0.2)^2]
=
\mathbb E_X\!\left[
\mathbb E_{Y,Z}[(\widehat\theta-0.2)^2\mid X]
\right].
\]
% lean: CausalSmith.Experimentation.BivariateSobolevDesignCapacity.ht_gaussian_joint_mean

\item
The squared-sine integral is
\[
\begin{aligned}
\int_0^1\sin^2(2\pi v_{\mathrm{int}})\,dv_{\mathrm{int}}
&=
\int_0^1\frac{1-\cos(4\pi v_{\mathrm{int}})}2\,dv_{\mathrm{int}}\\
&=
\left[
\frac{v_{\mathrm{int}}}{2}
-\frac{\sin(4\pi v_{\mathrm{int}})}{8\pi}
\right]_0^1
=\frac12.
\end{aligned}
\]
Together with the treatment-effect mean already computed, this gives
\[
\operatorname{Var}_X(\tau_{\mathrm G}(X_i))
=
(0.1)^2\,2\int_0^1
\sin^2(2\pi v_{\mathrm{int}})\,dv_{\mathrm{int}}
=
(0.1)^2.
\]
Independence of the original covariate vectors consequently yields
\[
\begin{aligned}
\mathbb E_X[
(\overline\tau_{\mathrm G}(X)-0.2)^2]
&=
\operatorname{Var}_X(\overline\tau_{\mathrm G}(X))\\
&=
\frac1{n^2}\sum_{i=1}^n
\operatorname{Var}_X(\tau_{\mathrm G}(X_i))
=
\frac{(0.1)^2}{n}.
\end{aligned}
\]
By the definition of the finite signing loss,
\[
\frac4{n^2}
\mathbb E_X\!\left[
\mathbb E_\pi[S_m(X,Z)^2\mid X]
\right]
=
\frac{\mathcal L_{n,d,s}(\pi,m)}{n}.
\]
Integrating the conditional centered second moment thus gives
\[
\operatorname{Var}_{X,Y,Z}(\widehat\theta)
=
\frac{(0.1)^2}{n}
+\frac{\mathcal L_{n,d,s}(\pi,m)}{n}
+\frac4n.
\]
Multiplying by \(n\) and subtracting the evaluated benchmark proves
\[
n\operatorname{Var}_{X,Y,Z}(\widehat\theta)-V^*
=
\mathcal L_{n,d,s}(\pi,m).
\]
% lean: CausalSmith.Experimentation.BivariateSobolevDesignCapacity.ht_gaussian_outer_centered_second_moment
\end{enumerate}
\end{proof}

The midpoint vector \leanref{sym:q}{\ensuremath{q}} captures the schedule's contribution to assignment variance, while the conditional assignment covariance \leanref{sym:Cpi}{\ensuremath{C_\pi(X)}} records the dependence induced by the design. The simultaneous quantifier in \cref{obj:lem:ht-transfer} makes this relationship a fixed-schedule identity: a single probability-one set of covariate samples supports the conclusion for every real schedule. Conditional unbiasedness targets the realized contrast \(\theta_n\); under the specified Gaussian completion, the sampling-law calculation instead averages over covariates, potential outcomes, and assignment and has target \(\theta=0.2\).

For this completion, \(V^*=4.01\) supplies a numerical illustration of the causal interpretation of signing loss. Thus the assignment guarantee in \cref{obj:thm:log-free-bivariate-capacity} controls the actual scaled HT variance above that benchmark. The corresponding transfer for arbitrary square-integrable outcome laws uses each law's own benchmark. The relevant published excess-variance identity is \citet[Proposition 2.3, equation (2.4), under Assumption 2.1 and Definition 2.2]{Cytrynbaum2026Limits}. The next statement specifies that benchmark, characterizes the attainable prognostic functions, and identifies the resulting minimax criterion.

\begin{lemmav}{lem:published-prognostic-capacity}[Published prognostic image and capacity]*
For a single-unit law \(P\), let \(U\in[0,1]^d\) be its covariate and \(O(0),O(1)\) its potential outcomes. Write
\[
m_P(u)=\mathbb E_P\!\left[\frac{O(1)+O(0)}2\mid U=u\right],
\qquad
\vartheta(P)=\mathbb E_P[O(1)-O(0)],
\]
using a finite-valued Borel version of the conditional mean. Let \(\mathcal Q_{d,s}\) consist of laws with uniform covariates, finite potential-outcome second moments, and \(m_P\) equal almost everywhere to a member of the centered prognostic class \(\mathcal H_{d,s}\) appearing in \cref{obj:def:criterion}. For independent copies across the original \(n\) units, define
\[
\widehat T_P=\frac2n\sum_{i=1}^n Z_iO_i((Z_i+1)/2),
\]
\[
\mathsf V(P)=
\operatorname{Var}_P\!\left(\mathbb E_P[O(1)-O(0)\mid U]\right)
+2\mathbb E_P\operatorname{Var}_P(O(1)\mid U)
+2\mathbb E_P\operatorname{Var}_P(O(0)\mid U),
\qquad
\mathcal V_n(\pi,P)=n\operatorname{Var}_{P,\pi}(\widehat T_P)-\mathsf V(P).
\]
The scaled signing loss is
\[
\mathcal L_{n,d,s}(\pi,m)
=\frac4n\mathbb E_{\pi}\left(\sum_{i=1}^n Z_i m(X_i)\right)^2,
\]
where \(X_1,\ldots,X_n\) are independent uniform covariate vectors and \(Z\) is drawn from the original-sample assignment kernel \(\pi\).

For every integer \(d\geq2\) and \(0<s\leq1\), the prognostic image is exactly the Sobolev image, with both images saturated under Lebesgue null representatives:
\[
\left\{f:\exists P\in\mathcal Q_{d,s},\ f=m_P\ \text{almost everywhere}\right\}
=
\left\{f:\exists m\in\mathcal H_{d,s},\ f=m\ \text{almost everywhere}\right\}.
\]
The stated Gaussian completions of members of \(\mathcal H_{d,s}\) form a proper subclass of \(\mathcal Q_{d,s}\): the specified two-point outcome law belongs to \(\mathcal Q_{d,s}\) and lies outside that Gaussian subclass.

For every \(n,d\geq2\), a Borel probability assignment kernel satisfies published admissibility precisely when it belongs to the design class in \cref{obj:def:design-class}. Here published admissibility means conditional mean-zero assignment signs almost everywhere under the original uniform covariate law, together with conditional independence of the assignment signs and the potential-outcome array given the observed covariates, for every single-unit law with bounded strictly positive covariate density and finite potential-outcome second moments.

For every \(n,d\geq2\), \(0<s\leq1\), and \(\pi\in\mathcal D_{n,d,s}\),
\[
0\leq\mathcal V_n(\pi,P)
=\mathcal L_{n,d,s}(\pi,m_P),
\qquad P\in\mathcal Q_{d,s},
\]
and
\[
\sup_{P\in\mathcal Q_{d,s}}\mathcal V_n(\pi,P)
=
\sup_{m\in\mathcal H_{d,s}}\mathcal L_{n,d,s}(\pi,m).
\]
Consequently, the full-design capacity agrees with \cref{obj:def:criterion}:
\[
\inf_{\pi\in\mathcal D_{n,d,s}}
\sup_{P\in\mathcal Q_{d,s}}\mathcal V_n(\pi,P)
=R_s(n,d).
\]

Fix comparison parameters satisfying
\begin{itemize}
\item \textbf{(Smoothness.)} \(0<s\leq1\).
\item \textbf{(Lower density bound.)} \(0<\lambda\leq1\).
\item \textbf{(Upper density bound.)} \(1\leq\Lambda<\infty\).
\end{itemize}
Let \(\mathcal A_{d,s}\) be the correctly specified, no-priority-block, order-two class of \citet[Theorem B.8(ii)]{Cytrynbaum2026Limits}. Its laws have finite potential-outcome second moments, covariate densities satisfying \(\lambda\leq f(u)\leq\Lambda\), and prognostic functions satisfying
\[
\int m_P(u)^2\,dP_U(u)\leq1,
\qquad
m_P(u)=\beta+m(u)\quad\text{for almost every }u,
\]
for a permitted intercept \(\beta\) and a centered \(m\in\mathcal H_{d,s}\) under the same component norm convention. The lower comparison constants and saturated lower comparison scale satisfy
\[
0<c_1\leq c_s,
\qquad
b_s(n,d)=\min\{1,(\tbinom d2/n)^s\}.
\]
and, for every \(n,d\geq2\),
\[
c_s b_s(n,d)
\leq R_s(n,d)
\leq
\inf_{\pi\in\mathcal D_{n,d,s}}
\sup_{P\in\mathcal A_{d,s}}\mathcal V_n(\pi,P).
\]
For every integer sequence \(d_n\geq2\), vanishing capacity on this bounded-density class implies
\[
\inf_{\pi\in\mathcal D_{n,d_n,s}}
\sup_{P\in\mathcal A_{d_n,s}}\mathcal V_n(\pi,P)
\longrightarrow0
\quad\Longrightarrow\quad
d_n=o(\sqrt n).
\]
\end{lemmav}
% CAUSALSMITH-CITED-SCOPE-BEGIN lem:published-prognostic-capacity
\verificationfootnotetext{\textbf{Formalization scope.} This result uses the cited conclusion \citep{Cytrynbaum2026Limits} (Proposition 2.3 (Excess Variance), equation (2.4), with Assumption 2.1 and Definition 2.2) and \citep{BartlettMendelson2002Complexities} (Theorem 12(4), with Rademacher complexity as in Definition 2). The cited source proof is not formalized here: Lean verifies its use and all remaining steps. If that cited conclusion is false, the portions of this result that depend on it are not certified.}
% CAUSALSMITH-CITED-SCOPE-END lem:published-prognostic-capacity

\begin{proof}[Proof of \cref{obj:lem:published-prognostic-capacity}]
\begin{enumerate}
\item Fix \(d\geq2\) and \(0<s\leq1\). Every law in \(\mathcal Q_{d,s}\) has prognosis equal almost everywhere to a member of \(\mathcal H_{d,s}\), so transitivity of almost-everywhere equality gives the forward inclusion of the stated images.

For the reverse inclusion, use the Gaussian completion in \cref{obj:synth_4}. Its single-unit law can be written as
\[
\begin{gathered}
\tau(u)=0.2+0.1\sqrt2\sin(2\pi u_1),
\qquad u\in[0,1]^d,\\
\operatorname{Law}(U,\varepsilon_0,\varepsilon_1)
=\operatorname{Unif}([0,1]^d)\otimes
\mathcal N(0,1)\otimes\mathcal N(0,1),\\
P_m=\operatorname{Law}\!\left(
U,\,
m(U)-\frac{\tau(U)}2+\varepsilon_0,\,
m(U)+\frac{\tau(U)}2+\varepsilon_1
\right),
\qquad m\in\mathcal H_{d,s}.
\end{gathered}
\]
The function \(m\) is square integrable, \(\tau\) is bounded, and both Gaussian errors are square integrable. Closure of \(L^2\) under finite linear combinations therefore gives
\[
\mathbb E_{P_m}[O(0)^2]<\infty,
\qquad
\mathbb E_{P_m}[O(1)^2]<\infty.
\]
The displayed construction is a conditional outcome law given \(U\). Its conditional midpoint mean is
\[
m_{P_m}(u)
=m(u)+\frac{\mathbb E[\varepsilon_0]+\mathbb E[\varepsilon_1]}2
=m(u)
\quad\text{for almost every }u.
\]
Consequently \(P_m\in\mathcal Q_{d,s}\). Every null representative of \(m\) is thus a null representative of \(m_{P_m}\), proving the reverse image inclusion and the asserted saturation.
% lean: CausalSmith.Experimentation.BivariateSobolevDesignCapacity.publishedPrognosticImage_eq_sobolevImage

\item Define the Gaussian subclass and the two-point witness by
\[
\begin{gathered}
\mathcal Q^{\mathrm G}_{d,s}
=\{P_m:m\in\mathcal H_{d,s}\},\\
\operatorname{Law}(U,\zeta_{\mathrm{tp}})
=\operatorname{Unif}([0,1]^d)\otimes
\frac{\delta_{-1}+\delta_1}{2},
\qquad
P_{\mathrm{tp}}=\operatorname{Law}(U,\zeta_{\mathrm{tp}},\zeta_{\mathrm{tp}}).
\end{gathered}
\]
The preceding step gives
\[
\mathcal Q^{\mathrm G}_{d,s}\subseteq\mathcal Q_{d,s}.
\]
For the witness,
\[
m_{P_{\mathrm{tp}}}(u)=\mathbb E[\zeta_{\mathrm{tp}}]=0,
\qquad
\mathbb E_{P_{\mathrm{tp}}}[O(0)^2]
=\mathbb E_{P_{\mathrm{tp}}}[O(1)^2]=1.
\]
The zero function has zero canonical components and zero pooled component budget, hence belongs to \(\mathcal H_{d,s}\). Thus \(P_{\mathrm{tp}}\in\mathcal Q_{d,s}\).

The outcome diagonal distinguishes this law from every Gaussian completion:
\[
P_{\mathrm{tp}}\{O(0)=O(1)\}=1.
\]
Indeed, in the Gaussian construction, conditioning on \(U\) and \(\varepsilon_0\) gives
\[
\Pr_{P_m}\!\left(O(0)=O(1)\mid U,\varepsilon_0\right)
=
\Pr\!\left(\varepsilon_1=\varepsilon_0-\tau(U)\mid U,\varepsilon_0\right)
=0,
\]
because the independent Gaussian error \(\varepsilon_1\) assigns probability zero to every singleton. Integrating yields
\[
P_m\{O(0)=O(1)\}=0.
\]
Therefore
\[
P_{\mathrm{tp}}\notin\mathcal Q^{\mathrm G}_{d,s},
\qquad
\mathcal Q^{\mathrm G}_{d,s}\subsetneq\mathcal Q_{d,s}.
\]
% lean: CausalSmith.Experimentation.BivariateSobolevDesignCapacity.gaussianLawClass_ssubset_frozenUniformClass

\item\label{proof:published-prognostic-capacity:excess} Fix \(n,d\geq2\) and \(\pi\in\mathcal D_{n,d,s}\). We first handle its almost-everywhere fairness. For a deterministic covariate array \(\mathbf x\in([0,1]^d)^n\), define
\[
\pi^{\mathrm{fair}}(\mathbf x)=
\begin{cases}
\pi(\mathbf x),
&
\displaystyle
\sum_{z\in\mathcal Z_n}z_i\pi(\mathbf x,\{z\})=0
\quad\text{for every }i,\\[4pt]
\displaystyle
\bigotimes_{i=1}^n\frac{\delta_{-1}+\delta_1}{2},
&\text{otherwise}.
\end{cases}
\]
Each displayed conditional mean is a finite Borel sum. This defines a Borel probability kernel that is fair at every covariate array and satisfies
\[
\pi^{\mathrm{fair}}(\mathbf x)=\pi(\mathbf x)
\quad\text{for }
\operatorname{Unif}([0,1]^d)^{\otimes n}\text{-almost every }\mathbf x.
\]

A covariate-based kernel also supplies conditional assignment independence. More explicitly, for an arbitrary covariate law and pre-assignment outcome kernel with domains
\[
\mu_X\in\mathcal P\bigl(([0,1]^d)^n\bigr),
\qquad
\kappa:([0,1]^d)^n\longrightarrow\mathcal P(\mathbb R^{n\times2}),
\]
the experiment is
\[
\operatorname{Law}(X,Y,Z)(d\mathbf x,dY,dz)
=\mu_X(d\mathbf x)\,\kappa(\mathbf x,dY)\,\pi(\mathbf x,dz).
\]
Its conditional distributions satisfy
\[
\operatorname{Law}(Z\mid X,Y)=\pi(X)
=\operatorname{Law}(Z\mid X),
\qquad
\operatorname{Law}(Y,Z\mid X)=\kappa(X)\otimes\pi(X).
\]
The first equality follows from the construction; integrating out \(Y\) gives the second conditional distribution. Their equality gives conditional independence, also for \(\pi^{\mathrm{fair}}\).

For any single-unit law \(P\) with bounded strictly positive covariate density and finite outcome second moments, the independent covariate sample law satisfies
\[
\operatorname{Law}_{P}\bigl((U_i)_{i=1}^n\bigr)
\ll \operatorname{Unif}([0,1]^d)^{\otimes n}.
\]
Hence the replacement preserves the entire experiment law:
\[
\operatorname{Law}_{P,\pi}\!\left(
(U_i,O_i(0),O_i(1))_{i=1}^n,Z\right)
=
\operatorname{Law}_{P,\pi^{\mathrm{fair}}}\!\left(
(U_i,O_i(0),O_i(1))_{i=1}^n,Z\right).
\]
Apply the published excess-variance identity
\citep[Proposition 2.3, equation (2.4), with Assumption 2.1 and Definition 2.2]{Cytrynbaum2026Limits}
to the pointwise fair replacement and transfer it through this equality. It gives
\[
\mathcal V_n(\pi,P)
=
\frac4n\,
\mathbb E_{P,\pi}
\left(\sum_{i=1}^n Z_i m_P(U_i)\right)^2
\geq0.
\]

Now take \(P\in\mathcal Q_{d,s}\), and choose \(m\in\mathcal H_{d,s}\) with \(m_P=m\) almost everywhere. Independence and uniformity of the unit covariates give
\[
\operatorname{Law}_{P,\pi}\bigl((U_i)_{i=1}^n,Z\bigr)
=
\operatorname{Unif}([0,1]^d)^{\otimes n}(dX)\,\pi(X,dZ).
\]
Moreover, the finite sample avoids the exceptional null set simultaneously:
\[
\Pr_{P,\pi}\!\left(
m_P(U_i)=m(U_i)\text{ for every }i
\right)=1.
\]
Square integrability of \(m\) makes the finite signed sum square integrable. Substituting the preceding two identities yields
\[
0\leq\mathcal V_n(\pi,P)
=\mathcal L_{n,d,s}(\pi,m_P)
=\mathcal L_{n,d,s}(\pi,m).
\]
This also proves invariance of the signing loss under the null representatives used here.
% lean: CausalSmith.Experimentation.BivariateSobolevDesignCapacity.published_uniform_excess_eq_loss

\item For each fixed admissible \(\pi\), the forward prognostic-image inclusion and the preceding identity give
\[
\sup_{P\in\mathcal Q_{d,s}}\mathcal V_n(\pi,P)
\leq
\sup_{m\in\mathcal H_{d,s}}\mathcal L_{n,d,s}(\pi,m).
\]
Conversely, the image equality supplies, for every \(m\in\mathcal H_{d,s}\), a law \(P\in\mathcal Q_{d,s}\) with \(m=m_P\) almost everywhere. For that law,
\[
\mathcal L_{n,d,s}(\pi,m)
=\mathcal V_n(\pi,P)
\leq
\sup_{P\in\mathcal Q_{d,s}}\mathcal V_n(\pi,P).
\]
Taking the function supremum proves
\[
\sup_{P\in\mathcal Q_{d,s}}\mathcal V_n(\pi,P)
=
\sup_{m\in\mathcal H_{d,s}}\mathcal L_{n,d,s}(\pi,m).
\]
% lean: CausalSmith.Experimentation.BivariateSobolevDesignCapacity.publishedWorst_eq_of_image_and_excess

\item The kernel construction in \cref{proof:published-prognostic-capacity:excess} establishes the independence requirement of \cref{obj:ass:assignment-independence} for every Borel probability assignment kernel. Thus a published-admissible kernel, which already satisfies fairness under the original uniform sample, satisfies both requirements of \cref{obj:def:design-class}.

Conversely, a kernel in that design class has the required fairness. In the experiment generated from independent copies of any single-unit law, drawing signs from \(\pi(X)\) gives
\[
\operatorname{Law}\!\left(
Z\mid X,(O_i(0),O_i(1))_{i=1}^n
\right)
=\pi(X)
=\operatorname{Law}(Z\mid X).
\]
This proves the published conditional-independence requirement throughout its stated law class. Consequently,
\[
\pi\text{ satisfies published admissibility}
\quad\Longleftrightarrow\quad
\pi\in\mathcal D_{n,d,s}.
\]
% lean: CausalSmith.Experimentation.BivariateSobolevDesignCapacity.published_admissible_iff_designClass

\item\label{proof:published-prognostic-capacity:capacity} Taking the infimum of the fixed-design equality over the common design class gives
\[
\inf_{\pi\in\mathcal D_{n,d,s}}
\sup_{P\in\mathcal Q_{d,s}}\mathcal V_n(\pi,P)
=
\inf_{\pi\in\mathcal D_{n,d,s}}
\sup_{m\in\mathcal H_{d,s}}\mathcal L_{n,d,s}(\pi,m)
=R_s(n,d),
\]
where the last equality is \cref{obj:def:criterion}.
% lean: CausalSmith.Experimentation.BivariateSobolevDesignCapacity.publishedCapacity_eq_of_worst

\item\label{proof:published-prognostic-capacity:comparison} Fix \(0<s\leq1\) and density bounds \(0<\lambda\leq1\leq\Lambda<\infty\). For \(P\in\mathcal Q_{d,s}\), choose \(m\in\mathcal H_{d,s}\) with \(m_P=m\) almost everywhere. Uniform density meets the comparison bounds:
\[
\frac{dP_U}{du}=1,
\qquad
\lambda\leq1\leq\Lambda.
\]
The final second-moment assertion of \cref{obj:lem:exact-reflection-budget} supplies
\[
\int m_P(u)^2\,dP_U(u)
=
\int_{[0,1]^d}m(u)^2\,du
\leq1.
\]
The permitted intercept can be chosen as zero:
\[
m_P(u)=0+m(u)\quad\text{almost everywhere}.
\]
Together with the finite outcome second moments, these statements verify every condition for membership in \(\mathcal A_{d,s}\). Therefore
\[
\mathcal Q_{d,s}\subseteq\mathcal A_{d,s}.
\]
For each admissible design, inclusion increases the law supremum. Taking infima and using \cref{proof:published-prognostic-capacity:capacity} gives
\[
R_s(n,d)
=
\inf_{\pi\in\mathcal D_{n,d,s}}
\sup_{P\in\mathcal Q_{d,s}}\mathcal V_n(\pi,P)
\leq
\inf_{\pi\in\mathcal D_{n,d,s}}
\sup_{P\in\mathcal A_{d,s}}\mathcal V_n(\pi,P).
\]
All excesses in these classes are nonnegative by the general identity in \cref{proof:published-prognostic-capacity:excess}.
% lean: CausalSmith.Experimentation.BivariateSobolevDesignCapacity.publishedCapacity_mono_class

\item With the expected absolute-supremum contraction inequality
\citep[Theorem 12(4), with Definition 2]{BartlettMendelson2002Complexities}
as its cited input, \cref{obj:thm:full-design-lower} supplies
\[
c_s b_s(n,d)\leq R_s(n,d),
\qquad
c_s=\frac{2}{(48+32\pi^2)16384^s},
\qquad
b_s(n,d)=\min\left\{1,\left(\frac{\binom d2}{n}\right)^s\right\}.
\]
The constant calibration is
\[
16384^s\leq16384,
\qquad
0<c_1=\frac{2}{(48+32\pi^2)16384}\leq c_s.
\]
Indeed, the power inequality follows from \(s\leq1\), and the denominator is positive. Combining the lower bound with \cref{proof:published-prognostic-capacity:comparison} proves
\[
c_s b_s(n,d)
\leq R_s(n,d)
\leq
\inf_{\pi\in\mathcal D_{n,d,s}}
\sup_{P\in\mathcal A_{d,s}}\mathcal V_n(\pi,P).
\]
% lean: CausalSmith.Experimentation.BivariateSobolevDesignCapacity.cLower_uniform_positive

\item Finally, let \(d_n\geq2\) for every \(n\), and suppose the displayed bounded-density capacities tend to zero. For \(n\geq2\), define
\[
\begin{gathered}
\mathfrak C_n
=
\inf_{\pi\in\mathcal D_{n,d_n,s}}
\sup_{P\in\mathcal A_{d_n,s}}\mathcal V_n(\pi,P),\\
B_n=\binom{d_n}{2},
\qquad
q_n=\frac{B_n}{n}.
\end{gathered}
\]
These quantities are used on the tail \(n\geq2\), which determines the limit. The comparison just proved gives
\[
0\leq c_s\min\{1,q_n^s\}\leq\mathfrak C_n.
\]
Since \(\mathfrak C_n\to0\) and \(c_s>0\), squeezing and division by \(c_s\) yield
\[
b_s(n,d_n)=\min\{1,q_n^s\}\longrightarrow0.
\]
Eventually this minimum is less than one, so on that tail
\[
q_n^s=b_s(n,d_n),
\qquad
q_n=b_s(n,d_n)^{1/s}\longrightarrow0,
\]
using \(s>0\).

The exact pair count and \(d_n\geq2\) imply
\[
B_n=\frac{d_n(d_n-1)}2,
\qquad
\frac{d_n^2}{4}\leq B_n\leq\frac{d_n^2}{2}.
\]
For the lower inequality, \(d_n-1\geq d_n/2\); for the upper inequality, \(d_n-1\leq d_n\). Hence
\[
0\leq\frac{d_n^2}{n}\leq4q_n\longrightarrow0.
\]
Taking square roots on the positive-\(n\) tail gives
\[
\frac{d_n}{\sqrt n}
=\sqrt{\frac{d_n^2}{n}}\longrightarrow0,
\]
which is \(d_n=o(\sqrt n)\).
% lean: CausalSmith.Experimentation.BivariateSobolevDesignCapacity.dimension_necessity_of_scale_lower
\end{enumerate}
\end{proof}

The benchmark \(\mathsf V(P)\) in \cref{obj:lem:published-prognostic-capacity} combines variation in the conditional mean treatment effect with the two conditional potential-outcome variances. It therefore adjusts to the outcome law before excess variance is compared across designs. With independent original units and admissible assignment, the remaining scaled variance depends on the law through its conditional mean midpoint \(m_P\). Finite second moments ensure that the estimator and benchmark are well defined.

Equality of prognostic images gives the transfer its full uniform-law scope. Every member of \(\mathcal H_{d,s}\) is attained by a Gaussian completion, and every law in \(\mathcal Q_{d,s}\) has a prognosis in that same class, modulo null sets. The two-point witness establishes that this outcome-law class also accommodates discrete potential outcomes. Together with the excess-variance identity, the image equality yields equality of the worst-case criteria for each admissible design and hence equality of the minimax capacities. Consequently, the uniform-law upper and lower comparisons concern the actual HT excess variance throughout \(\mathcal Q_{d,s}\).

The larger-class comparison uses the inclusion recorded in \cref{obj:prop:published-class-inclusion}. Under the common component norm convention, the uniform laws with centered prognosis satisfy the comparison class's prognostic restrictions with intercept zero; their density is one and satisfies the displayed density bounds. Thus \(\mathcal Q_{d,s}\subseteq\mathcal A_{d,s}\), and enlarging the outcome-law supremum preserves the uniform-class lower bound. For the class associated with \citet[Theorem B.8(ii)]{Cytrynbaum2026Limits}, \cref{obj:lem:published-prognostic-capacity} therefore supplies the saturated converse and the necessary dimension condition \(d_n=o(\sqrt n)\) for vanishing minimax scaled excess.


\section{Spectral budget and assignment bounds}

The attaining assignment in \cref{obj:def:capacity-handle} combines a componentwise regularity bound with ordered spectral balance. The exact reflection inequality preserves the pooled budget when the nonperiodic prognostic function is represented on the torus. The common shift then converts expected prognostic imbalance into a sum of frequency-specific second moments. Finally, the rounding guarantee controls those moments while preserving fair assignment on the original units.

\subsection{Exact reflection and the pooled budget}

For the canonical components in \cref{obj:synth_1}, introduce the appendix-local reflected functions
\[
h_j(y)=g_j(m)(b_1(y)),
\qquad
h_{j\ell}(y)=g_{j\ell}(m)(b_2(y)),
\]
on \(\mathbb T_2^1\) and \(\mathbb T_2^2\), respectively. Their component Fourier coefficients, with normalized Haar integration, are
\[
\widehat h_j(a)
=\int_{\mathbb T_2^1}h_j(y)e^{-\mathrm i\pi ay}\,dy,
\qquad a\in\mathbb Z,
\]
and
\[
\widehat h_{j\ell}(a)
=\int_{\mathbb T_2^2}h_{j\ell}(y)e^{-\mathrm i\pi a^{\mathsf T}y}\,dy,
\qquad a\in\mathbb Z^2.
\]
These coefficients connect each component's restriction norm to its contribution to the full reflected spectrum. The following statement records the exact normalization before the broader transfer result.

\begin{lemmav}{lem:exact-reflection-budget}[Exact reflection budget]
Let \(p\), \(s\), and the real-valued cube function \(g\) satisfy:
\begin{itemize}
\item \textbf{(Dimension.)} \(p\in\{1,2\}\).
\item \textbf{(Smoothness.)} \(0<s\leq1\).
\item \textbf{(Restriction norm.)} \(g\in L^2([0,1]^p)\), with a Borel representative, and \(N_{p,s}(g)^2<\infty\), where
\[
N_{p,s}(g)^2
=
\inf_{\substack{G\in L^1(\mathbb R^p;\mathbb C)\cap L^2(\mathbb R^p;\mathbb C)\\
G=g\ \text{a.e. on }[0,1]^p}}
\int_{\mathbb R^p}
\left(1+\frac{\|\omega\|_2^2}{p}\right)^s
\left|
(2\pi)^{-p/2}
\int_{\mathbb R^p}G(u)e^{-\mathrm i\omega^{\mathsf T}u}\,du
\right|^2\,d\omega .
\]
\end{itemize}
On the period-two torus \(\mathbb T_2^p=(\mathbb R/2\mathbb Z)^p\), define the coordinatewise reflection map and reflected function by
\[
(b_p(y))_j=\min\{y_j,2-y_j\},
\qquad y\in[0,2)^p,
\qquad h(y)=g(b_p(y)).
\]
For \(a\in\mathbb Z^p\), define the normalized Fourier coefficient
\[
\widehat h(a)
=
2^{-p}\int_{[0,2)^p}h(y)e^{-\mathrm i\pi a^{\mathsf T}y}\,dy .
\]
Then
\[
\sum_{a\in\mathbb Z^p}
\left(1+\frac{\pi^2\|a\|_2^2}{p}\right)^s
|\widehat h(a)|^2
\leq N_{p,s}(g)^2 .
\]

For every integer \(d\geq2\), every \(0<s\leq1\), and every \(m\in\mathcal H_{d,s}\) as defined in \cref{obj:def:sobolev-class}, let \(F\) be its reflected function on \(\mathbb T_2^d\), with
\[
F(y)=m(b_d(y)),
\qquad
(b_d(y))_j=\min\{y_j,2-y_j\},
\qquad y\in[0,2)^d,
\]
and define
\[
\widehat F(k)
=
2^{-d}\int_{[0,2)^d}F(y)e^{-\mathrm i\pi k^{\mathsf T}y}\,dy,
\qquad
\operatorname{supp}(k)=\{j:k_j\neq0\},
\qquad k\in\mathbb Z^d.
\]
For nonzero \(k\), put
\[
t(k)=\frac{\|k\|_2^2}{|\operatorname{supp}(k)|}.
\]
Then
\[
\widehat F(0)=0,
\qquad
\widehat F(k)=0
\quad\text{whenever }|\operatorname{supp}(k)|>2,
\]
and
\[
\sum_{\substack{k\in\mathbb Z^d\\1\leq|\operatorname{supp}(k)|\leq2}}
(1+\pi^2t(k))^s|\widehat F(k)|^2\leq1,
\qquad
\int_{[0,1]^d}|m(x)|^2\,dx\leq1.
\]
\end{lemmav}

\begin{proof}[Proof of \cref{obj:lem:exact-reflection-budget}]
We first prove the component estimate at \(s=1\), treating the one-dimensional and two-dimensional cases in that order, and then prove the estimate for \(0<s<1\).

\begin{enumerate}
\item \textit{Normalization and inverse truncations.}
For \(p\in\{1,2\}\), write
\[
d\mu_p(y)=2^{-p}\,dy,
\qquad
(\mathscr D_pG)(y)=G(b_p(y)),
\qquad
\widehat\zeta(a)=
\int_{\mathbb T_2^p}\zeta(y)e^{-\mathrm i\pi a^{\mathsf T}y}\,d\mu_p(y),
\]
where \(G\in L^2(\mathbb R^p;\mathbb C)\), \(\zeta\in L^2(\mathbb T_2^p;\mathbb C)\), and \(a\in\mathbb Z^p\). These operators act on equivalence classes. Indeed, for \(\epsilon\in\{0,1\}^p\), define the reflection branch
\[
\bigl(\mathscr R_\epsilon(x)\bigr)_j
=
\begin{cases}
x_j,&\epsilon_j=0,\\
2-x_j,&\epsilon_j=1,
\end{cases}
\qquad x\in[0,1]^p.
\]
Each branch preserves Lebesgue measure between the corresponding boxes, and
\(b_p(\mathscr R_\epsilon(x))=x\) almost everywhere. Consequently, for every nonnegative measurable \(\varphi\) on the cube,
\[
\int_{\mathbb T_2^p}\varphi(b_p(y))\,d\mu_p(y)
=
2^{-p}\sum_{\epsilon\in\{0,1\}^p}
\int_{[0,1]^p}\varphi(x)\,dx
=
\int_{[0,1]^p}\varphi(x)\,dx.
\]
In particular, folding preserves null sets in the required direction and
\[
\|\mathscr D_pG\|_{L^2(\mu_p)}^2
=
\int_{[0,1]^p}|G(u)|^2\,du
\leq
\|G\|_{L^2(\mathbb R^p)}^2.
\]
The normalized characters form a complete orthonormal Fourier basis, so
\[
\sum_{a\in\mathbb Z^p}|\widehat\zeta(a)|^2
=
\|\zeta\|_{L^2(\mu_p)}^2.
\]

Use the angular Fourier transform
\[
(\mathcal F_pG)(\omega)
=
(2\pi)^{-p/2}
\int_{\mathbb R^p}G(u)e^{-\mathrm i\omega^{\mathsf T}u}\,du
\]
on \(L^1\cap L^2\), with its unitary extension to \(L^2\).
For a measurable profile \(\psi\in L^2(\mathbb R^p;\mathbb C)\) and an integer
\(n_{\mathrm{cut}}\geq0\), define
\[
\psi_{n_{\mathrm{cut}}}(\omega)
=
\mathbf1_{\{\|\omega\|_2\leq n_{\mathrm{cut}}\}}\psi(\omega),
\qquad
G_{\psi,n_{\mathrm{cut}}}(u)
=
(2\pi)^{-p/2}
\int_{\mathbb R^p}
e^{\mathrm i\omega^{\mathsf T}u}
\psi_{n_{\mathrm{cut}}}(\omega)\,d\omega.
\]
Every frequency moment of the truncated profile is integrable. Explicitly, for every integer \(\ell_{\mathrm{mom}}\geq0\), Cauchy--Schwarz gives
\[
\int_{\mathbb R^p}
\|\omega\|_2^{\ell_{\mathrm{mom}}}
|\psi_{n_{\mathrm{cut}}}(\omega)|\,d\omega
\leq
\left(
\int_{\{\|\omega\|_2\leq n_{\mathrm{cut}}\}}
\|\omega\|_2^{2\ell_{\mathrm{mom}}}\,d\omega
\right)^{1/2}
\|\psi\|_2
<\infty.
\]
At \(n_{\mathrm{cut}}=0\), the cutoff is zero almost everywhere because \(p>0\).
Plancherel and spatial negation give
\[
\|G_{\psi,n_{\mathrm{cut}}}\|_2^2
=
\|\psi_{n_{\mathrm{cut}}}\|_2^2,
\qquad
\|G_{\psi,n_{\mathrm{cut}}}-G_{\psi,n'_{\mathrm{cut}}}\|_2^2
=
\|\psi_{n_{\mathrm{cut}}}-\psi_{n'_{\mathrm{cut}}}\|_2^2.
\]
Dominated convergence makes these inverse truncations Cauchy in \(L^2\).

To identify their limit explicitly, let \(\mathcal F_{\mathrm{cyc},p}\) denote the unitary \(L^2\) extension of
\[
(\mathcal F_{\mathrm{cyc},p}\psi)(\eta)
=
\int_{\mathbb R^p}\psi(\omega)
e^{-2\pi\mathrm i\eta^{\mathsf T}\omega}\,d\omega,
\qquad \psi\in L^1\cap L^2,\quad \eta\in\mathbb R^p,
\]
and define, as an \(L^2\) class,
\[
(\mathcal I_p\psi)(u)
=
(2\pi)^{-p/2}
(\mathcal F_{\mathrm{cyc},p}\psi)\!\left(-\frac{u}{2\pi}\right).
\]
The squared prefactor \((2\pi)^{-p}\) cancels the dilation Jacobian \((2\pi)^p\). Thus
\[
\|\mathcal I_p\psi\|_2^2=\|\psi\|_2^2,
\qquad
\|\mathcal I_p\psi-\mathcal I_p\chi\|_2^2=\|\psi-\chi\|_2^2,
\qquad \psi,\chi\in L^2(\mathbb R^p;\mathbb C).
\]
The defining integral identifies
\(G_{\psi,n_{\mathrm{cut}}}=\mathcal I_p\psi_{n_{\mathrm{cut}}}\) almost everywhere, whence
\[
\|G_{\psi,n_{\mathrm{cut}}}-\mathcal I_p\psi\|_2^2
=
\|\psi_{n_{\mathrm{cut}}}-\psi\|_2^2
\longrightarrow0.
\]

For an admissible extension \(G\in L^1\cap L^2\), its original spatial class is recovered by this construction. Fubini, applied to the integrable product of the truncated profile and an \(L^1\cap L^2\) test function \(\varphi\), gives
\[
\int_{\mathbb R^p}\overline{\varphi(u)}
G_{\psi,n_{\mathrm{cut}}}(u)\,du
=
\int_{\mathbb R^p}
\overline{(\mathcal F_p\varphi)(\omega)}
\psi_{n_{\mathrm{cut}}}(\omega)\,d\omega.
\]
Taking \(\psi=\mathcal F_pG\) and \(\varphi=G\), and expanding the squared distance using this pairing and the two Plancherel identities, yields
\[
\|G-G_{\mathcal F_pG,n_{\mathrm{cut}}}\|_2^2
=
\|\mathcal F_pG-(\mathcal F_pG)_{n_{\mathrm{cut}}}\|_2^2
\longrightarrow0.
\]
Uniqueness of the \(L^2\) limit therefore gives
\[
\mathcal I_p(\mathcal F_pG)=G
\quad\text{almost everywhere}.
\]
% lean: CausalSmith.Experimentation.BivariateSobolevDesignCapacity.reflectionInverseL2Representative_fourier_ae_original

\item \textit{The one-dimensional endpoint.}
Differentiating the inverse truncation under the integral is justified by the frequency moments established above. For \(1\leq j\leq p\),
\[
\partial_jG_{\psi,n_{\mathrm{cut}}}(u)
=
(2\pi)^{-p/2}
\int_{\mathbb R^p}
\mathrm i\omega_j e^{\mathrm i\omega^{\mathsf T}u}
\psi_{n_{\mathrm{cut}}}(\omega)\,d\omega.
\]
These derivatives are continuous and square integrable. Applying the inverse-integral energy identity to the profile
\(\mathrm i\omega_j\psi_{n_{\mathrm{cut}}}\) gives
\[
\int_{\mathbb R^p}
|\partial_jG_{\psi,n_{\mathrm{cut}}}(u)|^2\,du
=
\int_{\mathbb R^p}
\omega_j^2|\psi_{n_{\mathrm{cut}}}(\omega)|^2\,d\omega.
\]
Summing the coordinate identities proves the exact first-order identity
\[
\int_{\mathbb R^p}
\left(
|G_{\psi,n_{\mathrm{cut}}}(u)|^2+
\frac1p\sum_{j=1}^p
|\partial_jG_{\psi,n_{\mathrm{cut}}}(u)|^2
\right)\,du
=
\int_{\mathbb R^p}
\left(1+\frac{\|\omega\|_2^2}{p}\right)
|\psi_{n_{\mathrm{cut}}}(\omega)|^2\,d\omega.
\]

Now take \(p=1\), identifying a one-dimensional integer frequency with \(a\in\mathbb Z\). Define the signed reflected derivative on the period-two circle by
\[
V_{\psi,n_{\mathrm{cut}},1}(y)
=
\begin{cases}
G'_{\psi,n_{\mathrm{cut}}}(y),&0<y<1,\\
-G'_{\psi,n_{\mathrm{cut}}}(2-y),&1<y<2,\\
0,&y\in\{0,1\},
\end{cases}
\]
using representatives in \([0,2)\). For a continuously differentiable periodic test function \(\varphi\), integration by parts on the two intervals has total boundary term
\[
\bigl[G_{\psi,n_{\mathrm{cut}}}(y)\varphi(y)\bigr]_0^1
+
\bigl[G_{\psi,n_{\mathrm{cut}}}(2-y)\varphi(y)\bigr]_1^2
=
G_{\psi,n_{\mathrm{cut}}}(0)\bigl(\varphi(2)-\varphi(0)\bigr)
=
0.
\]
It follows that
\[
\int_{\mathbb T_2^1}
(\mathscr D_1G_{\psi,n_{\mathrm{cut}}})(y)\varphi'(y)\,d\mu_1(y)
=
-\int_{\mathbb T_2^1}
V_{\psi,n_{\mathrm{cut}},1}(y)\varphi(y)\,d\mu_1(y).
\]
Taking \(\varphi(y)=e^{-\mathrm i\pi ay}\) gives the exact multiplier identity
\[
\widehat V_{\psi,n_{\mathrm{cut}},1}(a)
=
\mathrm i\pi a\,
\widehat{\mathscr D_1G_{\psi,n_{\mathrm{cut}}}}(a).
\]
Parseval for the reflected function and its signed derivative, followed by folding, therefore gives
\[
\begin{aligned}
\sum_{a\in\mathbb Z}(1+\pi^2a^2)
\left|\widehat{\mathscr D_1G_{\psi,n_{\mathrm{cut}}}}(a)\right|^2
&=
\int_0^1
\left(
|G_{\psi,n_{\mathrm{cut}}}(u)|^2+
|G'_{\psi,n_{\mathrm{cut}}}(u)|^2
\right)\,du\\
&\leq
\int_{\mathbb R}
\left(
|G_{\psi,n_{\mathrm{cut}}}(u)|^2+
|G'_{\psi,n_{\mathrm{cut}}}(u)|^2
\right)\,du\\
&=
\int_{\mathbb R}(1+\omega^2)
|\psi_{n_{\mathrm{cut}}}(\omega)|^2\,d\omega\\
&\leq
\int_{\mathbb R}(1+\omega^2)|\psi(\omega)|^2\,d\omega.
\end{aligned}
\]

For an admissible extension \(G\), take \(\psi=\mathcal F_1G\).
The \(L^2\) contraction of \(\mathscr D_1\) and the spatial convergence above imply
\[
\mathscr D_1G_{\mathcal F_1G,n_{\mathrm{cut}}}
\longrightarrow\mathscr D_1G
\quad\text{in }L^2(\mu_1).
\]
Every Fourier coefficient converges, since its absolute difference is bounded by the \(L^2\) distance. For every finite frequency set \(\mathcal A_{\mathrm{fin}}\subset\mathbb Z\), passing to the limit in its finite weighted sum gives
\[
\sum_{a\in\mathcal A_{\mathrm{fin}}}(1+\pi^2a^2)
|\widehat{\mathscr D_1G}(a)|^2
\leq
\int_{\mathbb R}(1+\omega^2)|\mathcal F_1G(\omega)|^2\,d\omega.
\]
Taking the supremum over finite sets proves the same inequality for the full sum. Folding transports \(G=g\) almost everywhere on the cube to
\(\mathscr D_1G=h\) almost everywhere on the torus. Infimizing over admissible extensions gives
\[
\sum_{a\in\mathbb Z}(1+\pi^2a^2)|\widehat h(a)|^2
\leq N_{1,1}(g)^2.
\]
Extensions with infinite first-order energy satisfy the intermediate inequality in the extended sense as well.
% lean: CausalSmith.Experimentation.BivariateSobolevDesignCapacity.componentFourierBudget_one_firstOrder_le

\item \textit{The two-dimensional endpoint.}
For \(p=2\), define, for \(j\in\{1,2\}\),
\[
\sigma_j(y)=
\begin{cases}
1,&0<y_j<1,\\
-1,&1<y_j<2,\\
0,&y_j\in\{0,1\},
\end{cases}
\qquad
V_{\psi,n_{\mathrm{cut}},j}(y)
=
\sigma_j(y)
(\partial_jG_{\psi,n_{\mathrm{cut}}})(b_2(y)).
\]
Apply the one-dimensional integration-by-parts computation on each coordinate slice and integrate in the other coordinate. The boundary terms cancel on each slice, giving
\[
\int_{\mathbb T_2^2}
(\mathscr D_2G_{\psi,n_{\mathrm{cut}}})(y)
\partial_j\varphi(y)\,d\mu_2(y)
=
-\int_{\mathbb T_2^2}
V_{\psi,n_{\mathrm{cut}},j}(y)\varphi(y)\,d\mu_2(y).
\]
For \(a=(a_1,a_2)\in\mathbb Z^2\), the character test yields
\[
\widehat V_{\psi,n_{\mathrm{cut}},j}(a)
=
\mathrm i\pi a_j
\widehat{\mathscr D_2G_{\psi,n_{\mathrm{cut}}}}(a).
\]
Joint Parseval, obtained from the product Fourier basis, combines the zeroth-order energy and the two derivative energies with factor \(1/2\):
\[
\begin{aligned}
&\sum_{a\in\mathbb Z^2}
\left(1+\frac{\pi^2(a_1^2+a_2^2)}2\right)
\left|\widehat{\mathscr D_2G_{\psi,n_{\mathrm{cut}}}}(a)\right|^2\\
&\quad=
\int_{[0,1]^2}
\left(
|G_{\psi,n_{\mathrm{cut}}}(u)|^2+
\frac{
|\partial_1G_{\psi,n_{\mathrm{cut}}}(u)|^2+
|\partial_2G_{\psi,n_{\mathrm{cut}}}(u)|^2}{2}
\right)\,du\\
&\quad\leq
\int_{\mathbb R^2}
\left(
|G_{\psi,n_{\mathrm{cut}}}(u)|^2+
\frac{
|\partial_1G_{\psi,n_{\mathrm{cut}}}(u)|^2+
|\partial_2G_{\psi,n_{\mathrm{cut}}}(u)|^2}{2}
\right)\,du\\
&\quad=
\int_{\mathbb R^2}
\left(1+\frac{\|\omega\|_2^2}{2}\right)
|\psi_{n_{\mathrm{cut}}}(\omega)|^2\,d\omega\\
&\quad\leq
\int_{\mathbb R^2}
\left(1+\frac{\|\omega\|_2^2}{2}\right)|\psi(\omega)|^2\,d\omega.
\end{aligned}
\]
For \(\psi=\mathcal F_2G\), spatial \(L^2\) convergence and the reflection contraction again give convergence of every reflected coefficient. Passing first to each finite frequency sum and then to their supremum proves
\[
\sum_{a\in\mathbb Z^2}
\left(1+\frac{\pi^2\|a\|_2^2}{2}\right)
|\widehat{\mathscr D_2G}(a)|^2
\leq
\int_{\mathbb R^2}
\left(1+\frac{\|\omega\|_2^2}{2}\right)
|\mathcal F_2G(\omega)|^2\,d\omega.
\]
The equality \(\mathscr D_2G=h\) almost everywhere for every admissible extension, followed by the extension infimum, yields
\[
\sum_{a\in\mathbb Z^2}
\left(1+\frac{\pi^2\|a\|_2^2}{2}\right)|\widehat h(a)|^2
\leq N_{2,1}(g)^2.
\]
% lean: CausalSmith.Experimentation.BivariateSobolevDesignCapacity.componentFourierBudget_two_firstOrder_le

\item \textit{The common operator for fractional smoothness.}
Assume now that \(0<s<1\). On the complete frequency \(L^2\) space, define
\[
\mathscr T_p\psi=\mathscr D_p(\mathcal I_p\psi),
\qquad p\in\{1,2\}.
\]
This is a common linear operator for the two endpoints. For \(\vartheta\in[0,1]\), define the extended gauges
\[
\|\psi\|_{\mathrm{freq},p,\vartheta}^2
=
\int_{\mathbb R^p}
\left(1+\frac{\|\omega\|_2^2}{p}\right)^\vartheta
|\psi(\omega)|^2\,d\omega,
\]
\[
\|\zeta\|_{\mathrm{per},p,\vartheta}^2
=
\sum_{a\in\mathbb Z^p}
\left(1+\frac{\pi^2\|a\|_2^2}{p}\right)^\vartheta
|\widehat\zeta(a)|^2,
\]
for \(\psi\in L^2(\mathbb R^p;\mathbb C)\) and
\(\zeta\in L^2(\mathbb T_2^p;\mathbb C)\); a divergent expression has value \(+\infty\).

The inverse isometry and folding prove
\[
\|\mathscr T_p\psi\|_{\mathrm{per},p,0}
\leq
\|\psi\|_{\mathrm{freq},p,0}.
\]
For every frequency \(L^2\) class, the inverse truncations converge to
\(\mathcal I_p\psi\), so
\[
\mathscr D_pG_{\psi,n_{\mathrm{cut}}}
\longrightarrow\mathscr T_p\psi
\quad\text{in }L^2(\mu_p).
\]
Apply the smooth spectral estimate in the appropriate dimension and pass to finite sums as above. This proves the second endpoint on the same operator:
\[
\|\mathscr T_p\psi\|_{\mathrm{per},p,1}
\leq
\|\psi\|_{\mathrm{freq},p,1}.
\]
The inequality includes profiles with infinite first-order gauge.
% lean: CausalSmith.Experimentation.BivariateSobolevDesignCapacity.reflectionProfileOperator_firstOrder_norm_le

\item \textit{The normalized interpolation energy.}
For two extended gauges \(\mathsf n_0,\mathsf n_1\) on a common additive space, define, for \(\rho>0\),
\[
\mathfrak K^2(\rho;\mathsf n_0,\mathsf n_1;\zeta)
=
\inf_{\zeta=\zeta_0+\zeta_1}
\left\{
\mathsf n_0(\zeta_0)^2+
\rho^2\mathsf n_1(\zeta_1)^2
\right\},
\]
and, for \(0<s<1\),
\[
\mathfrak J_s(\mathsf n_0,\mathsf n_1;\zeta)
=
\frac{2\sin(\pi s)}{\pi}
\int_0^\infty
\rho^{-1-2s}
\mathfrak K^2(\rho;\mathsf n_0,\mathsf n_1;\zeta)\,d\rho.
\]
Every source decomposition maps to a target decomposition under \(\mathscr T_p\). The two endpoint contractions therefore imply
\[
\begin{aligned}
&\mathfrak K^2\!\left(
\rho;
\|\cdot\|_{\mathrm{per},p,0},
\|\cdot\|_{\mathrm{per},p,1};
\mathscr T_p\psi
\right)\\
&\qquad\leq
\mathfrak K^2\!\left(
\rho;
\|\cdot\|_{\mathrm{freq},p,0},
\|\cdot\|_{\mathrm{freq},p,1};
\psi
\right).
\end{aligned}
\]
Integration against the positive normalized measure gives the same contraction for \(\mathfrak J_s\).

We evaluate the source functional exactly. Set
\[
\mathsf w_{\mathrm{freq},p}(\omega)
=
1+\frac{\|\omega\|_2^2}{p},
\qquad
\psi_{0,\rho}(\omega)
=
\frac{\rho^2\mathsf w_{\mathrm{freq},p}(\omega)}
{1+\rho^2\mathsf w_{\mathrm{freq},p}(\omega)}\psi(\omega),
\]
\[
\psi_{1,\rho}(\omega)
=
\frac{1}{1+\rho^2\mathsf w_{\mathrm{freq},p}(\omega)}\psi(\omega).
\]
These are an admissible \(L^2\) decomposition of \(\psi\). Their cost satisfies
\[
|\psi_{0,\rho}(\omega)|^2+
\rho^2\mathsf w_{\mathrm{freq},p}(\omega)
|\psi_{1,\rho}(\omega)|^2
=
\frac{\rho^2\mathsf w_{\mathrm{freq},p}(\omega)}
{1+\rho^2\mathsf w_{\mathrm{freq},p}(\omega)}
|\psi(\omega)|^2
\leq|\psi(\omega)|^2.
\]
For every other decomposition
\(\psi=\chi_0+\chi_1\), completing the square gives
\[
\begin{aligned}
|\chi_0|^2+\rho^2\mathsf w_{\mathrm{freq},p}|\chi_1|^2
&=
\frac{\rho^2\mathsf w_{\mathrm{freq},p}}
{1+\rho^2\mathsf w_{\mathrm{freq},p}}|\psi|^2\\
&\quad+
(1+\rho^2\mathsf w_{\mathrm{freq},p})
\left|
\chi_1-\frac{\psi}{1+\rho^2\mathsf w_{\mathrm{freq},p}}
\right|^2.
\end{aligned}
\]
Thus the displayed decomposition attains the infimum, including when the competing decomposition has infinite cost, and
\[
\mathfrak K^2\!\left(
\rho;
\|\cdot\|_{\mathrm{freq},p,0},
\|\cdot\|_{\mathrm{freq},p,1};
\psi
\right)
=
\int_{\mathbb R^p}
\frac{\rho^2\mathsf w_{\mathrm{freq},p}(\omega)}
{1+\rho^2\mathsf w_{\mathrm{freq},p}(\omega)}
|\psi(\omega)|^2\,d\omega.
\]
For every \(\beta>0\), substitution
\(\tau=\rho\sqrt{\beta}\) gives
\[
\frac{2\sin(\pi s)}{\pi}
\int_0^\infty
\rho^{-1-2s}\frac{\rho^2\beta}{1+\rho^2\beta}\,d\rho
=
\beta^s\frac{2\sin(\pi s)}{\pi}
\int_0^\infty\frac{\tau^{1-2s}}{1+\tau^2}\,d\tau
=
\beta^s,
\]
because the last integral is \(\pi/(2\sin(\pi s))\).
Tonelli consequently proves
\[
\mathfrak J_s\!\left(
\|\cdot\|_{\mathrm{freq},p,0},
\|\cdot\|_{\mathrm{freq},p,1};
\psi
\right)
=
\|\psi\|_{\mathrm{freq},p,s}^2.
\]
Combining this identity with the operator contraction yields
\[
\mathfrak J_s\!\left(
\|\cdot\|_{\mathrm{per},p,0},
\|\cdot\|_{\mathrm{per},p,1};
\mathscr T_p\psi
\right)
\leq
\|\psi\|_{\mathrm{freq},p,s}^2.
\]
% lean: CausalSmith.Experimentation.BivariateSobolevDesignCapacity.reflectionProfileOperator_kNormSq_le_energy

\item \textit{Transport through Fourier coefficients.}
Define the coefficient map and its sequence gauges by
\[
\mathscr C_p\zeta=(\widehat\zeta(a))_{a\in\mathbb Z^p},
\qquad
\mathsf w_{\mathrm{per},p}(a)
=
1+\frac{\pi^2\|a\|_2^2}{p},
\]
\[
\|\gamma\|_{\mathrm{seq},p,0}^2
=
\sum_{a\in\mathbb Z^p}|\gamma(a)|^2,
\qquad
\|\gamma\|_{\mathrm{seq},p,1}^2
=
\sum_{a\in\mathbb Z^p}
\mathsf w_{\mathrm{per},p}(a)|\gamma(a)|^2,
\qquad \gamma\in\ell^2(\mathbb Z^p;\mathbb C).
\]
Parseval and the first-order gauge definition give
\[
\|\mathscr C_p\zeta\|_{\mathrm{seq},p,0}
=
\|\zeta\|_{\mathrm{per},p,0},
\qquad
\|\mathscr C_p\zeta\|_{\mathrm{seq},p,1}
=
\|\zeta\|_{\mathrm{per},p,1}.
\]
Mapping each torus decomposition through \(\mathscr C_p\) therefore gives
\[
\mathfrak J_s\!\left(
\|\cdot\|_{\mathrm{seq},p,0},
\|\cdot\|_{\mathrm{seq},p,1};
\mathscr C_p\zeta
\right)
\leq
\mathfrak J_s\!\left(
\|\cdot\|_{\mathrm{per},p,0},
\|\cdot\|_{\mathrm{per},p,1};
\zeta
\right).
\]
The same pointwise minimizing decomposition and scalar integral from the previous step, now with counting measure and weight
\(\mathsf w_{\mathrm{per},p}\), show that
\[
\mathfrak J_s\!\left(
\|\cdot\|_{\mathrm{seq},p,0},
\|\cdot\|_{\mathrm{seq},p,1};
\mathscr C_p\zeta
\right)
=
\sum_{a\in\mathbb Z^p}
\mathsf w_{\mathrm{per},p}(a)^s|\widehat\zeta(a)|^2.
\]
All these weights are finite and at least one, so the minimizing decompositions belong to the ambient \(\ell^2\) space. Applying the two inequalities to \(\zeta=\mathscr T_p\psi\) yields
\[
\sum_{a\in\mathbb Z^p}
\left(1+\frac{\pi^2\|a\|_2^2}{p}\right)^s
|\widehat{\mathscr T_p\psi}(a)|^2
\leq
\int_{\mathbb R^p}
\left(1+\frac{\|\omega\|_2^2}{p}\right)^s
|\psi(\omega)|^2\,d\omega.
\]

For an admissible extension \(G\), the inverse identification in the first step gives
\[
\mathscr T_p(\mathcal F_pG)=\mathscr D_pG=h
\quad\text{almost everywhere}.
\]
Hence
\[
\sum_{a\in\mathbb Z^p}
\left(1+\frac{\pi^2\|a\|_2^2}{p}\right)^s
|\widehat h(a)|^2
\leq
\int_{\mathbb R^p}
\left(1+\frac{\|\omega\|_2^2}{p}\right)^s
|\mathcal F_pG(\omega)|^2\,d\omega.
\]
Infimizing over all admissible extensions proves the component assertion for \(0<s<1\). Together with the two endpoint cases, this establishes it throughout \(0<s\leq1\).
% lean: CausalSmith.Experimentation.BivariateSobolevDesignCapacity.componentFourierBudget_le_sobolevNormSq

\item \textit{Coefficient additivity and the order-two support bound.}
Fix \(d\geq2\) and \(m\in\mathcal H_{d,s}\). By \cref{obj:def:sobolev-class}, \(m\) is centered and square integrable and satisfies \cref{obj:ass:order-two,obj:ass:pooled-budget}. Use the canonical components defined in \cref{obj:synth_1}. Define their ambient reflections by
\[
h_j^{\mathrm{amb}}(y)
=
g_j(m)((b_d(y))_j),
\qquad
h_{j\ell}^{\mathrm{amb}}(y)
=
g_{j\ell}(m)((b_d(y))_j,(b_d(y))_\ell),
\qquad j<\ell,
\]
and their coefficients by
\[
\widehat h_j^{\mathrm{amb}}(k)
=
2^{-d}\int_{[0,2)^d}
h_j^{\mathrm{amb}}(y)e^{-\mathrm i\pi k^{\mathsf T}y}\,dy,
\]
\[
\widehat h_{j\ell}^{\mathrm{amb}}(k)
=
2^{-d}\int_{[0,2)^d}
h_{j\ell}^{\mathrm{amb}}(y)e^{-\mathrm i\pi k^{\mathsf T}y}\,dy,
\qquad k\in\mathbb Z^d.
\]

The requisite component \(L^2\) regularity follows from conditional Cauchy--Schwarz. In the integration notation of \cref{obj:synth_1},
\[
\int_0^1|g_j(m)(x_j)|^2\,dx_j
\leq
\int_{[0,1]^d}|m(x)|^2\,dx,
\]
and
\[
\int_{[0,1]^2}
\left|
\int_{[0,1]^{d-2}}
m(x_j,x_\ell,u_{-\{j,\ell\}})\,du_{-\{j,\ell\}}
\right|^2\,dx_j\,dx_\ell
\leq
\int_{[0,1]^d}|m(x)|^2\,dx.
\]
Subtracting the two lifted main effects therefore gives an \(L^2\) pair effect. For \(d=2\), the inner zero-dimensional integral is evaluation, and the second inequality is equality. Folding preserves this regularity. The almost-everywhere decomposition supplied by \cref{obj:ass:order-two} thus permits integration term by term:
\[
\widehat F(k)
=
\sum_{j=1}^d\widehat h_j^{\mathrm{amb}}(k)
+
\sum_{j<\ell}\widehat h_{j\ell}^{\mathrm{amb}}(k).
\]

To establish coordinate locality, fix a coordinate \(\jmath\) with
\(k_\jmath\neq0\), and define a torus translation by
\[
(\tau_{k,\jmath})_\iota
=
\begin{cases}
-1/k_\jmath\pmod{2},&\iota=\jmath,\\
0\pmod{2},&\iota\neq\jmath,
\end{cases}
\qquad 1\leq\iota\leq d.
\]
If a reflected component is independent of coordinate \(\jmath\), it is invariant under this translation, whereas its coefficient character acquires the factor
\[
e^{-\mathrm i\pi k^{\mathsf T}\tau_{k,\jmath}}=-1.
\]
Translation invariance of Haar measure then makes the component coefficient equal to its own negative, hence zero. In particular,
\[
\widehat h_j^{\mathrm{amb}}(k)=0
\quad\text{if some }k_\jmath\neq0\text{ with }\jmath\neq j,
\]
\[
\widehat h_{j\ell}^{\mathrm{amb}}(k)=0
\quad\text{if some }k_\jmath\neq0\text{ with }\jmath\notin\{j,\ell\}.
\]
If \(|\operatorname{supp}(k)|>2\), choose three distinct active coordinates
\(\jmath_1,\jmath_2,\jmath_3\). Every main component excludes at least one of them, and every pair component also excludes at least one. All terms in the coefficient decomposition vanish, proving
\[
\widehat F(k)=0
\qquad\text{when }|\operatorname{supp}(k)|>2.
\]
% lean: CausalSmith.Experimentation.BivariateSobolevDesignCapacity.Fhat_zero_of_orderTwo_from_components

\item \textit{Centering and unique component coefficients.}
Fubini applied to the canonical marginal formulas gives
\[
\int_0^1g_j(m)(\upsilon)\,d\upsilon
=
\int_{[0,1]^d}m(x)\,dx
=
0.
\]
Integrating the raw pair marginal in its second coordinate recovers the first main marginal. Consequently, for almost every \(x_j\),
\[
\int_0^1g_{j\ell}(m)(x_j,\upsilon)\,d\upsilon
=
g_j(m)(x_j)-g_j(m)(x_j)
-\int_0^1g_\ell(m)(\upsilon)\,d\upsilon
=
0.
\]
Interchanging the two marginal coordinates gives, for almost every \(x_\ell\),
\[
\int_0^1g_{j\ell}(m)(\upsilon,x_\ell)\,d\upsilon
=
g_\ell(m)(x_\ell)
-\int_0^1g_j(m)(\upsilon)\,d\upsilon
-g_\ell(m)(x_\ell)
=
0.
\]
These identities also apply at \(d=2\), with the evaluation convention for the raw pair marginal.

Folding preserves the coordinate integrals. Thus the main reflected coefficient at zero vanishes. For a pair effect, if \(k_\ell=0\), integrating first in coordinate \(\ell\) gives zero by the first pair-centering identity; if \(k_j=0\), use the other identity. Combining these facts with coordinate locality yields
\[
\widehat h_j^{\mathrm{amb}}(k)\neq0
\ \Longrightarrow\
\operatorname{supp}(k)=\{j\},
\]
\[
\widehat h_{j\ell}^{\mathrm{amb}}(k)\neq0
\ \Longrightarrow\
\operatorname{supp}(k)=\{j,\ell\}.
\]
The supports of distinct canonical components are therefore disjoint. Substituting these support statements into coefficient additivity proves
\[
\widehat F(k)
=
\begin{cases}
\widehat h_j^{\mathrm{amb}}(k),
&\operatorname{supp}(k)=\{j\},\\
\widehat h_{j\ell}^{\mathrm{amb}}(k),
&\operatorname{supp}(k)=\{j,\ell\},\quad j<\ell.
\end{cases}
\]
For the pair case, ordering the two supporting coordinates increasingly selects the unique pair in the decomposition.
% lean: CausalSmith.Experimentation.BivariateSobolevDesignCapacity.Fhat_orderTwo_pair_component

\item \textit{Pooling and local coefficient transport.}
For the auxiliary sums over all ambient frequencies, extend the stated complexity by
\[
t(0)=0,
\qquad
t(k)=\frac{\|k\|_2^2}{|\operatorname{supp}(k)|}
\quad(k\neq0).
\]
At a singleton-support frequency, the unique main coefficient just identified supplies the full coefficient; at a pair-support frequency, the unique ordered pair supplies it. At every other frequency the indicator below is zero. Hence, pointwise in \(k\),
\[
\begin{aligned}
&\mathbf1_{\{1\leq|\operatorname{supp}(k)|\leq2\}}
(1+\pi^2t(k))^s|\widehat F(k)|^2\\
&\quad\leq
\sum_{j=1}^d
(1+\pi^2t(k))^s|\widehat h_j^{\mathrm{amb}}(k)|^2
+
\sum_{j<\ell}
(1+\pi^2t(k))^s|\widehat h_{j\ell}^{\mathrm{amb}}(k)|^2.
\end{aligned}
\]
Summing and interchanging the nonnegative frequency sums with the finite component sums gives
\[
\begin{aligned}
&\sum_{\substack{k\in\mathbb Z^d\\1\leq|\operatorname{supp}(k)|\leq2}}
(1+\pi^2t(k))^s|\widehat F(k)|^2\\
&\quad\leq
\sum_{j=1}^d\sum_{k\in\mathbb Z^d}
(1+\pi^2t(k))^s|\widehat h_j^{\mathrm{amb}}(k)|^2\\
&\qquad+
\sum_{j<\ell}\sum_{k\in\mathbb Z^d}
(1+\pi^2t(k))^s|\widehat h_{j\ell}^{\mathrm{amb}}(k)|^2.
\end{aligned}
\]

Define the local reflections and coefficients by
\[
h_j^{\mathrm{loc}}(y)
=
g_j(m)(b_1(y)),
\qquad
\widehat h_j^{\mathrm{loc}}(\alpha)
=
\frac12\int_{[0,2)}
h_j^{\mathrm{loc}}(y)e^{-\mathrm i\pi\alpha y}\,dy,
\qquad \alpha\in\mathbb Z,
\]
\[
h_{j\ell}^{\mathrm{loc}}(y)
=
g_{j\ell}(m)((b_2(y))_1,(b_2(y))_2),
\qquad
\widehat h_{j\ell}^{\mathrm{loc}}(a)
=
\frac14\int_{[0,2)^2}
h_{j\ell}^{\mathrm{loc}}(y)e^{-\mathrm i\pi a^{\mathsf T}y}\,dy,
\qquad a\in\mathbb Z^2.
\]
For \(\alpha\in\mathbb Z\) and \(a\in\mathbb Z^2\), define the ambient frequency embeddings
\[
(\iota_j(\alpha))_\jmath
=
\begin{cases}
\alpha,&\jmath=j,\\
0,&\jmath\neq j,
\end{cases}
\qquad
(\iota_{j\ell}(a))_\jmath
=
\begin{cases}
a_1,&\jmath=j,\\
a_2,&\jmath=\ell,\\
0,&\jmath\notin\{j,\ell\}.
\end{cases}
\]
Coordinate projection preserves product Haar probability, and folding commutes with selecting coordinates. Integrating the unused coordinates therefore gives
\[
\widehat h_j^{\mathrm{amb}}(\iota_j(\alpha))
=
\widehat h_j^{\mathrm{loc}}(\alpha),
\qquad
\widehat h_{j\ell}^{\mathrm{amb}}(\iota_{j\ell}(a))
=
\widehat h_{j\ell}^{\mathrm{loc}}(a).
\]
At the frequencies where the corresponding ambient coefficients can be nonzero,
\[
t(\iota_j(\alpha))=\alpha^2
\quad(\alpha\neq0),
\qquad
t(\iota_{j\ell}(a))=\frac{a_1^2+a_2^2}{2}
\quad(a_1a_2\neq0).
\]
Indeed, the ambient squared frequency length is the local squared length, and the support cardinality is respectively one or two.

The ambient support statements restrict each sum to these embedded frequencies. Restriction to the supporting coordinates is injective, so enlarging to all local frequencies gives
\[
\sum_{k\in\mathbb Z^d}
(1+\pi^2t(k))^s|\widehat h_j^{\mathrm{amb}}(k)|^2
\leq
\sum_{\alpha\in\mathbb Z}
(1+\pi^2\alpha^2)^s|\widehat h_j^{\mathrm{loc}}(\alpha)|^2,
\]
\[
\sum_{k\in\mathbb Z^d}
(1+\pi^2t(k))^s|\widehat h_{j\ell}^{\mathrm{amb}}(k)|^2
\leq
\sum_{a\in\mathbb Z^2}
\left(1+\frac{\pi^2\|a\|_2^2}{2}\right)^s
|\widehat h_{j\ell}^{\mathrm{loc}}(a)|^2.
\]
Thus the ambient bound is bounded in turn by the sum of the local component Fourier budgets.
% lean: CausalSmith.Experimentation.BivariateSobolevDesignCapacity.reflectedBudget_le_local_component_budgets

\item \textit{The pooled budget and cube energy.}
The local canonical components are Borel and square integrable by the marginal argument above. Their restriction norms are finite because the finite, nonnegative sum in \cref{obj:ass:pooled-budget} is bounded by one. Applying the component contraction to each main effect and each pair effect yields
\[
\begin{aligned}
&\sum_{\substack{k\in\mathbb Z^d\\1\leq|\operatorname{supp}(k)|\leq2}}
(1+\pi^2t(k))^s|\widehat F(k)|^2\\
&\quad\leq
\sum_{j=1}^d\sum_{\alpha\in\mathbb Z}
(1+\pi^2\alpha^2)^s|\widehat h_j^{\mathrm{loc}}(\alpha)|^2\\
&\qquad+
\sum_{j<\ell}\sum_{a\in\mathbb Z^2}
\left(1+\frac{\pi^2\|a\|_2^2}{2}\right)^s
|\widehat h_{j\ell}^{\mathrm{loc}}(a)|^2\\
&\quad\leq
\sum_{j=1}^dN_{1,s}(g_j(m))^2
+
\sum_{j<\ell}N_{2,s}(g_{j\ell}(m))^2
\leq1,
\end{aligned}
\]
where the final inequality is precisely \cref{obj:ass:pooled-budget}.
% lean: CausalSmith.Experimentation.BivariateSobolevDesignCapacity.reflectedBudget_le_one_of_sobolevClass

Folding sends Haar probability to cube Lebesgue probability in dimension \(d\), by the same coordinatewise reflection calculation used above. Centering therefore gives
\[
\widehat F(0)
=
\int_{\mathbb T_2^d}F(y)\,d\mu_d(y)
=
\int_{[0,1]^d}m(x)\,dx
=
0,
\qquad d\mu_d(y)=2^{-d}\,dy.
\]
Since \(d\geq2\), Parseval applies in positive dimension. Also, \(s>0\) and \(t(k)\geq0\) give
\[
(1+\pi^2t(k))^s\geq1.
\]
Using the zero coefficient and the already established vanishing for support size greater than two, we conclude
\[
\begin{aligned}
\int_{[0,1]^d}|m(x)|^2\,dx
&=
\int_{\mathbb T_2^d}|F(y)|^2\,d\mu_d(y)\\
&=
\sum_{k\in\mathbb Z^d}|\widehat F(k)|^2\\
&=
\sum_{\substack{k\in\mathbb Z^d\\1\leq|\operatorname{supp}(k)|\leq2}}
|\widehat F(k)|^2\\
&\leq
\sum_{\substack{k\in\mathbb Z^d\\1\leq|\operatorname{supp}(k)|\leq2}}
(1+\pi^2t(k))^s|\widehat F(k)|^2
\leq1.
\end{aligned}
\]
This establishes all the asserted conclusions.
% lean: CausalSmith.Experimentation.BivariateSobolevDesignCapacity.cube_second_moment_le_of_reflectedBudget
\end{enumerate}
\end{proof}

The component inequality in \cref{obj:lem:exact-reflection-budget} preserves the squared restriction budget with constant one throughout \(0<s\leq1\). Canonical centering assigns main effects to frequencies supported on their single coordinate and pair effects to frequencies supported on both of their coordinates. Thus the full spectral budget retains the component allocation specified by \cref{obj:ass:pooled-budget}. Its exact normalization supplies the budget used in the numerical attaining bound.

The accompanying support and centering conclusions identify the spectral rows relevant to the population criterion. The reflected function has frequencies supported on one or two coordinates and has zero constant coefficient. Its cube \(L^2\) bound also controls the independent-sign branch of \cref{obj:def:capacity-handle}.

\subsection{Common-shift spectral transfer}

The randomized lift in \cref{obj:synth_2} makes each lifted covariate Haar-uniform. A common translation preserves the product Haar law and makes the transformed sample independent of the shift. This independence is the basis of the spectral identity: the shift supplies a common random phase for each frequency while assignment depends on the transformed sample. The next statement joins the reflection bound to this identity for the allowed measurable prognostic functions.

\begin{lemmav}{lem:reflection-transfer}[Reflection and spectral transfer]
Write \(\mathbb T_2^p=(\mathbb R/2\mathbb Z)^p\) for the period-two torus, equipped with normalized Haar measure, and let \(b_p\) fold each coordinate onto \([0,1]\), by taking its absolute value in a representative from \([-1,1]\). For a cube function \(g\), its even reflection \(h\) and normalized Fourier coefficients are
\[
h=g\circ b_p,\qquad
\widehat h(a)=\int_{\mathbb T_2^p}h(w)
  \exp(-\mathrm i\pi a\cdot w)\,dw,
\qquad a\in\mathbb Z^p.
\]
The squared restriction norm is
\[
N_{p,s}(g)^2
=
\inf_{\substack{G\in L^1(\mathbb R^p;\mathbb C)\cap L^2(\mathbb R^p;\mathbb C)\\
G=g\ {\rm a.e.}\ {\rm on}\ [0,1]^p}}
\int_{\mathbb R^p}
\left|
(2\pi)^{-p/2}
\int_{\mathbb R^p}\exp(-\mathrm i\omega\cdot u)G(u)\,du
\right|^2
\left(1+\frac{\|\omega\|^2}{p}\right)^s\,d\omega.
\]
For every \(p\in\{1,2\}\) and \(0<s\leq1\), there exists a constant \(c>0\), depending on \(p,s\), such that every real cube function \(g\), under the usual measurability and square-integrability conditions and with \(N_{p,s}(g)^2<\infty\), satisfies
\[
\sum_{a\in\mathbb Z^p}
\left(1+\frac{\pi^2\|a\|^2}{p}\right)^s
|\widehat h(a)|^2
\leq c\,N_{p,s}(g)^2.
\]

For every \(d\geq2\), \(0<s\leq1\), and \(m\in\mathcal H_{d,s}\) as in \cref{obj:def:sobolev-class}, define the reflected function \(F\), its Fourier coefficients, and the frequency complexity by
\[
F=m\circ b_d,\qquad
\widehat F(k)=\int_{\mathbb T_2^d}F(w)
  \exp(-\mathrm i\pi k\cdot w)\,dw,
\]
\[
\operatorname{supp}(k)=\{j:k_j\neq0\},\qquad
t(k)=
\begin{cases}
\displaystyle\frac{\sum_j k_j^2}{|\operatorname{supp}(k)|},&k\neq0,\\
0,&k=0.
\end{cases}
\]
Then, with \(C_s=3072\),
\[
\sum_{\substack{k\in\mathbb Z^d\\
1\leq|\operatorname{supp}(k)|\leq2}}
(1+\pi^2t(k))^s|\widehat F(k)|^2
\leq C_s.
\]

For every pair of nonnegative integers \(n,d\), let \(X\) satisfy \cref{obj:ass:uniform-draw} on an arbitrary measurable sampling space. Independently of \(X\), reflect each coordinate of each \(X_i\) by an independent fair sign, and add an independent common Haar shift \(T\in\mathbb T_2^d\), obtaining \(W_i\) modulo \(2\). The joint law of \(W=(W_i)_{i=1}^n\) and \(T\) is the product of \(n+1\) normalized Haar laws. Thus the \(W_i\) are independent Haar draws and \(W\) is independent of \(T\).

The following identity holds under these conditions:
\begin{itemize}
\item \textbf{(Parameters.)} \(n\geq2\), \(d\geq2\), and \(0<s\leq1\).
\item \textbf{(Prognostic function.)} \(m\in\mathcal H_{d,s}\) as in \cref{obj:def:sobolev-class}.
\item \textbf{(Sampling.)} \(X\) satisfies \cref{obj:ass:uniform-draw} on an arbitrary measurable sampling space, and \(W,T\) are constructed as above.
\item \textbf{(Assignment.)} \(\pi\) is any Borel probability kernel from \((\mathbb T_2^d)^n\) to \(\mathcal Z_n=\{-1,1\}^n\), and, conditional on the sample, reflection signs, and shift, \(Z\) has law \(\pi(W)\).
\end{itemize}
With the torus character \(e_k(w)=\exp(\mathrm i\pi k\cdot w)\), one has
\[
\mathbb E\left|\sum_{i=1}^n Z_i m(X_i)\right|^2
=
\sum_{k\in\mathbb Z^d}
|\widehat F(k)|^2\,
\mathbb E\left|\sum_{i=1}^n Z_i e_k(W_i)\right|^2.
\]
On the left, expectation includes the original sample, independent reflection signs, common shift, and conditional assignment law; on the right, it includes the independent Haar sample \(W\) and conditional assignment law \(\pi(W)\). The equality is understood in the extended nonnegative reals.
\end{lemmav}

\begin{proof}[Proof of \cref{obj:lem:reflection-transfer}]
\begin{enumerate}
\item For \(p\in\{1,2\}\) and \(0<s\leq1\), \cref{obj:lem:exact-reflection-budget} gives, under the stated measurability, square-integrability, and finite restriction-norm hypotheses,
\[
\sum_{a\in\mathbb Z^p}
\left(1+\frac{\pi^2\|a\|^2}{p}\right)^s
|\widehat h(a)|^2
\leq N_{p,s}(g)^2.
\]
Thus the first assertion holds with \(c=1>0\).
% lean: CausalSmith.Experimentation.BivariateSobolevDesignCapacity.reflection_transfer

\item For \(d\geq2\), \(0<s\leq1\), and
\(m\in\mathcal H_{d,s}\), the pooled conclusion of
\cref{obj:lem:exact-reflection-budget} gives
\[
\sum_{\substack{k\in\mathbb Z^d\\
1\leq|\operatorname{supp}(k)|\leq2}}
(1+\pi^2t(k))^s|\widehat F(k)|^2
\leq1\leq3072=C_s.
\]
This proves the second assertion.
% lean: CausalSmith.Experimentation.BivariateSobolevDesignCapacity.exact_reflection_budget

\item We next establish the joint sampling law for arbitrary nonnegative
integers \(n,d\). Write
\[
\mu_{\mathrm H,d}
=\text{normalized Haar measure on }\mathbb T_2^d,
\qquad
\nu_{\mathrm H,n,d}=\mu_{\mathrm H,d}^{\otimes n}.
\]
Let the independent reflection signs and the reflected sample be
\[
\eta=(\eta_{ij})_{\substack{1\leq i\leq n\\1\leq j\leq d}}
\in\{-1,1\}^{n\times d},
\qquad
\mathbb P(\eta_{ij}=1)=\mathbb P(\eta_{ij}=-1)=\frac12,
\]
\[
\widetilde X_i=(\eta_{ij}X_{ij})_{j=1}^d
\pmod{2},
\qquad
W_i=\widetilde X_i+T.
\]
The sampling measure has total mass one because its pushforward under
\(X\) is the probability law in \cref{obj:ass:uniform-draw}.

For every nonnegative Borel function
\(\varphi:\mathbb T_2^1\to[0,\infty]\), a single reflected coordinate
satisfies
\[
\begin{aligned}
\mathbb E\varphi(\widetilde X_{ij})
&=\frac12\int_0^1\varphi(u)\,du
  +\frac12\int_0^1\varphi(2-u)\,du\\
&=\frac12\int_0^2\varphi(u)\,du
 =\int_{\mathbb T_2^1}\varphi\,d\mu_{\mathrm H,1}.
\end{aligned}
\]
Here the change of variables in the second integral combines the two
halves of the period-two circle; endpoints have Lebesgue measure zero.
Taking products over coordinates and units, using independence of the
uniform coordinates and reflection signs, gives
\[
\operatorname{Law}(\widetilde X)=\nu_{\mathrm H,n,d},
\qquad
\operatorname{Law}(\widetilde X,T)
=\nu_{\mathrm H,n,d}\otimes\mu_{\mathrm H,d}.
\]
For each \(\tau\in\mathbb T_2^d\), coordinatewise Haar translation
invariance gives
\[
\bigl(\mathbf w\mapsto
(\mathbf w_i+\tau)_{i=1}^n\bigr)_{\#}
\nu_{\mathrm H,n,d}
=\nu_{\mathrm H,n,d}.
\]
Consequently, for Borel sets
\(E_{\mathrm H}\subseteq(\mathbb T_2^d)^n\) and
\(C_{\mathrm H}\subseteq\mathbb T_2^d\),
\[
\begin{aligned}
\mathbb P(W\in E_{\mathrm H},T\in C_{\mathrm H})
&=\int_{C_{\mathrm H}}
  \int_{(\mathbb T_2^d)^n}
  \mathbf1_{E_{\mathrm H}}
  \bigl((\mathbf w_i+\tau)_{i=1}^n\bigr)
  \,d\nu_{\mathrm H,n,d}(\mathbf w)
  \,d\mu_{\mathrm H,d}(\tau)\\
&=\nu_{\mathrm H,n,d}(E_{\mathrm H})
  \mu_{\mathrm H,d}(C_{\mathrm H}).
\end{aligned}
\]
Thus
\[
\operatorname{Law}(W,T)
=\nu_{\mathrm H,n,d}\otimes\mu_{\mathrm H,d}.
\]
For \(n=0\) or \(d=0\), empty products are point masses on the
corresponding unique empty tuple, and the same product and translation
identities apply.
% lean: CausalSmith.Experimentation.BivariateSobolevDesignCapacity.reflectionLaw_map_shiftedSample

\item Now impose the hypotheses of the final assertion. Folding either
reflection of a cube coordinate recovers that coordinate:
\[
b_1(u\bmod2)=u,
\qquad
b_1(-u\bmod2)=u,
\qquad 0\leq u\leq1.
\]
For the second identity, \(u>0\) has representative \(2-u\), whose fold
is \(u\); at \(u=0\), the representative and its fold are both zero.
Since the sampled coordinates belong to the cube almost surely,
\[
b_d(\widetilde X_i)=X_i,
\qquad
F(W_i-T)=F(\widetilde X_i)=m(X_i)
\quad\text{almost surely},\qquad 1\leq i\leq n.
\]
The folding map and \(F\) are Borel measurable. Integrating the
conditional assignment law \(\pi(W)\), and then using the joint law
proved above, therefore yields
\[
\begin{aligned}
\mathbb E\left|\sum_{i=1}^n Z_i m(X_i)\right|^2
={}&
\int_{(\mathbb T_2^d)^n}\int_{\mathbb T_2^d}
\int_{\mathcal Z_n}
\left|\sum_{i=1}^n z_i F(\mathbf w_i-\tau)\right|^2
\,\pi(\mathbf w,dz)\,
d\mu_{\mathrm H,d}(\tau)\,
d\nu_{\mathrm H,n,d}(\mathbf w).
\end{aligned}
\]
Because \(F\) is real-valued, the absolute value here also equals the
complex modulus of the same sum.
% lean: CausalSmith.Experimentation.BivariateSobolevDesignCapacity.reflection_kernel_fourier_identity

\item Tonelli's theorem first interchanges the sign and shift integrals:
\[
\begin{aligned}
&\int_{\mathbb T_2^d}\int_{\mathcal Z_n}
\left|\sum_{i=1}^n z_i F(\mathbf w_i-\tau)\right|^2
\,\pi(\mathbf w,dz)\,d\mu_{\mathrm H,d}(\tau)\\
&\qquad=
\int_{\mathcal Z_n}\int_{\mathbb T_2^d}
\left|\sum_{i=1}^n z_i F(\mathbf w_i-\tau)\right|^2
\,d\mu_{\mathrm H,d}(\tau)\,\pi(\mathbf w,dz).
\end{aligned}
\]
We evaluate the inner shift integral by computing its Fourier
coefficients.

The folding normalization gives
\[
\int_{\mathbb T_2^d}|F(y)|^2\,d\mu_{\mathrm H,d}(y)
=\int_{[0,1]^d}|m(x)|^2\,dx<\infty.
\]
Indeed, fix the remaining coordinates
\(x_{\mathrm{rest}}\in[0,1]^{d-1}\) and define the first-coordinate slice by
\[
m_{\mathrm{slice},x_{\mathrm{rest}}}:[0,1]\to\mathbb R,
\qquad
m_{\mathrm{slice},x_{\mathrm{rest}}}(u_{\mathrm{slice}})
=m(u_{\mathrm{slice}},x_{\mathrm{rest}}),
\qquad u_{\mathrm{slice}}\in[0,1].
\]
The two reflected halves contribute
\[
\frac12\int_0^1
|m_{\mathrm{slice},x_{\mathrm{rest}}}(u_{\mathrm{slice}})|^2
\,du_{\mathrm{slice}}
+\frac12\int_1^2
|m_{\mathrm{slice},x_{\mathrm{rest}}}(2-u_{\mathrm{slice}})|^2
\,du_{\mathrm{slice}}
=\int_0^1
|m_{\mathrm{slice},x_{\mathrm{rest}}}(u_{\mathrm{slice}})|^2
\,du_{\mathrm{slice}}.
\]
The same change of variables applies to a slice in any coordinate.
Integrating over the remaining coordinates and iterating this identity
in each coordinate, using Tonelli's theorem, proves the displayed
\(d\)-dimensional identity.

For every
\(\mathbf w\in(\mathbb T_2^d)^n\) and \(z\in\mathcal Z_n\), define
\[
S_{\mathbf w,z}(\tau)
=\sum_{i=1}^n z_i F(\mathbf w_i-\tau),
\qquad \tau\in\mathbb T_2^d,
\]
\[
\widehat S_{\mathbf w,z}(k)
=\int_{\mathbb T_2^d}
e_{-k}(\tau)S_{\mathbf w,z}(\tau)
\,d\mu_{\mathrm H,d}(\tau),
\qquad k\in\mathbb Z^d.
\]
The map \(\tau\mapsto\mathbf w_i-\tau\) preserves Haar measure, so
\[
\int_{\mathbb T_2^d}|F(\mathbf w_i-\tau)|^2
\,d\mu_{\mathrm H,d}(\tau)
=\int_{\mathbb T_2^d}|F(y)|^2\,d\mu_{\mathrm H,d}(y)<\infty.
\]
Each summand, and hence \(S_{\mathbf w,z}\), belongs to \(L^2\).

Changing variables \(y=\mathbf w_i-\tau\) and using the character
identity
\[
e_{-k}(\mathbf w_i-y)=e_{-k}(\mathbf w_i)e_k(y)
\]
gives
\[
\begin{aligned}
\int_{\mathbb T_2^d}
e_{-k}(\tau)F(\mathbf w_i-\tau)\,d\mu_{\mathrm H,d}(\tau)
&=e_{-k}(\mathbf w_i)
  \int_{\mathbb T_2^d}e_k(y)F(y)\,d\mu_{\mathrm H,d}(y)\\
&=e_{-k}(\mathbf w_i)\widehat F(-k).
\end{aligned}
\]
Linearity over the finite sum therefore gives
\[
\widehat S_{\mathbf w,z}(k)
=\widehat F(-k)\sum_{i=1}^n z_i e_{-k}(\mathbf w_i).
\]

The normalized characters form an orthonormal basis of
\(L^2(\mu_{\mathrm H,d})\). Their orthogonality follows coordinatewise:
\[
\int_{\mathbb T_2^d}
e_k(y)\overline{e_{k'}(y)}\,d\mu_{\mathrm H,d}(y)
=
\prod_{j=1}^d
\left(\frac12\int_0^2
  \exp\bigl(\mathrm i\pi(k_j-k'_j)u\bigr)\,du\right)
=
\begin{cases}
1,&k=k',\\
0,&k\neq k',
\end{cases}
\quad k,k'\in\mathbb Z^d.
\]
Trigonometric polynomials are uniformly dense in continuous functions
on the torus, and continuous functions are dense in \(L^2\); thus
\[
\overline{\operatorname{span}\{e_k:k\in\mathbb Z^d\}}^{\,L^2}
=L^2(\mu_{\mathrm H,d}).
\]
The Hilbert-space norm identity for this orthonormal basis, followed by
the bijective reindexing \(k\mapsto-k\), now gives
\[
\begin{aligned}
\int_{\mathbb T_2^d}|S_{\mathbf w,z}(\tau)|^2
\,d\mu_{\mathrm H,d}(\tau)
&=\sum_{k\in\mathbb Z^d}|\widehat S_{\mathbf w,z}(k)|^2\\
&=\sum_{k\in\mathbb Z^d}
|\widehat F(-k)|^2
\left|\sum_{i=1}^n z_i e_{-k}(\mathbf w_i)\right|^2\\
&=\sum_{k\in\mathbb Z^d}
|\widehat F(k)|^2
\left|\sum_{i=1}^n z_i e_k(\mathbf w_i)\right|^2.
\end{aligned}
\]
All terms are nonnegative, so this identity also holds as an equality
of extended nonnegative integrals and sums.
% lean: CausalSmith.Experimentation.BivariateSobolevDesignCapacity.commonShift_parseval_lintegral

\item Substitute the shift identity into the transported expectation.
Tonelli's theorem moves the nonnegative frequency sum through both the
assignment-kernel integral and the Haar-sample integral, giving
\[
\begin{aligned}
\mathbb E\left|\sum_{i=1}^n Z_i m(X_i)\right|^2
&=
\int_{(\mathbb T_2^d)^n}\int_{\mathcal Z_n}
\sum_{k\in\mathbb Z^d}
|\widehat F(k)|^2
\left|\sum_{i=1}^n z_i e_k(\mathbf w_i)\right|^2
\,\pi(\mathbf w,dz)\,d\nu_{\mathrm H,n,d}(\mathbf w)\\
&=
\sum_{k\in\mathbb Z^d}|\widehat F(k)|^2
\int_{(\mathbb T_2^d)^n}\int_{\mathcal Z_n}
\left|\sum_{i=1}^n z_i e_k(\mathbf w_i)\right|^2
\,\pi(\mathbf w,dz)\,d\nu_{\mathrm H,n,d}(\mathbf w)\\
&=
\sum_{k\in\mathbb Z^d}|\widehat F(k)|^2
\mathbb E\left|\sum_{i=1}^n Z_i e_k(W_i)\right|^2.
\end{aligned}
\]
The last expectation uses the independent Haar sample and conditional
assignment law \(\pi(W)\), as asserted.
% lean: CausalSmith.Experimentation.BivariateSobolevDesignCapacity.commonShift_kernel_parseval
\end{enumerate}
\end{proof}

The transfer in \cref{obj:lem:reflection-transfer} separates population regularity from assignment performance. Squared Fourier coefficients describe the reflected prognosis; the expected squared signed character sums describe the assignment rule. Orthogonality of the common-shift phases removes cross-frequency terms, even when the assignment signs depend on the entire transformed sample. The dependence condition is therefore substantive: conditional assignment uses \(W\), while the independent shift supplies the averaging that diagonalizes the criterion.

The identity applies to every Borel representative allowed by \cref{obj:def:sobolev-class}. Its square-integrable formulation accommodates Fourier expansions in the \(L^2\) sense, and uniform sampling makes changes on cube null sets immaterial to the original-sample loss. Reflection consequently carries nonperiodic component regularity into the spectral assignment criterion with the boundary behavior permitted by the restriction norm. For the attaining comparison, the sharper unit budget in \cref{obj:lem:exact-reflection-budget} supplies the normalization within this transfer.

\subsection{Ordered rounding and the attaining bound}

The spectral rows are bounded real functions ordered by increasing frequency complexity. Their assignment guarantees must hold conditional on each realized input matrix, so that auxiliary transformations can subsequently be averaged. The following statement gives the measurable realization, finite completion, fairness, and second-moment bounds for the procedure in \cref{obj:def:ordered-boundary-rounding}.

\begin{lemmav}{lem:ordered-prefix-rounding}[Ordered prefix rounding guarantees]
Let \(n\geq0\) be an integer, let \(\Lambda\in\mathbb R\), and let \(a_1,a_2,\ldots\in\mathbb R^n\) be deterministic ordered rows satisfying
\[
|(a_h)_i|\leq\Lambda
\qquad(h\geq1,\ 1\leq i\leq n).
\]
Supply the first \(K=\lfloor n/4\rfloor\) rows to the procedure in \cref{obj:def:ordered-boundary-rounding}, and write \(Z\in\{-1,1\}^n\) for its output signs. The output map is jointly Borel measurable in the input matrix \(a\in\mathbb R^{K\times n}\) and the real seed vector of length \(3n\). For every seed vector in \([0,1]^{3n}\), the fractional state after \(n\) moves satisfies
\[
u_i=Z_i\qquad(1\leq i\leq n).
\]
With the \(3n\) seeds drawn independently and uniformly from \([0,1]\),
\[
\mathbb E[Z_i]=0
\qquad(1\leq i\leq n),
\qquad
\mathbb E\!\left[(a_h^{\mathsf T}Z)^2\right]
\leq32\Lambda^2\min\{h,n\}
\qquad(h\geq1).
\]
\end{lemmav}

\begin{proof}[Proof of \cref{obj:lem:ordered-prefix-rounding}]
\begin{enumerate}
\item
We first establish joint measurability. For an arbitrary input matrix \(a\in\mathbb R^{K\times n}\) and seed vector
\[
\boldsymbol\eta=(\eta_1,\ldots,\eta_{3n})\in\mathbb R^{3n},
\]
run \cref{obj:def:ordered-boundary-rounding}, using \(\eta_{2\tau+1}\) for direction selection and \(\eta_{2\tau+2}\) for the boundary coin at move \(\tau\), where \(0\leq\tau<n\). A terminal coordinate round uses the second of these seeds. After exhaustion, pad the procedure with identity moves.

Write \(u^{(\tau)}\) for the fractional vector after \(\tau\) moves, and define its active set and count by
\[
u^{(0)}=0,\qquad
\mathcal I_\tau=\{i\in\{1,\ldots,n\}:|u_i^{(\tau)}|<1\},
\qquad
r_\tau=|\mathcal I_\tau|.
\]
More explicitly, on a nonterminal move the recursion is
\[
u^{(\tau+1)}
=
\begin{cases}
u^{(\tau)}+\delta_b^+v_b,
&\displaystyle \eta_{2\tau+2}<\frac{\delta_b^-}{\delta_b^++\delta_b^-},\\[4pt]
u^{(\tau)}-\delta_b^-v_b,
&\displaystyle \eta_{2\tau+2}\geq\frac{\delta_b^-}{\delta_b^++\delta_b^-},
\end{cases}
\]
where \(b\), the basis vectors, and the boundary lengths are computed from the current state as specified in \cref{obj:def:ordered-boundary-rounding}. On a terminal move, put
\[
i_{\mathrm{term},\tau}=\min\mathcal I_\tau
\]
and replace that coordinate by \(+1\) exactly when
\(\eta_{2\tau+2}<(1+u_{i_{\mathrm{term},\tau}}^{(\tau)})/2\), and by \(-1\) otherwise. This index is used exclusively when the active set is nonempty. An empty active set gives \(u^{(\tau+1)}=u^{(\tau)}\). The output is
\[
Z_i(a,\boldsymbol\eta)
=
\begin{cases}
1,&u_i^{(n)}(a,\boldsymbol\eta)>0,\\
-1,&u_i^{(n)}(a,\boldsymbol\eta)\leq0.
\end{cases}
\]

All active-set and phase tests are Borel. Ordered Gram--Schmidt is a finite recursion: each residual is formed by sums and products, the zero-residual branch is Borel, and normalization on the positive-norm branch is Borel. Consequently the constraint projection and the ordered basis obtained from its projected coordinate vectors are Borel functions of the input and state. Boundary lengths are finite minima of sign-guarded ratios, and the finite inverse-distribution recursion uses Borel threshold tests. For the total formulas, an empty minimum is assigned zero and division by zero is assigned zero. For arbitrary real selection seeds, the same threshold recursion selects the first interval below zero and the last interval at or above one.

Thus induction over the finite move recursion gives joint Borel measurability of \(u^{(n)}\). In particular,
\[
\{(a,\boldsymbol\eta):Z_i(a,\boldsymbol\eta)=1\}
=
\{(a,\boldsymbol\eta):u_i^{(n)}(a,\boldsymbol\eta)>0\}
\]
is Borel for every coordinate. Since the sign space is finite, the output map is jointly Borel. When \(n=0\), it is the constant map into the singleton empty sign space.
% lean: CausalSmith.Experimentation.BivariateSobolevDesignCapacity.orderedBoundaryRounding_measurable

\item
We next prove termination. At a nonterminal move, define
\[
r_{\mathrm{con},\tau}
=\operatorname{rank}\bigl((a_{\ell i})_{1\leq\ell\leq q_1,\ i\in\mathcal I_\tau}\bigr),
\qquad
r_{\mathrm{null},\tau}=\dim(\operatorname{im}P).
\]
The orthonormalized constraint rows and the nullspace basis together span the active coordinate subspace. Hence
\[
r_{\mathrm{con},\tau}+r_{\mathrm{null},\tau}=r_\tau,
\qquad
r_{\mathrm{con},\tau}\leq q_1,
\qquad
r_{\mathrm{null},\tau}\geq r_\tau-q_1.
\]
At a new nonterminal phase, \(r_\tau\geq4\) and
\(q_1=\lfloor r_\tau/4\rfloor<r_\tau\). During a continuing phase with starting count \(r\),
\[
q_1=\lfloor r/4\rfloor
\leq\lfloor r/2\rfloor<r_\tau.
\]
Thus \(r_{\mathrm{null},\tau}>0\) on every nonterminal move.

Every basis vector is nonzero and supported on active coordinates. Therefore the sets defining its two boundary lengths are nonempty, and
\[
0<\delta_b^+<\infty,\qquad
0<\delta_b^-<\infty.
\]
The boundary minima ensure
\[
|u_i^{(\tau)}+\delta_b^+(v_b)_i|\leq1,
\qquad
|u_i^{(\tau)}-\delta_b^-(v_b)_i|\leq1
\]
for every coordinate, with equality for at least one previously active coordinate on each branch. Already inactive coordinates have zero direction entries. A terminal round likewise fixes one active coordinate at a boundary. Thus, whenever \(r_\tau>0\), a move preserves the cube and strictly decreases the active count; when \(r_\tau=0\), the move is the identity. Induction yields
\[
|u_i^{(\tau)}|\leq1,
\qquad
r_\tau\leq\max\{n-\tau,0\}
\qquad(0\leq\tau\leq n).
\]
In particular \(r_n=0\), so \(|u_i^{(n)}|=1\). If \(u_i^{(n)}>0\), then \(u_i^{(n)}=1=Z_i\); otherwise \(u_i^{(n)}=-1=Z_i\). This proves
\[
u_i^{(n)}=Z_i
\qquad(1\leq i\leq n).
\]
The argument applies to every real seed vector, and hence to every seed vector in the stated cube.
% lean: CausalSmith.Experimentation.BivariateSobolevDesignCapacity.roundingIteration_terminal_eq_sign

\item
Now supply \(a_\ell^{\mathsf T}\), \(1\leq\ell\leq K\), as the input rows and draw the seeds independently and uniformly. Define
\[
\mathcal F_\tau=\sigma(\eta_1,\ldots,\eta_{2\tau}),
\qquad
\Delta^{(\tau)}=u^{(\tau+1)}-u^{(\tau)}
\qquad(0\leq\tau<n).
\]
Conditional on the current state and the selected basis index \(b\), the boundary coin gives
\[
\mathbb E[\Delta^{(\tau)}\mid\mathcal F_\tau,b]
=
\left(
\frac{\delta_b^-\delta_b^+}{\delta_b^++\delta_b^-}
-
\frac{\delta_b^+\delta_b^-}{\delta_b^++\delta_b^-}
\right)v_b
=0.
\]
For a terminal coordinate round, its conditional mean is
\[
\frac{1+u_i^{(\tau)}}2
-
\frac{1-u_i^{(\tau)}}2
=u_i^{(\tau)}.
\]
Identity moves also have zero increment. Integrating the fresh coin and then the remaining seeds therefore gives
\[
\mathbb E[u_i^{(\tau+1)}]
=\mathbb E[u_i^{(\tau)}].
\]
Starting from \(u^{(0)}=0\) and using the terminal identity from the preceding step, we obtain
\[
\mathbb E[Z_i]=\mathbb E[u_i^{(n)}]=0.
\]
% lean: CausalSmith.Experimentation.BivariateSobolevDesignCapacity.orderedBoundaryRounding_fair

\item
Fix an integer \(h\geq1\). Define the predictable unprotected events by
\[
\mathcal E_\tau=\{r_\tau<4h\}
\qquad(0\leq\tau<n).
\]
These events are Borel and depend exclusively on seeds read before move \(\tau\).

We establish the row-to-energy comparison on each move. On a nonterminal move, the scalar boundary-coin second moment is
\[
\frac{\delta_b^-(\delta_b^+)^2+\delta_b^+(\delta_b^-)^2}
{\delta_b^++\delta_b^-}
=\delta_b^+\delta_b^-=t_b.
\]
Averaging over the direction probabilities \(t_b^{-1}/S\) consequently gives
\[
\mathbb E\!\left[(a_h^{\mathsf T}\Delta^{(\tau)})^2\mid\mathcal F_\tau\right]
=
S^{-1}\sum_{b=1}^{r_{\mathrm{null},\tau}}
(a_h^{\mathsf T}v_b)^2.
\]
Here \(S\) and the basis refer to the current move. Orthonormality and active support imply
\[
\begin{aligned}
0
&\leq
\left\|
\bigl((a_h)_i\mathbf1_{\{i\in\mathcal I_\tau\}}\bigr)_{i=1}^n
-
\sum_{b=1}^{r_{\mathrm{null},\tau}}
(a_h^{\mathsf T}v_b)v_b
\right\|_2^2\\
&=
\sum_{i\in\mathcal I_\tau}(a_h)_i^2
-
\sum_{b=1}^{r_{\mathrm{null},\tau}}(a_h^{\mathsf T}v_b)^2.
\end{aligned}
\]
Thus the row-entry bound yields
\[
\mathbb E\!\left[(a_h^{\mathsf T}\Delta^{(\tau)})^2\mid\mathcal F_\tau\right]
\leq S^{-1}\Lambda^2r_\tau.
\]
The same boundary-coin calculation, together with zero conditional mean and unit basis norms, gives
\[
\mathbb E\!\left[
\|u^{(\tau+1)}\|_2^2-\|u^{(\tau)}\|_2^2
\mid\mathcal F_\tau\right]
=S^{-1}r_{\mathrm{null},\tau}\geq0.
\]

At a phase refresh,
\(2\lfloor r_\tau/4\rfloor<r_\tau\). On a continuing phase,
\[
2q_1=2\lfloor r/4\rfloor
\leq\lfloor r/2\rfloor<r_\tau.
\]
Combining \(2q_1<r_\tau\) with
\(r_{\mathrm{null},\tau}\geq r_\tau-q_1\) gives
\[
r_\tau\leq4r_{\mathrm{null},\tau}.
\]
Therefore the conditional row variance is at most \(4\Lambda^2\) times the conditional norm-energy increase.

On a terminal move, the corresponding two quantities are
\[
\mathbb E\!\left[(a_h^{\mathsf T}\Delta^{(\tau)})^2\mid\mathcal F_\tau\right]
=(a_h)_{i_{\mathrm{term},\tau}}^2
\bigl(1-(u_{i_{\mathrm{term},\tau}}^{(\tau)})^2\bigr),
\]
and
\[
\mathbb E\!\left[
\|u^{(\tau+1)}\|_2^2-\|u^{(\tau)}\|_2^2
\mid\mathcal F_\tau\right]
=1-(u_{i_{\mathrm{term},\tau}}^{(\tau)})^2\geq0.
\]
These satisfy the same comparison. Both quantities vanish on an identity move.

Multiplying by the predictable indicator and integrating therefore gives
\[
\mathbb E\!\left[
\bigl(\mathbf1_{\mathcal E_\tau}a_h^{\mathsf T}\Delta^{(\tau)}\bigr)^2
\right]
\leq
4\Lambda^2
\mathbb E\!\left[
\mathbf1_{\mathcal E_\tau}
\bigl(\|u^{(\tau+1)}\|_2^2-\|u^{(\tau)}\|_2^2\bigr)
\right].
\]
For each coordinate, expanding the square and integrating the fresh zero-mean coin gives
\[
\mathbb E\!\left[
\bigl(\mathbf1_{\mathcal E_\tau}\Delta_i^{(\tau)}\bigr)^2
\right]
=
\mathbb E\!\left[
\mathbf1_{\mathcal E_\tau}
\bigl((u_i^{(\tau+1)})^2-(u_i^{(\tau)})^2\bigr)
\right].
\]
Summing first over coordinates and then over moves proves
\[
\sum_{\tau=0}^{n-1}
\mathbb E\!\left[
\bigl(\mathbf1_{\mathcal E_\tau}a_h^{\mathsf T}\Delta^{(\tau)}\bigr)^2
\right]
\leq
4\Lambda^2
\sum_{\tau=0}^{n-1}\sum_{i=1}^n
\mathbb E\!\left[
\bigl(\mathbf1_{\mathcal E_\tau}\Delta_i^{(\tau)}\bigr)^2
\right].
\]
% lean: CausalSmith.Experimentation.BivariateSobolevDesignCapacity.roundingIteration_masked_row_variance_sum_le

\item
Define the first unprotected time by
\[
\tau_*=\min\{\tau\in\{0,\ldots,n\}:r_\tau<4h\}.
\]
This finite minimum exists because \(r_n=0\) and \(h\geq1\). It is measurable. Since active counts decrease,
\[
\mathbf1_{\mathcal E_\tau}
=\mathbf1_{\{\tau_*\leq\tau\}}
\qquad(0\leq\tau<n),
\]
and
\[
r_{\tau_*}<4h,\qquad
r_{\tau_*}\leq n,\qquad
r_{\tau_*}\leq4\min\{h,n\}.
\]
The coordinate-energy identity from the preceding step now telescopes over this predictable terminal interval:
\[
\begin{aligned}
\sum_{\tau=0}^{n-1}\sum_{i=1}^n
\mathbb E\!\left[
\bigl(\mathbf1_{\mathcal E_\tau}\Delta_i^{(\tau)}\bigr)^2
\right]
&=
\mathbb E\!\left[
\sum_{\tau=0}^{n-1}
\mathbf1_{\{\tau_*\leq\tau\}}
\bigl(\|u^{(\tau+1)}\|_2^2-\|u^{(\tau)}\|_2^2\bigr)
\right]\\
&=
\mathbb E\!\left[\|u^{(n)}\|_2^2-\|u^{(\tau_*)}\|_2^2\right].
\end{aligned}
\]
Coordinates inactive at \(\tau_*\) remain fixed, while every endpoint coordinate has square at most one. Hence, pointwise,
\[
\|u^{(n)}\|_2^2-\|u^{(\tau_*)}\|_2^2
=
\sum_{i\in\mathcal I_{\tau_*}}
\bigl((u_i^{(n)})^2-(u_i^{(\tau_*)})^2\bigr)
\leq r_{\tau_*}.
\]
Consequently,
\[
\sum_{\tau=0}^{n-1}\sum_{i=1}^n
\mathbb E\!\left[
\bigl(\mathbf1_{\mathcal E_\tau}\Delta_i^{(\tau)}\bigr)^2
\right]
\leq\mathbb E[r_{\tau_*}]
\leq4\min\{h,n\}.
\]
The telescope also applies when \(\tau_*=n\), with an empty selected interval.
% lean: CausalSmith.Experimentation.BivariateSobolevDesignCapacity.roundingIteration_unprotected_coordinate_variance_sum_le

\item
It remains to identify the final row second moment with the selected move variances. When \(r_\tau\geq4h\), we have \(h\leq K\). At a phase refresh,
\[
q_1=\lfloor r_\tau/4\rfloor\geq h.
\]
During a continuing phase, its starting count is at least \(r_\tau\), so its retained prefix again contains row \(h\). Every direction is orthogonal to that row. Thus
\[
r_\tau\geq4h
\quad\Longrightarrow\quad
a_h^{\mathsf T}\Delta^{(\tau)}=0.
\]

Zero conditional increment means also give
\[
\mathbb E\!\left[
(a_h^{\mathsf T}u^{(\tau)})
(a_h^{\mathsf T}\Delta^{(\tau)})
\right]=0.
\]
Expanding the row square therefore yields
\[
\mathbb E[(a_h^{\mathsf T}u^{(\tau+1)})^2]
-
\mathbb E[(a_h^{\mathsf T}u^{(\tau)})^2]
=
\mathbb E[(a_h^{\mathsf T}\Delta^{(\tau)})^2].
\]
Summing from the zero initial state, using \(u^{(n)}=Z\), and discarding the identically zero protected increments gives
\[
\mathbb E[(a_h^{\mathsf T}Z)^2]
=
\sum_{\tau=0}^{n-1}
\mathbb E\!\left[
\bigl(\mathbf1_{\mathcal E_\tau}a_h^{\mathsf T}\Delta^{(\tau)}\bigr)^2
\right].
\]
Combining the two preceding bounds, with \(4\Lambda^2\geq0\), proves
\[
\begin{aligned}
\mathbb E[(a_h^{\mathsf T}Z)^2]
&\leq
4\Lambda^2
\sum_{\tau=0}^{n-1}\sum_{i=1}^n
\mathbb E\!\left[
\bigl(\mathbf1_{\mathcal E_\tau}\Delta_i^{(\tau)}\bigr)^2
\right]\\
&\leq16\Lambda^2\min\{h,n\}\\
&\leq32\Lambda^2\min\{h,n\}.
\end{aligned}
\]
For \(h>K\), every move is unprotected because \(r_\tau\leq n<4h\), so the same argument covers all subsequent rows. When \(n=0\), all row sums and move sums are empty and the bound reads \(0\leq0\).
% lean: CausalSmith.Experimentation.BivariateSobolevDesignCapacity.orderedBoundaryRounding_row_bound
\end{enumerate}
\end{proof}

Joint Borel measurability in \cref{obj:lem:ordered-prefix-rounding} makes the finite randomized map suitable for covariate-dependent assignment. The ordered projection and basis conventions in \cref{obj:def:ordered-boundary-rounding} specify its choices on matrices of every rank, including boundary cases. Completion holds for every allowed seed vector; fairness and the second-moment bounds concern independent uniform seeds. Integrating this map over the rounding seeds therefore gives a Borel assignment law with fair marginals for every fixed input matrix.

The boundary probabilities preserve coordinate means, while the active constraints protect an ordered prefix of rows. Earlier rows receive longer protection as the active set contracts. The bound \(32\Lambda^2\min\{h,n\}\) quantifies this allocation and applies to the entire ordered row family, including rows beyond the supplied prefix. This last feature connects the finite input matrix to the full spectral budget.

For the prescribed cosine and sine rows, the entry bound is \(\Lambda=\sqrt2\). Conditional application of \cref{obj:lem:ordered-prefix-rounding} controls their signed second moments, and \cref{obj:lem:reflection-transfer} transfers those controls to the original prognostic imbalance. Ordering by frequency complexity aligns the strongest row bounds with the spectral weights in \cref{obj:lem:exact-reflection-budget}. Together, these are the assignment and regularity ingredients of the spectral-branch comparison in \cref{obj:thm:log-free-bivariate-capacity}.

The construction remains admissible after averaging over independent reflection coins and the common shift. Conditional fairness holds for every transformed input, and the randomization uses the original covariates and independent auxiliary seeds as specified in \cref{obj:def:capacity-handle}. Hence the resulting original-sample assignment belongs to \cref{obj:def:design-class}.

When \(n\leq B\), where \(B=\binom d2\) is the coordinate-pair count, the independent-sign branch is controlled by the cube \(L^2\) bound in \cref{obj:lem:exact-reflection-budget}. When \(B<n\), the ordered spectral branch gives the smoothness-dependent comparison. The assembled guarantee, stated in \cref{obj:thm:log-free-bivariate-capacity}, is
\[
\sup_{m\in\mathcal H_{d,s}}
\mathcal L_{n,d,s}(\pi^\star_{n,d},m)
\leq
3072\min\{1,(B/n)^s\},
\qquad n,d\geq2,\quad 0<s\leq1.
\]
Here the assignment \(\pi^\star_{n,d}\) is the construction in \cref{obj:def:capacity-handle}, and the loss is the criterion specified in \cref{obj:thm:log-free-bivariate-capacity}. Smoothness enters through the spectral weights used to evaluate this single assignment procedure, yielding simultaneous attainment across the stated smoothness range.


\section{Independent-prior converse}

The converse combines a finite prior supported on the prognostic class with a lower bound for signing independent isotropic vectors. The prior in \cref{obj:def:cosine-prior} is drawn before the covariate sample, so its average loss provides a comparison with the worst-case loss over fixed functions in \cref{obj:def:criterion}.

To specify the \leanref{S-5}{pair-cosine coefficient family}, let \leanref{sym:L}{\ensuremath{L}} be a positive integer frequency cutoff. Index the features by \leanref{sym:alpha}{\ensuremath{\alpha=(j,\ell,a,b)}}, where \(1\leq j<\ell\leq d\) and \(1\leq a,b\leq L\), and set
\[
\phi_\alpha(x)=2\cos(\pi a x_j)\cos(\pi b x_\ell).
\]
The full feature vector \leanref{sym:V}{\ensuremath{V(x)}} collects these coordinates and has dimension \leanref{sym:M}{\ensuremath{M=\binom{d}{2}L^2}}. Draw the coefficient array \leanref{sym:xi}{\ensuremath{\xi}} with independent fair sign entries, independently of the covariates and assignment randomization. The normalization constant \leanref{sym:C0}{\ensuremath{C_0=\sqrt{48+32\pi^2}}} scales the signed feature sum by \(C_0L^s\sqrt M\). The following statement gives the resulting function explicitly and records the class membership and second-moment identity needed for the converse.

\begin{lemmav}{lem:cosine-prior}[Cosine prior legality and isotropy]
Let the parameters satisfy:
\begin{itemize}
\item \textbf{(Dimension.)} \(d\) is an integer with \(d\geq2\).
\item \textbf{(Frequency cutoff.)} \(L\) is an integer with \(L\geq1\).
\item \textbf{(Smoothness.)} \(s\) is real with \(0<s\leq1\).
\end{itemize}
Set \(M=\binom{d}{2}L^2\), the number of pair-cosine features, and \(C_0=\sqrt{48+32\pi^2}\), the prior normalization constant. For every coefficient-sign array \(\xi\) with entries in \(\{-1,1\}\), the function
\[
m_\xi(x)=
\frac{1}{C_0L^s\sqrt M}
\sum_{1\leq j<\ell\leq d}\sum_{a=1}^{L}\sum_{b=1}^{L}
\xi_{j\ell ab}\,2\cos(\pi a x_j)\cos(\pi b x_\ell),
\qquad x\in[0,1]^d,
\]
is centered under the product uniform measure and belongs to \(\mathcal H_{d,s}\) as defined in \cref{obj:def:sobolev-class}.

For every integer \(n\geq1\) and every covariate array \(X\) on a sampling space satisfying \cref{obj:ass:uniform-draw}, let \(X_1\) denote its first covariate vector and define the pair-cosine feature vector by
\[
V(x)=
\bigl(2\cos(\pi a x_j)\cos(\pi b x_\ell)\bigr)_{
1\leq j<\ell\leq d,\;1\leq a,b\leq L}.
\]
Then
\[
\mathbb E\!\left[V(X_1)V(X_1)^{\mathsf T}\right]=I_M,
\]
where \(I_M\) is the \(M\)-dimensional identity matrix. This identity includes distinct pair features sharing an original coordinate.
\end{lemmav}

\begin{proof}[Proof of \cref{obj:lem:cosine-prior}]
\begin{enumerate}
\item \emph{Centering and square integrability.}
Fix the stated parameters and a coefficient-sign array \(\xi\). Write
\[
\begin{aligned}
\mathcal I
&=\{(j,\ell,a,b):1\leq j<\ell\leq d,\ 1\leq a,b\leq L\},\\
B&=\binom d2,\qquad M=BL^2,\qquad
\gamma_{\mathrm{prior}}=\frac{1}{C_0L^s\sqrt M},\\
V_{j\ell ab}(x)
&=2\cos(\pi a x_j)\cos(\pi b x_\ell).
\end{aligned}
\]
Choosing an index amounts to choosing a coordinate pair and two frequencies, so
\[
|\mathcal I|=\binom d2\,L^2=M.
\]
Since \(d\geq2\) and \(L\geq1\), we have \(B\geq1\), \(M>0\), and
\(\gamma_{\mathrm{prior}}>0\).

For every positive integer frequency \(a\),
\[
\int_0^1\cos(\pi a t)\,dt
=\frac{\sin(\pi a)}{\pi a}=0.
\]
Product integration therefore gives, for every \((j,\ell,a,b)\in\mathcal I\),
\[
\int_{[0,1]^d}V_{j\ell ab}(x)\,dx
=2\left(\int_0^1\cos(\pi a t)\,dt\right)
  \left(\int_0^1\cos(\pi b t)\,dt\right)
=0.
\]
The function \(m_\xi\) is a finite linear combination of continuous bounded
features, hence is Borel measurable and square integrable. Integrating its
finite sum yields
\[
\int_{[0,1]^d}m_\xi(x)\,dx
=\gamma_{\mathrm{prior}}
  \sum_{\alpha\in\mathcal I}\xi_\alpha
  \int_{[0,1]^d}V_\alpha(x)\,dx
=0.
\]
% lean: CausalSmith.Experimentation.BivariateSobolevDesignCapacity.mXi_mean_zero

\item \emph{Canonical components and reconstruction.}
Use the same finite cosine expression to evaluate \(m_\xi\) on
\(\mathbb R^d\). For \(\tau_j,\tau_\ell\in\mathbb R\), the canonical marginal
formulas are
\[
\begin{aligned}
g_j(m_\xi)(\tau_j)
&=\int_{[0,1]^d}
  \left.m_\xi(x)\right|_{x_j=\tau_j}\,dx,\\
g_{j\ell}(m_\xi)(\tau_j,\tau_\ell)
&=\int_{[0,1]^d}
  \left.m_\xi(x)\right|_{x_j=\tau_j,\ x_\ell=\tau_\ell}\,dx
  -g_j(m_\xi)(\tau_j)-g_\ell(m_\xi)(\tau_\ell),
  \qquad j<\ell.
\end{aligned}
\]
Here substitution fixes the indicated coordinates before integration; the
corresponding unused integration variables contribute factors of one.

For a feature indexed by
\(\alpha=(j_\alpha,\ell_\alpha,a_\alpha,b_\alpha)\), fixing coordinate \(j\)
leaves a centered cosine factor whenever \(j=j_\alpha\) or
\(j=\ell_\alpha\). If \(j\) equals neither coordinate, the integral is the
feature's mean. Thus, in all three cases,
\[
g_j(V_\alpha)(\tau_j)=0,
\qquad
g_j(m_\xi)(\tau_j)=0.
\]
After fixing \(j<\ell\), a feature survives integration precisely when its
coordinate pair is \((j,\ell)\). In every other configuration, at least one
of its two positive-frequency factors remains integrated and has mean
zero. The ordering of both pairs excludes a reversed match. Consequently,
\[
g_{j\ell}(V_\alpha)(\tau_j,\tau_\ell)
=
\begin{cases}
2\cos(\pi a_\alpha\tau_j)\cos(\pi b_\alpha\tau_\ell),
  &(j_\alpha,\ell_\alpha)=(j,\ell),\\
0,&(j_\alpha,\ell_\alpha)\ne(j,\ell),
\end{cases}
\]
and finite-sum linearity gives
\[
g_{j\ell}(m_\xi)(\tau_j,\tau_\ell)
=\gamma_{\mathrm{prior}}
  \sum_{a,b=1}^L\xi_{j\ell ab}
  \,2\cos(\pi a\tau_j)\cos(\pi b\tau_\ell).
\]
Summing these pair effects counts each feature exactly once:
\[
\sum_{j=1}^d g_j(m_\xi)(x_j)
+\sum_{j<\ell}g_{j\ell}(m_\xi)(x_j,x_\ell)
=\gamma_{\mathrm{prior}}
  \sum_{\alpha\in\mathcal I}\xi_\alpha V_\alpha(x)
=m_\xi(x).
\]
This proves the decomposition required in \cref{obj:ass:order-two},
pointwise and hence almost everywhere.
% lean: CausalSmith.Experimentation.BivariateSobolevDesignCapacity.mXi_orderTwo

\item \emph{A cutoff extension and its weak derivatives.}
To establish the component restriction bound, take an arbitrary real array
\[
c^{\mathrm{loc}}\in\mathbb R^{\{1,\ldots,L\}^2}
\]
and define, for
\(u_{\mathrm{loc}}=(u_1^{\mathrm{loc}},u_2^{\mathrm{loc}})\in\mathbb R^2\),
\[
\begin{aligned}
U_{\mathrm{loc}}(u_{\mathrm{loc}})
&=\sum_{a,b=1}^L c^{\mathrm{loc}}_{ab}
  \,2\cos(\pi a u_1^{\mathrm{loc}})
     \cos(\pi b u_2^{\mathrm{loc}}),\\
E_{\mathrm{coef}}
&=\sum_{a,b=1}^L(c^{\mathrm{loc}}_{ab})^2.
\end{aligned}
\]
Define the scalar cutoff and a representative of its almost-everywhere
derivative on the whole real line by
\[
q_{\mathrm{cut}}(t)=
\begin{cases}
0,&t<-1,\\
t+1,&-1\leq t<0,\\
1,&0\leq t\leq1,\\
2-t,&1<t\leq2,\\
0,&t>2,
\end{cases}
\qquad
q'_{\mathrm{cut}}(t)=
\begin{cases}
1,&-1<t<0,\\
-1,&1<t<2,\\
0,&\text{otherwise}.
\end{cases}
\]
In particular, the chosen derivative representative is zero at all four
joins. Put
\[
\begin{aligned}
\chi_{\mathrm{cut}}(u_{\mathrm{loc}})
&=q_{\mathrm{cut}}(u_1^{\mathrm{loc}})
  q_{\mathrm{cut}}(u_2^{\mathrm{loc}}),\\
\chi_{\mathrm{cut},1}(u_{\mathrm{loc}})
&=q'_{\mathrm{cut}}(u_1^{\mathrm{loc}})
  q_{\mathrm{cut}}(u_2^{\mathrm{loc}}),\\
\chi_{\mathrm{cut},2}(u_{\mathrm{loc}})
&=q_{\mathrm{cut}}(u_1^{\mathrm{loc}})
  q'_{\mathrm{cut}}(u_2^{\mathrm{loc}}),\\
G_{\mathrm{cut}}&=\chi_{\mathrm{cut}}U_{\mathrm{loc}},\\
D_\rho
&=\chi_{\mathrm{cut},\rho}U_{\mathrm{loc}}
  +\chi_{\mathrm{cut}}\partial_\rho U_{\mathrm{loc}},
  \qquad \rho\in\{1,2\}.
\end{aligned}
\]
These formulas give
\[
0\leq\chi_{\mathrm{cut}}\leq1,
\qquad
\chi_{\mathrm{cut},1}^2+\chi_{\mathrm{cut},2}^2\leq2.
\]
The extension \(G_{\mathrm{cut}}\) equals \(U_{\mathrm{loc}}\) on
\([0,1]^2\), is continuous, and is supported in \([-1,2]^2\).
Both \(G_{\mathrm{cut}}\) and \(D_\rho\) belong to
\(L^1(\mathbb R^2)\cap L^2(\mathbb R^2)\), since they are bounded and
compactly supported.

Fix the other coordinate and let \(\varphi\in C^1(\mathbb R;\mathbb R)\).
On each of the intervals \((-1,0)\), \((0,1)\), and \((1,2)\), ordinary
integration by parts applies to the corresponding slice of
\(G_{\mathrm{cut}}\), whose derivative there is the slice of \(D_\rho\).
Adding the three identities cancels the contributions at \(0\) and \(1\)
by continuity. The contributions at \(-1\) and \(2\) vanish because the
cutoff is zero. Extending the integrals to the real line therefore gives
\[
\int_{\mathbb R}
 \left.G_{\mathrm{cut}}(u_{\mathrm{loc}})
 \right|_{u_\rho^{\mathrm{loc}}=t}
 \varphi'(t)\,dt
=
-\int_{\mathbb R}
 \left.D_\rho(u_{\mathrm{loc}})
 \right|_{u_\rho^{\mathrm{loc}}=t}
 \varphi(t)\,dt.
\]
This holds for each \(\rho\in\{1,2\}\); hence \(D_\rho\) is the weak
coordinate derivative of \(G_{\mathrm{cut}}\).
% lean: CausalSmith.Experimentation.BivariateSobolevDesignCapacity.pairCosineExtensionPartial_weak_slice

\item \emph{Spatial energy on the nine support cells.}
For an integer \(\kappa\), integration of the cosine primitive gives
\[
\int_0^1\cos(\pi\kappa t)\,dt
=\mathbf1_{\{\kappa=0\}},
\qquad
\int_{-1}^2\cos(\pi\kappa t)\,dt
=3\mathbf1_{\{\kappa=0\}}.
\]
Indeed, for \(\kappa=0\) these are the interval lengths; for
\(\kappa\ne0\), all endpoint sine values vanish. For positive integer
frequencies \(a,b\), the identities
\[
\begin{aligned}
\cos(\pi a t)\cos(\pi b t)
&=\frac12\bigl(\cos(\pi(a-b)t)+\cos(\pi(a+b)t)\bigr),\\
\sin(\pi a t)\sin(\pi b t)
&=\frac12\bigl(\cos(\pi(a-b)t)-\cos(\pi(a+b)t)\bigr)
\end{aligned}
\]
therefore imply
\[
\begin{aligned}
\int_0^1\cos(\pi a t)\cos(\pi b t)\,dt
&=\int_0^1\sin(\pi a t)\sin(\pi b t)\,dt
 =\frac12\mathbf1_{\{a=b\}},\\
\int_{-1}^2\cos(\pi a t)\cos(\pi b t)\,dt
&=\int_{-1}^2\sin(\pi a t)\sin(\pi b t)\,dt
 =\frac32\mathbf1_{\{a=b\}}.
\end{aligned}
\]
Expanding the finite squares and using product integration eliminates
every off-diagonal frequency term. On the support square this yields
\[
\int_{[-1,2]^2}|U_{\mathrm{loc}}|^2
=4\left(\frac32\right)^2 E_{\mathrm{coef}}
=9E_{\mathrm{coef}}.
\]
Differentiating the polynomial gives
\[
\begin{aligned}
\partial_1U_{\mathrm{loc}}(u_{\mathrm{loc}})
&=-2\pi\sum_{a,b=1}^L
 a\,c^{\mathrm{loc}}_{ab}
 \sin(\pi a u_1^{\mathrm{loc}})
 \cos(\pi b u_2^{\mathrm{loc}}),\\
\partial_2U_{\mathrm{loc}}(u_{\mathrm{loc}})
&=-2\pi\sum_{a,b=1}^L
 b\,c^{\mathrm{loc}}_{ab}
 \cos(\pi a u_1^{\mathrm{loc}})
 \sin(\pi b u_2^{\mathrm{loc}}).
\end{aligned}
\]
The unit-square derivative energies satisfy
\[
\begin{aligned}
\int_{[0,1]^2}|\partial_1U_{\mathrm{loc}}|^2
&=\pi^2\sum_{a,b=1}^L a^2(c^{\mathrm{loc}}_{ab})^2
 \leq\pi^2L^2E_{\mathrm{coef}},\\
\int_{[0,1]^2}|\partial_2U_{\mathrm{loc}}|^2
&=\pi^2\sum_{a,b=1}^L b^2(c^{\mathrm{loc}}_{ab})^2
 \leq\pi^2L^2E_{\mathrm{coef}}.
\end{aligned}
\]
On \([-1,2]^2\), the corresponding energies are nine times these
frequency sums. In particular,
\[
\begin{aligned}
\int_{[-1,2]^2}|\partial_1U_{\mathrm{loc}}|^2
&=9\pi^2\sum_{a,b=1}^L a^2(c^{\mathrm{loc}}_{ab})^2
 \leq9\pi^2L^2E_{\mathrm{coef}},\\
\int_{[-1,2]^2}|\partial_2U_{\mathrm{loc}}|^2
&=9\pi^2\sum_{a,b=1}^L b^2(c^{\mathrm{loc}}_{ab})^2
 \leq9\pi^2L^2E_{\mathrm{coef}},\\
\int_{[-1,2]^2}\|\nabla U_{\mathrm{loc}}\|_2^2
&\leq18\pi^2L^2E_{\mathrm{coef}}.
\end{aligned}
\]
Since \(0\leq\chi_{\mathrm{cut}}\leq1\),
\[
\|G_{\mathrm{cut}}\|_{L^2(\mathbb R^2)}^2
\leq9E_{\mathrm{coef}}.
\]
For each coordinate, apply
\((r_1+r_2)^2\leq2r_1^2+2r_2^2\), valid for real \(r_1,r_2\), to the
two terms defining \(D_\rho\). Summing the two coordinate inequalities
and using the cutoff derivative bound gives
\[
\frac12(D_1^2+D_2^2)
\leq
\chi_{\mathrm{cut}}^2
 \bigl(|\partial_1U_{\mathrm{loc}}|^2+
       |\partial_2U_{\mathrm{loc}}|^2\bigr)
+2|U_{\mathrm{loc}}|^2.
\]
Thus
\[
D_1^2+D_2^2
\leq2\|\nabla U_{\mathrm{loc}}\|_2^2+4|U_{\mathrm{loc}}|^2
\]
on the support square. Integration proves
\[
\|D_1\|_{L^2(\mathbb R^2)}^2+
\|D_2\|_{L^2(\mathbb R^2)}^2
\leq(36+36\pi^2L^2)E_{\mathrm{coef}}.
\]
% lean: CausalSmith.Experimentation.BivariateSobolevDesignCapacity.pairCosineExtensionPartial_spatial_energy_le

\item \emph{Fourier endpoint bounds.}
For an integrable function
\(\widetilde G:\mathbb R^2\to\mathbb C\), use the unitary angular-frequency
Fourier transform
\[
(\mathcal F_2\widetilde G)(\omega_{\mathrm{loc}})
=\frac{1}{2\pi}\int_{\mathbb R^2}
 e^{-\mathrm i\langle\omega_{\mathrm{loc}},u_{\mathrm{loc}}\rangle}
 \widetilde G(u_{\mathrm{loc}})\,du_{\mathrm{loc}},
\qquad \omega_{\mathrm{loc}}\in\mathbb R^2.
\]
For \(\lambda\in[0,1]\) and
\(\widetilde G\in L^1(\mathbb R^2)\cap L^2(\mathbb R^2)\), define
\[
\mathcal E_\lambda(\widetilde G)
=\int_{\mathbb R^2}
 |\mathcal F_2\widetilde G(\omega_{\mathrm{loc}})|^2
 \left(1+\frac{\|\omega_{\mathrm{loc}}\|_2^2}{2}\right)^\lambda
 \,d\omega_{\mathrm{loc}}.
\]
The slice integration-by-parts identities, applied to the real and
imaginary parts of the Fourier exponential, give
\[
\mathcal F_2D_\rho(\omega_{\mathrm{loc}})
=\mathrm i\,\omega_\rho^{\mathrm{loc}}\,
 \mathcal F_2G_{\mathrm{cut}}(\omega_{\mathrm{loc}}),
\qquad \rho\in\{1,2\}.
\]
Plancherel's identity consequently yields
\[
\begin{aligned}
\mathcal E_0(G_{\mathrm{cut}})
&=\|G_{\mathrm{cut}}\|_2^2
 \leq9E_{\mathrm{coef}}
 \leq16E_{\mathrm{coef}},\\
\mathcal E_1(G_{\mathrm{cut}})
&=\|G_{\mathrm{cut}}\|_2^2
 +\frac12\bigl(\|D_1\|_2^2+\|D_2\|_2^2\bigr)\\
&\leq(27+18\pi^2L^2)E_{\mathrm{coef}}.
\end{aligned}
\]
The calibration uses \(L^2\geq1\):
\[
27+18\pi^2L^2
\leq(48+32\pi^2)L^2=C_0^2L^2.
\]
Hence the required endpoint estimates are
\[
\mathcal E_0(G_{\mathrm{cut}})\leq16E_{\mathrm{coef}},
\qquad
\mathcal E_1(G_{\mathrm{cut}})\leq C_0^2L^2E_{\mathrm{coef}}.
\]
% lean: CausalSmith.Experimentation.BivariateSobolevDesignCapacity.pairCosineExtension_fourier_endpoints

\item \emph{Fractional energy and the restriction norm.}
Define the nonnegative energy densities on \(\mathbb R^2\) by
\[
\begin{aligned}
\eta_0(\omega_{\mathrm{loc}})
&=|\mathcal F_2G_{\mathrm{cut}}(\omega_{\mathrm{loc}})|^2,\\
\eta_1(\omega_{\mathrm{loc}})
&=\eta_0(\omega_{\mathrm{loc}})
 \left(1+\frac{\|\omega_{\mathrm{loc}}\|_2^2}{2}\right).
\end{aligned}
\]
With the convention \(r^0=1\) for \(r\geq0\), their pointwise factorization is
\[
\eta_0^{\,1-s}\eta_1^{\,s}
=\eta_0
 \left(1+\frac{\|\omega_{\mathrm{loc}}\|_2^2}{2}\right)^s.
\]
The integral form of Hölder's inequality with nonnegative weights
\(1-s\) and \(s\), whose sum is one, therefore gives
\[
\mathcal E_s(G_{\mathrm{cut}})
\leq
\mathcal E_0(G_{\mathrm{cut}})^{1-s}
\mathcal E_1(G_{\mathrm{cut}})^s.
\]
This form includes endpoint weights and zero energy. Since
\(16\leq C_0^2\), the endpoint bounds imply
\[
\begin{aligned}
\mathcal E_s(G_{\mathrm{cut}})
&\leq
 (C_0^2E_{\mathrm{coef}})^{1-s}
 (C_0^2E_{\mathrm{coef}}L^2)^s\\
&=C_0^2L^{2s}E_{\mathrm{coef}}.
\end{aligned}
\]
In particular, when \(E_{\mathrm{coef}}=0\), every coefficient and the
extension vanish, and this bound is zero.

The squared restriction norm is the extension infimum
\[
N_{2,s}(U_{\mathrm{loc}}|_{[0,1]^2})^2
=
\inf_{\substack{
 \widetilde G\in L^1(\mathbb R^2;\mathbb C)\cap L^2(\mathbb R^2;\mathbb C)\\
 \widetilde G=U_{\mathrm{loc}}\ \text{a.e.\ on }[0,1]^2}}
 \mathcal E_s(\widetilde G).
\]
The continuous compactly supported extension \(G_{\mathrm{cut}}\) is
admissible in this infimum. We obtain
\[
N_{2,s}(U_{\mathrm{loc}}|_{[0,1]^2})^2
\leq
\mathcal E_s(G_{\mathrm{cut}})
\leq C_0^2L^{2s}
 \sum_{a,b=1}^L(c^{\mathrm{loc}}_{ab})^2.
\]
% lean: CausalSmith.Experimentation.BivariateSobolevDesignCapacity.pairCosinePolynomial_restriction_bound

\item \emph{Component budget estimates.}
For each \(j<\ell\), define the normalized pair coefficients and their
energy by
\[
c^{\mathrm{prior}}_{j\ell ab}
=\gamma_{\mathrm{prior}}\xi_{j\ell ab},
\qquad
E^{\mathrm{prior}}_{j\ell}
=\sum_{a,b=1}^L(c^{\mathrm{prior}}_{j\ell ab})^2,
\qquad 1\leq a,b\leq L.
\]
Every sign has square one, so
\[
E^{\mathrm{prior}}_{j\ell}
=\sum_{a,b=1}^L\gamma_{\mathrm{prior}}^2
=L^2\gamma_{\mathrm{prior}}^2.
\]
The canonical pair formula and the restriction bound now give
\[
N_{2,s}(g_{j\ell}(m_\xi))^2
\leq C_0^2L^{2s}E^{\mathrm{prior}}_{j\ell}
=C_0^2L^{2s}L^2\gamma_{\mathrm{prior}}^2.
\]
The main effects are zero; the zero extension certifies their zero
restriction norms. Hence
\[
\sum_{j=1}^dN_{1,s}(g_j(m_\xi))^2=0.
\]
Summing the common pair bound over the \(B\) coordinate pairs yields
\[
\sum_{j=1}^dN_{1,s}(g_j(m_\xi))^2
+\sum_{j<\ell}N_{2,s}(g_{j\ell}(m_\xi))^2
\leq B\,C_0^2L^{2s}L^2\gamma_{\mathrm{prior}}^2.
\]
% lean: CausalSmith.Experimentation.BivariateSobolevDesignCapacity.mXi_pooledBudget

\item \emph{Cube inner products.}
Take
\[
\alpha=(j,\ell,a,b)\in\mathcal I,
\qquad
\beta=(j',\ell',a',b')\in\mathcal I.
\]
Product integration factors the feature inner product as
\[
\begin{aligned}
\int_{[0,1]^d}V_\alpha(x)V_\beta(x)\,dx
=4\prod_{\iota=1}^d\int_0^1
&\cos(\pi a t)^{\mathbf1_{\{\iota=j\}}}
 \cos(\pi b t)^{\mathbf1_{\{\iota=\ell\}}}\\
&{}\cdot
 \cos(\pi a't)^{\mathbf1_{\{\iota=j'\}}}
 \cos(\pi b't)^{\mathbf1_{\{\iota=\ell'\}}}\,dt .
\end{aligned}
\]
First suppose \(j=j'\). If also \(\ell=\ell'\), the unit-interval cosine
integrals give
\[
\int_{[0,1]^d}V_\alpha V_\beta
=4\left(\frac12\mathbf1_{\{a=a'\}}\right)
  \left(\frac12\mathbf1_{\{b=b'\}}\right)
=\mathbf1_{\{\alpha=\beta\}}.
\]
If instead \(\ell\ne\ell'\), then \(\ell\) is distinct from both
\(j'=j\) and \(\ell'\). Its factor occurs exactly once and has integral
zero.

Now suppose \(j\ne j'\). If \(j=\ell'\), the ordering gives
\[
j'<\ell'=j<\ell,
\]
so \(\ell\) is distinct from both coordinates of \(\beta\) and occurs
exactly once. If \(j\ne\ell'\), coordinate \(j\) is distinct from both
coordinates of \(\beta\) and occurs exactly once. In each of these
unmatched-coordinate cases, the product factorization contains a
positive-frequency cosine mean and consequently gives
\[
\int_{[0,1]^d}V_\alpha(x)V_\beta(x)\,dx=0.
\]
Combining the cases proves
\[
\int_{[0,1]^d}V_\alpha(x)V_\beta(x)\,dx
=\mathbf1_{\{\alpha=\beta\}},
\qquad \alpha,\beta\in\mathcal I.
\]
% lean: CausalSmith.Experimentation.BivariateSobolevDesignCapacity.pairFeature_cube_inner_product

\item \emph{Normalization and class membership.}
Expanding the finite square and applying the cube inner-product identity
gives
\[
\begin{aligned}
\int_{[0,1]^d}m_\xi(x)^2\,dx
&=\gamma_{\mathrm{prior}}^2
 \sum_{\alpha,\beta\in\mathcal I}
 \xi_\alpha\xi_\beta\mathbf1_{\{\alpha=\beta\}}\\
&=\gamma_{\mathrm{prior}}^2
 \sum_{\alpha\in\mathcal I}\xi_\alpha^2
=\gamma_{\mathrm{prior}}^2M
=\frac{1}{C_0^2L^{2s}}.
\end{aligned}
\]
Thus the normalized coefficient energy satisfies
\[
C_0^2L^{2s}
 \sum_{\alpha\in\mathcal I}
  (\gamma_{\mathrm{prior}}\xi_\alpha)^2
=C_0^2L^{2s}\gamma_{\mathrm{prior}}^2M
=1.
\]
Since \(M=BL^2\), this also calibrates the common pair budget:
\[
B\,C_0^2L^{2s}L^2\gamma_{\mathrm{prior}}^2
=C_0^2L^{2s}\gamma_{\mathrm{prior}}^2M
=1.
\]
The component estimates therefore establish
\[
\sum_{j=1}^dN_{1,s}(g_j(m_\xi))^2
+\sum_{j<\ell}N_{2,s}(g_{j\ell}(m_\xi))^2
\leq1,
\]
as required in \cref{obj:ass:pooled-budget}. Together with measurability,
square integrability, centering, and the canonical order-two
reconstruction, this proves
\[
m_\xi\in\mathcal H_{d,s}
\]
by \cref{obj:def:sobolev-class}.
% lean: CausalSmith.Experimentation.BivariateSobolevDesignCapacity.mXi_pooledBudget

\item \emph{Transfer to the first sampled covariate.}
Let \(n\geq1\) and let \(X\) satisfy \cref{obj:ass:uniform-draw}.
Taking the first coordinate marginal of its product sampling law gives
\[
\operatorname{Law}(X_1)=\operatorname{Unif}([0,1]^d).
\]
The feature products are measurable and bounded, so integration against
this marginal yields
\[
\mathbb E[V_\alpha(X_1)V_\beta(X_1)]
=\int_{[0,1]^d}V_\alpha(x)V_\beta(x)\,dx
=\mathbf1_{\{\alpha=\beta\}},
\qquad \alpha,\beta\in\mathcal I.
\]
There are \(M\) feature coordinates. Taking these expectations entrywise
therefore proves
\[
\mathbb E[V(X_1)V(X_1)^{\mathsf T}]=I_M.
\]
The unmatched-coordinate cases include distinct pairs with one shared
original coordinate, so the identity covers those pairs as well.
% lean: CausalSmith.Experimentation.BivariateSobolevDesignCapacity.uniformDraw_pairFeature_inner_product
\end{enumerate}
\end{proof}

Class membership holds for every coefficient array, making the entire finite prior admissible under the pooled budget in \cref{obj:def:sobolev-class}. The isotropy conclusion treats all coordinate pairs together, including overlapping pairs. It therefore supplies the full dimension \(M\) for the signing comparison.

For that comparison, introduce an appendix-local random vector \(v\in\mathbb R^{r_0}\), where \(r_0\) is a positive integer, with identity second-moment matrix \(I_{r_0}\). Let \(v_1,\ldots,v_n\) be independent copies, and let \(z\in\{-1,1\}^n\) denote a candidate assignment sign vector. The expected minimum over these signs is controlled by the following bound.

\begin{lemmav}{lem:isotropic-row-signing}[Isotropic row signing bound]*
Let \(r_0\) be a positive integer and \(n\) a nonnegative integer. Let \(v_1,\ldots,v_n\) be independent copies of a random vector \(v\in\mathbb R^{r_0}\), under the usual measurability and square-integrability conditions, with
\[
\mathbb E[vv^{\mathsf T}]=I_{r_0},
\]
where \(I_{r_0}\) is the identity matrix. Then
\[
\mathbb E\min_{z\in\{-1,1\}^n}
\left\|\sum_{i=1}^n z_i v_i\right\|_2^2
\geq nr_0-n^2-16n\sqrt{nr_0}.
\]
In particular, if \(r_0\geq4096n\), then
\[
\mathbb E\min_{z\in\{-1,1\}^n}
\left\|\sum_{i=1}^n z_i v_i\right\|_2^2
\geq \frac{nr_0}{2}.
\]
\end{lemmav}
% CAUSALSMITH-CITED-SCOPE-BEGIN lem:isotropic-row-signing
\verificationfootnotetext{\textbf{Formalization scope.} This result uses the cited conclusion \citep{BartlettMendelson2002Complexities} (Theorem 12(4), with Rademacher complexity as in Definition 2). The cited source proof is not formalized here: Lean verifies its use and all remaining steps. If that cited conclusion is false, the portions of this result that depend on it are not certified.}
% CAUSALSMITH-CITED-SCOPE-END lem:isotropic-row-signing

\begin{proof}[Proof of \cref{obj:lem:isotropic-row-signing}]
For a deterministic row tuple \(\boldsymbol v=(v_1,\ldots,v_n)\in(\mathbb R^{r_0})^n\), define
\[
\begin{aligned}
\mathsf D_{\rm row}(\boldsymbol v)
&:=\sum_{i=1}^n\|v_i\|_2^2,\\
\mathsf O_{\rm row}(\boldsymbol v,z)
&:=\sum_{\substack{1\leq i,i'\leq n\\i\ne i'}}
z_i z_{i'}\langle v_i,v_{i'}\rangle,
\qquad z\in\{-1,1\}^n,\\
\mathsf M_{\rm off}(\boldsymbol v)
&:=\max_{z\in\{-1,1\}^n}
|\mathsf O_{\rm row}(\boldsymbol v,z)|,\\
\mathsf M_{\rm sign}(\boldsymbol v)
&:=\min_{z\in\{-1,1\}^n}
\left\|\sum_{i=1}^n z_i v_i\right\|_2^2.
\end{aligned}
\]
For \(n=0\), the sign space consists of the empty tuple, and all four quantities are zero. Square-integrability ensures integrability of the finite extrema used below.

\begin{enumerate}
\item Suppose first that \(n\leq256\).
Expanding the squared norm gives, for every signing,
\[
\left\|\sum_{i=1}^n z_i v_i\right\|_2^2
=\mathsf D_{\rm row}(\boldsymbol v)
+\mathsf O_{\rm row}(\boldsymbol v,z).
\]
Consequently,
\[
\mathsf M_{\rm sign}(\boldsymbol v)
\geq\mathsf D_{\rm row}(\boldsymbol v)
-\mathsf M_{\rm off}(\boldsymbol v).
\]
Isotropy gives
\[
\mathbb E\|v\|_2^2
=\sum_{\alpha=1}^{r_0}\mathbb E(v_\alpha^2)=r_0,
\qquad
\mathbb E\mathsf D_{\rm row}(\boldsymbol v)=nr_0,
\]
and hence
\[
\mathbb E\mathsf M_{\rm sign}(\boldsymbol v)
\geq nr_0-\mathbb E\mathsf M_{\rm off}(\boldsymbol v).
\]
These identities and the diagonal-subtraction inequality hold for every \(n\).
% lean: CausalSmith.Experimentation.BivariateSobolevDesignCapacity.isotropic_signing_lower_of_offDiagonal

For distinct rows, independence and isotropy yield
\[
\begin{aligned}
\mathbb E\langle v_i,v_{i'}\rangle^2
&=\sum_{\alpha,\beta=1}^{r_0}
\mathbb E[(v_i)_\alpha(v_i)_\beta]\,
\mathbb E[(v_{i'})_\alpha(v_{i'})_\beta]\\
&=\sum_{\alpha,\beta=1}^{r_0}\mathbf 1_{\{\alpha=\beta\}}
=r_0.
\end{aligned}
\]
Thus Cauchy--Schwarz implies
\[
\mathbb E|\langle v_i,v_{i'}\rangle|
\leq\sqrt{r_0}.
\]
The triangle inequality, followed by the loose count of at most \(n^2\) ordered pairs, gives
\[
\mathsf M_{\rm off}(\boldsymbol v)
\leq\sum_{\substack{1\leq i,i'\leq n\\i\ne i'}}
|\langle v_i,v_{i'}\rangle|,
\qquad
\mathbb E\mathsf M_{\rm off}(\boldsymbol v)
\leq n^2\sqrt{r_0}.
\]
Therefore,
\[
\mathbb E\mathsf M_{\rm sign}(\boldsymbol v)
\geq nr_0-n^2\sqrt{r_0}.
\]
% lean: CausalSmith.Experimentation.BivariateSobolevDesignCapacity.isotropic_signing_lower_pair_count

Since \(n\leq256\),
\[
\sqrt n\leq16,\qquad
n=(\sqrt n)^2\leq16\sqrt n,
\qquad
n^2\sqrt{r_0}\leq16n\sqrt{nr_0}.
\]
It follows that
\[
\mathbb E\mathsf M_{\rm sign}(\boldsymbol v)
\geq nr_0-16n\sqrt{nr_0}
\geq nr_0-n^2-16n\sqrt{nr_0}.
\]
The same calculation covers \(n=0\).
% lean: CausalSmith.Experimentation.BivariateSobolevDesignCapacity.isotropic_small_sample_signing_lower

\item Suppose now that \(n>256\). We first establish the finite-class estimate needed for the partition argument.
Let
\[
n_{\rm loc}\in\mathbb N_0,\qquad
\boldsymbol v_{\rm loc}
=(v^{\rm loc}_1,\ldots,v^{\rm loc}_{n_{\rm loc}})
\in(\mathbb R^{r_0})^{n_{\rm loc}},
\]
where the local rows are independent copies of \(v\). Let
\[
\varnothing\ne\mathcal U_{\rm fin}\subset\mathbb R^{r_0}
\quad\text{be finite},\qquad
R_{\rm dir}:=\max_{u_{\rm dir}\in\mathcal U_{\rm fin}}
\|u_{\rm dir}\|_2.
\]
For \(u_{\rm dir}\in\mathcal U_{\rm fin}\), define
\[
\mathsf f_{u_{\rm dir}}:\mathbb R^{r_0}\to\mathbb R,
\qquad
\mathsf f_{u_{\rm dir}}(v):=|\langle u_{\rm dir},v\rangle|,
\]
and
\[
\mathsf C_{\rm fin}(\boldsymbol v_{\rm loc})
:=
\max_{u_{\rm dir}\in\mathcal U_{\rm fin}}
\left|
\sum_{i=1}^{n_{\rm loc}}
\bigl(\mathsf f_{u_{\rm dir}}(v_i^{\rm loc})
-\mathbb E\mathsf f_{u_{\rm dir}}(v)\bigr)
\right|.
\]
Isotropy and Cauchy--Schwarz give
\[
\mathbb E\langle u_{\rm dir},v\rangle^2
=u_{\rm dir}^{\mathsf T}\mathbb E[vv^{\mathsf T}]u_{\rm dir}
=\|u_{\rm dir}\|_2^2,
\qquad
\mathbb E\mathsf f_{u_{\rm dir}}(v)
\leq\|u_{\rm dir}\|_2.
\]
Adding and subtracting the means therefore yields the pointwise bound
\[
\max_{u_{\rm dir}\in\mathcal U_{\rm fin}}
\sum_{i=1}^{n_{\rm loc}}
\mathsf f_{u_{\rm dir}}(v_i^{\rm loc})
\leq n_{\rm loc}R_{\rm dir}
+\mathsf C_{\rm fin}(\boldsymbol v_{\rm loc}).
\]

Introduce an independent copy
\[
\boldsymbol v_{\rm ghost}
=(v^{\rm ghost}_1,\ldots,v^{\rm ghost}_{n_{\rm loc}})
\]
of the local sample, and independent fair signs
\[
\rho=(\rho_1,\ldots,\rho_{n_{\rm loc}})
\in\{-1,1\}^{n_{\rm loc}},
\qquad
\mathbb P(\rho_i=1)=\mathbb P(\rho_i=-1)=\tfrac12.
\]
Define
\[
\mathsf S_\rho(\boldsymbol v_{\rm loc})
:=\sum_{i=1}^{n_{\rm loc}}\rho_i v_i^{\rm loc}.
\]
Jensen's inequality, exchanging the two entries of each independent sample pair according to \(\rho_i\), and the triangle inequality give
\[
\begin{aligned}
\mathbb E\mathsf C_{\rm fin}(\boldsymbol v_{\rm loc})
&\leq
\mathbb E\max_{u_{\rm dir}\in\mathcal U_{\rm fin}}
\left|\sum_{i=1}^{n_{\rm loc}}
\bigl(\mathsf f_{u_{\rm dir}}(v_i^{\rm loc})
-\mathsf f_{u_{\rm dir}}(v_i^{\rm ghost})\bigr)\right|\\
&=
\mathbb E\max_{u_{\rm dir}\in\mathcal U_{\rm fin}}
\left|\sum_{i=1}^{n_{\rm loc}}\rho_i
\bigl(\mathsf f_{u_{\rm dir}}(v_i^{\rm loc})
-\mathsf f_{u_{\rm dir}}(v_i^{\rm ghost})\bigr)\right|\\
&\leq
2\mathbb E\max_{u_{\rm dir}\in\mathcal U_{\rm fin}}
\left|\sum_{i=1}^{n_{\rm loc}}\rho_i
\mathsf f_{u_{\rm dir}}(v_i^{\rm loc})\right|.
\end{aligned}
\]
% lean: CausalSmith.Experimentation.BivariateSobolevDesignCapacity.finite_first_moment_symmetrization

Apply the contraction inequality of
\citep[Theorem~12(4), with Rademacher complexity as in Definition~2]{BartlettMendelson2002Complexities}
to the finite class of linear projections and the map
\[
\varphi_{\rm abs}:\mathbb R\to\mathbb R,
\qquad
\varphi_{\rm abs}(t_{\rm abs}):=|t_{\rm abs}|.
\]
This map is \(1\)-Lipschitz and vanishes at zero. The class is its own countable dense subclass, and a common integrable envelope is
\[
\mathsf E_{\rm fin}(v)
:=\sum_{u_{\rm dir}\in\mathcal U_{\rm fin}}
|\langle u_{\rm dir},v\rangle|
\leq|\mathcal U_{\rm fin}|R_{\rm dir}\|v\|_2,
\qquad
\mathbb E\|v\|_2\leq\sqrt{r_0}.
\]
The envelope bound holds at every \(v\). Contraction supplies
\[
\mathbb E\max_{u_{\rm dir}\in\mathcal U_{\rm fin}}
\left|\sum_{i=1}^{n_{\rm loc}}\rho_i
|\langle u_{\rm dir},v_i^{\rm loc}\rangle|\right|
\leq
2\mathbb E\max_{u_{\rm dir}\in\mathcal U_{\rm fin}}
|\langle \mathsf S_\rho(\boldsymbol v_{\rm loc}),u_{\rm dir}\rangle|.
\]
Combining the preceding bounds proves
\[
\mathbb E\max_{u_{\rm dir}\in\mathcal U_{\rm fin}}
\sum_{i=1}^{n_{\rm loc}}|\langle u_{\rm dir},v_i^{\rm loc}\rangle|
\leq n_{\rm loc}R_{\rm dir}
+4\mathbb E\max_{u_{\rm dir}\in\mathcal U_{\rm fin}}
|\langle \mathsf S_\rho(\boldsymbol v_{\rm loc}),u_{\rm dir}\rangle|.
\]
All sums vanish when \(n_{\rm loc}=0\).
% lean: CausalSmith.Experimentation.BivariateSobolevDesignCapacity.finite_projection_sum_integral_le_exact_rad

\item Consider two independent groups of rows,
\[
\begin{gathered}
n_{\rm left},n_{\rm right}\in\mathbb N_0,\\
\boldsymbol v_{\rm left}
=(v^{\rm left}_1,\ldots,v^{\rm left}_{n_{\rm left}}),
\qquad
\boldsymbol v_{\rm right}
=(v^{\rm right}_1,\ldots,v^{\rm right}_{n_{\rm right}}),
\end{gathered}
\]
each consisting of independent copies of \(v\). For any local row tuple, define
\[
\mathsf K_{\rm row}(\boldsymbol v_{\rm loc})
:=
\max_{z_{\rm loc}\in\{-1,1\}^{n_{\rm loc}}}
\left\|\sum_{i=1}^{n_{\rm loc}}
(z_{\rm loc})_i v_i^{\rm loc}\right\|_2.
\]
Also define the finite direction set and cross-group quantity
\[
\begin{aligned}
\mathcal U_{\rm right}(\boldsymbol v_{\rm right})
&:=
\left\{\sum_{i'=1}^{n_{\rm right}}
(z_{\rm right})_{i'}v_{i'}^{\rm right}:
z_{\rm right}\in\{-1,1\}^{n_{\rm right}}\right\},\\
\mathsf H_{\rm cross}(\boldsymbol v_{\rm left},\boldsymbol v_{\rm right})
&:=
\max_{u_{\rm dir}\in\mathcal U_{\rm right}(\boldsymbol v_{\rm right})}
\sum_{i=1}^{n_{\rm left}}
|\langle u_{\rm dir},v_i^{\rm left}\rangle|.
\end{aligned}
\]
For an empty tuple, \(\mathsf K_{\rm row}=0\); if the right group is empty, \(\mathcal U_{\rm right}=\{0\}\). Thus \(\mathsf H_{\rm cross}=0\) whenever either group is empty.

Conditional on the right group, the finite-class estimate applies with maximal direction norm \(\mathsf K_{\rm row}(\boldsymbol v_{\rm right})\). Choosing each right-group sign to agree with its projection gives
\[
\max_{u_{\rm dir}\in\mathcal U_{\rm right}(\boldsymbol v_{\rm right})}
|\langle\mathsf S_\rho(\boldsymbol v_{\rm left}),u_{\rm dir}\rangle|
=
\sum_{i'=1}^{n_{\rm right}}
|\langle\mathsf S_\rho(\boldsymbol v_{\rm left}),v_{i'}^{\rm right}\rangle|.
\]
Consequently,
\[
\begin{aligned}
\mathbb E_{\rm left}\,
\mathsf H_{\rm cross}(\boldsymbol v_{\rm left},\boldsymbol v_{\rm right})
\leq{}&
n_{\rm left}\mathsf K_{\rm row}(\boldsymbol v_{\rm right})\\
&+4\mathbb E_{{\rm left},\rho}
\sum_{i'=1}^{n_{\rm right}}
|\langle\mathsf S_\rho(\boldsymbol v_{\rm left}),v_{i'}^{\rm right}\rangle|.
\end{aligned}
\]
% lean: CausalSmith.Experimentation.BivariateSobolevDesignCapacity.fixed_group_cross_integral_le_exact_rad

For any local sample, averaging over the fair signs cancels the off-diagonal terms:
\[
\mathbb E_\rho\|\mathsf S_\rho(\boldsymbol v_{\rm loc})\|_2^2
=\sum_{i=1}^{n_{\rm loc}}\|v_i^{\rm loc}\|_2^2.
\]
Thus
\[
\mathbb E\|\mathsf S_\rho(\boldsymbol v_{\rm loc})\|_2^2
=n_{\rm loc}r_0,
\qquad
\mathbb E\|\mathsf S_\rho(\boldsymbol v_{\rm loc})\|_2
\leq\sqrt{n_{\rm loc}r_0}.
\]
For every fixed \(u_{\rm dir}\in\mathbb R^{r_0}\), the projection estimate gives
\[
\mathbb E_{\rm right}
\sum_{i'=1}^{n_{\rm right}}
|\langle u_{\rm dir},v_{i'}^{\rm right}\rangle|
\leq n_{\rm right}\|u_{\rm dir}\|_2.
\]
Independence allows this estimate to be applied conditional on the left sample and its signs. Interchanging the integrable expectations yields
\[
\mathbb E
\sum_{i'=1}^{n_{\rm right}}
|\langle\mathsf S_\rho(\boldsymbol v_{\rm left}),v_{i'}^{\rm right}\rangle|
\leq
n_{\rm right}\sqrt{n_{\rm left}r_0}.
\]
After integrating the conditional cross-group bound, we obtain
\[
\mathbb E\mathsf H_{\rm cross}(\boldsymbol v_{\rm left},\boldsymbol v_{\rm right})
\leq
n_{\rm left}\mathbb E\mathsf K_{\rm row}(\boldsymbol v_{\rm right})
+4n_{\rm right}\sqrt{n_{\rm left}r_0}.
\]
% lean: CausalSmith.Experimentation.BivariateSobolevDesignCapacity.independent_group_cross_integral_le

\item We next bound the expected maximum signed norm.
Define the Euclidean unit ball
\[
\mathcal B_{\rm dir}
:=\{u_{\rm dir}\in\mathbb R^{r_0}:\|u_{\rm dir}\|_2\leq1\}.
\]
Norm duality and choosing the signs of the projections imply
\[
\mathsf K_{\rm row}(\boldsymbol v_{\rm loc})
=
\sup_{u_{\rm dir}\in\mathcal B_{\rm dir}}
\sum_{i=1}^{n_{\rm loc}}|\langle u_{\rm dir},v_i^{\rm loc}\rangle|.
\]
Indeed, for a fixed direction, the maximum absolute projection of a signed sum is the sum of the absolute projections; taking the supremum over unit directions gives its Euclidean norm.
% lean: CausalSmith.Experimentation.BivariateSobolevDesignCapacity.signingNormMax_eq_unit_projection_sum

For every nonempty finite \(\mathcal U_{\rm fin}\subset\mathcal B_{\rm dir}\),
\[
\max_{u_{\rm dir}\in\mathcal U_{\rm fin}}
|\langle\mathsf S_\rho(\boldsymbol v_{\rm loc}),u_{\rm dir}\rangle|
\leq\|\mathsf S_\rho(\boldsymbol v_{\rm loc})\|_2.
\]
The finite-class bound and the signed-norm first moment therefore give
\[
\mathbb E\max_{u_{\rm dir}\in\mathcal U_{\rm fin}}
\sum_{i=1}^{n_{\rm loc}}|\langle u_{\rm dir},v_i^{\rm loc}\rangle|
\leq n_{\rm loc}+4\sqrt{n_{\rm loc}r_0}.
\]
% lean: CausalSmith.Experimentation.BivariateSobolevDesignCapacity.finite_unit_projection_sum_integral_le

To pass to the whole unit ball, put
\[
\mathsf A_{\rm norm}
:=\mathbb E\sum_{i=1}^{n_{\rm loc}}\|v_i^{\rm loc}\|_2,
\qquad
0\leq\mathsf A_{\rm norm}<\infty.
\]
For an arbitrary \(\delta_{\rm net}>0\), set
\[
\varepsilon_{\rm net}
:=\frac{\delta_{\rm net}}{\mathsf A_{\rm norm}+1}>0.
\]
Compactness provides a nonempty finite set
\[
\mathcal U_{\rm net}\subset\mathcal B_{\rm dir},
\qquad
\forall u_{\rm dir}\in\mathcal B_{\rm dir}\ 
\exists u_{\rm net}\in\mathcal U_{\rm net}:
\|u_{\rm dir}-u_{\rm net}\|_2\leq\varepsilon_{\rm net}.
\]
For such a pair of directions,
\[
|\langle u_{\rm dir},v_i^{\rm loc}\rangle|
\leq
|\langle u_{\rm net},v_i^{\rm loc}\rangle|
+\varepsilon_{\rm net}\|v_i^{\rm loc}\|_2.
\]
Summing and taking the supremum gives
\[
\mathsf K_{\rm row}(\boldsymbol v_{\rm loc})
\leq
\max_{u_{\rm net}\in\mathcal U_{\rm net}}
\sum_{i=1}^{n_{\rm loc}}|\langle u_{\rm net},v_i^{\rm loc}\rangle|
+\varepsilon_{\rm net}\sum_{i=1}^{n_{\rm loc}}\|v_i^{\rm loc}\|_2.
\]
Hence
\[
\begin{aligned}
\mathbb E\mathsf K_{\rm row}(\boldsymbol v_{\rm loc})
&\leq n_{\rm loc}+4\sqrt{n_{\rm loc}r_0}
+\varepsilon_{\rm net}\mathsf A_{\rm norm}\\
&=n_{\rm loc}+4\sqrt{n_{\rm loc}r_0}
+\delta_{\rm net}\frac{\mathsf A_{\rm norm}}{\mathsf A_{\rm norm}+1}\\
&\leq n_{\rm loc}+4\sqrt{n_{\rm loc}r_0}+\delta_{\rm net}.
\end{aligned}
\]
Letting \(\delta_{\rm net}\) tend to zero proves
\[
\mathbb E\mathsf K_{\rm row}(\boldsymbol v_{\rm loc})
\leq n_{\rm loc}+4\sqrt{n_{\rm loc}r_0}.
\]
This also covers the empty sample.
% lean: CausalSmith.Experimentation.BivariateSobolevDesignCapacity.isotropic_signingNormMax_integral_le

Substitution into the preceding cross-group estimate now yields
\[
\begin{aligned}
\mathbb E\mathsf H_{\rm cross}(\boldsymbol v_{\rm left},\boldsymbol v_{\rm right})
\leq{}&
n_{\rm left}n_{\rm right}
+4n_{\rm left}\sqrt{n_{\rm right}r_0}\\
&+4n_{\rm right}\sqrt{n_{\rm left}r_0}.
\end{aligned}
\]
% lean: CausalSmith.Experimentation.BivariateSobolevDesignCapacity.independent_group_cross_integral_le_explicit

\item Introduce fair selectors, independent of the rows,
\[
\eta=(\eta_1,\ldots,\eta_n)\in\{0,1\}^n,
\qquad
\mathbb P(\eta_i=1)=\mathbb P(\eta_i=0)=\tfrac12,
\]
with independent coordinates. Define
\[
\begin{aligned}
\mathcal I_\eta&:=\{i\in\{1,\ldots,n\}:\eta_i=1\},&
\mathcal J_\eta&:=\{i\in\{1,\ldots,n\}:\eta_i=0\},\\
\mathsf N_{\rm sel}(\eta)&:=|\mathcal I_\eta|,&
\mathsf N_{\rm comp}(\eta)&:=|\mathcal J_\eta|,\\
\boldsymbol v_{\rm sel}&:=(v_i)_{i\in\mathcal I_\eta},&
\boldsymbol v_{\rm comp}&:=(v_i)_{i\in\mathcal J_\eta},
\end{aligned}
\]
where the subsequences are listed in increasing index order. For \(z\in\{-1,1\}^n\), define
\[
\begin{aligned}
\mathsf C_{\rm cross}(\boldsymbol v,z,\eta)
&:=\sum_{i\in\mathcal I_\eta}\sum_{i'\in\mathcal J_\eta}
z_i z_{i'}\langle v_i,v_{i'}\rangle\\
&=\left\langle
\sum_{i\in\mathcal I_\eta}z_i v_i,\,
\sum_{i'\in\mathcal J_\eta}z_{i'}v_{i'}
\right\rangle,\\
\mathsf M_{\rm cross}(\boldsymbol v,\eta)
&:=\max_{z\in\{-1,1\}^n}
|\mathsf C_{\rm cross}(\boldsymbol v,z,\eta)|.
\end{aligned}
\]
For fixed rows and selectors, the triangle inequality gives
\[
\begin{aligned}
|\mathsf C_{\rm cross}(\boldsymbol v,z,\eta)|
&\leq
\sum_{i\in\mathcal I_\eta}
\left|\left\langle
v_i,\sum_{i'\in\mathcal J_\eta}z_{i'}v_{i'}
\right\rangle\right|\\
&\leq\mathsf H_{\rm cross}(\boldsymbol v_{\rm sel},\boldsymbol v_{\rm comp}).
\end{aligned}
\]
For each fixed \(\eta\), the two subsequences are independent samples from the row law. The cross-group estimate thus implies
\[
\begin{aligned}
\mathbb E_{\boldsymbol v}\mathsf M_{\rm cross}(\boldsymbol v,\eta)
\leq{}&
\mathsf N_{\rm sel}(\eta)\mathsf N_{\rm comp}(\eta)\\
&+4\mathsf N_{\rm sel}(\eta)
\sqrt{\mathsf N_{\rm comp}(\eta)r_0}\\
&+4\mathsf N_{\rm comp}(\eta)
\sqrt{\mathsf N_{\rm sel}(\eta)r_0}.
\end{aligned}
\]
Since
\[
\mathsf N_{\rm sel}(\eta)+\mathsf N_{\rm comp}(\eta)=n,
\]
the square of their difference gives
\[
\mathsf N_{\rm sel}(\eta)\mathsf N_{\rm comp}(\eta)
\leq\frac{n^2}{4}.
\]
Both group sizes are at most \(n\), so
\[
\begin{aligned}
&4\mathsf N_{\rm sel}(\eta)
\sqrt{\mathsf N_{\rm comp}(\eta)r_0}
+4\mathsf N_{\rm comp}(\eta)
\sqrt{\mathsf N_{\rm sel}(\eta)r_0}\\
&\hspace{2em}\leq
4\bigl(\mathsf N_{\rm sel}(\eta)+\mathsf N_{\rm comp}(\eta)\bigr)
\sqrt{nr_0}
=4n\sqrt{nr_0}.
\end{aligned}
\]
Therefore every selector satisfies
\[
\mathbb E_{\boldsymbol v}\mathsf M_{\rm cross}(\boldsymbol v,\eta)
\leq\frac{n^2}{4}+4n\sqrt{nr_0}.
\]
% lean: CausalSmith.Experimentation.BivariateSobolevDesignCapacity.signingPartitionCrossMax_integral_le

Averaging this bound over the finite selector space and interchanging the integrable expectations gives
\[
\begin{aligned}
4\mathbb E_{\boldsymbol v}\mathbb E_\eta
\mathsf M_{\rm cross}(\boldsymbol v,\eta)
&=4\mathbb E_\eta\mathbb E_{\boldsymbol v}
\mathsf M_{\rm cross}(\boldsymbol v,\eta)\\
&\leq n^2+16n\sqrt{nr_0}.
\end{aligned}
\]
% lean: CausalSmith.Experimentation.BivariateSobolevDesignCapacity.signingPartitionCross_average_le

\item For distinct indices, independence of the selectors gives
\[
\mathbb P(i\in\mathcal I_\eta,\ i'\in\mathcal J_\eta)=\frac14,
\]
whereas an index cannot belong to both groups. Consequently, for every deterministic row tuple and signing,
\[
\mathsf O_{\rm row}(\boldsymbol v,z)
=4\mathbb E_\eta\mathsf C_{\rm cross}(\boldsymbol v,z,\eta).
\]
Taking absolute values and then maximizing yields
\[
\begin{aligned}
\mathsf M_{\rm off}(\boldsymbol v)
&=4\max_z
\left|\mathbb E_\eta\mathsf C_{\rm cross}(\boldsymbol v,z,\eta)\right|\\
&\leq4\max_z
\mathbb E_\eta|\mathsf C_{\rm cross}(\boldsymbol v,z,\eta)|\\
&\leq4\mathbb E_\eta\mathsf M_{\rm cross}(\boldsymbol v,\eta).
\end{aligned}
\]
Integrating and applying the partition-average estimate gives
\[
\mathbb E\mathsf M_{\rm off}(\boldsymbol v)
\leq4\mathbb E_{\boldsymbol v}\mathbb E_\eta
\mathsf M_{\rm cross}(\boldsymbol v,\eta)
\leq n^2+16n\sqrt{nr_0}.
\]
% lean: CausalSmith.Experimentation.BivariateSobolevDesignCapacity.signingOffDiagonal_integral_le_partition_average

The diagonal-subtraction inequality and its expected diagonal energy now give
\[
\begin{aligned}
\mathbb E\mathsf M_{\rm sign}(\boldsymbol v)
&\geq nr_0-\mathbb E\mathsf M_{\rm off}(\boldsymbol v)\\
&\geq nr_0-n^2-16n\sqrt{nr_0}.
\end{aligned}
\]
Together with the small-sample case, this proves the first assertion for every \(n\).
% lean: CausalSmith.Experimentation.BivariateSobolevDesignCapacity.isotropic_row_signing

\item Finally, assume \(r_0\geq4096n\). Multiplying this inequality by the nonnegative number \(r_0/4096\) gives
\[
nr_0\leq\frac{r_0^2}{4096}
=\left(\frac{r_0}{64}\right)^2,
\qquad
\sqrt{nr_0}\leq\frac{r_0}{64}.
\]
Multiplying by nonnegative factors also gives
\[
16n\sqrt{nr_0}\leq\frac{nr_0}{4},
\qquad
n^2\leq\frac{nr_0}{4096}.
\]
Hence
\[
n^2+16n\sqrt{nr_0}
\leq\left(\frac1{4096}+\frac14\right)nr_0
\leq\frac{nr_0}{2}.
\]
These inequalities include \(n=0\). Subtracting this error bound from the first assertion proves
\[
\mathbb E\mathsf M_{\rm sign}(\boldsymbol v)\geq\frac{nr_0}{2}.
\]
% lean: CausalSmith.Experimentation.BivariateSobolevDesignCapacity.isotropic_signing_error_le_half
\end{enumerate}
\end{proof}

The expected supremum estimate underlying \cref{obj:lem:isotropic-row-signing} uses the absolute-supremum contraction inequality of \citet[Theorem~12(4), with Rademacher complexity as in Definition~2]{BartlettMendelson2002Complexities}, applied to the absolute-value map. The classes used here are finite classes and the finite-dimensional Euclidean unit-ball linear class. For the latter, a countable dense subset approximates evaluations simultaneously on every finite sample, and the vector norm provides an integrable envelope under square integrability. These properties supply the measurability and integrability conditions for the cited inequality.

The minimum in \cref{obj:lem:isotropic-row-signing} allows the sign vector to depend on all realized vectors. Consequently, its lower bound also controls the expected squared imbalance of any randomized assignment evaluated on those vectors. When the vector dimension is at least \(4096n\), the bound retains at least half of the total second-moment scale \(nr_0\).

Apply this comparison with \(v_i=V(X_i)\) and \(r_0=M\). Independence of the covariate vectors follows from \cref{obj:ass:uniform-draw}, while their feature second moments are given by \cref{obj:lem:cosine-prior}. With \(B=\binom d2\), the coordinate-pair count, choose the cutoff specified in \cref{obj:def:cosine-prior},
\[
L=\max\left\{1,\left\lceil\sqrt{\frac{4096n}{B}}\right\rceil\right\}.
\]
Then \(M=BL^2\geq4096n\), and the sign-minimum comparison becomes
\[
\mathbb E_X\min_{z\in\{-1,1\}^n}
\left\|\sum_{i=1}^n z_iV(X_i)\right\|_2^2
\geq\frac{nM}{2}.
\]

For an admissible design \(\pi\) from \cref{obj:def:design-class}, averaging the squared prognostic imbalance over the independent coefficient signs gives
\[
\mathbb E_\xi\!\left[
\frac4n\mathbb E_{X,Z}
\left(\sum_{i=1}^n Z_i m_\xi(X_i)\right)^2
\right]
=
\frac{4}{nC_0^2L^{2s}M}
\mathbb E_{X,Z}
\left\|\sum_{i=1}^n Z_iV(X_i)\right\|_2^2.
\]
The assignment law depends on the covariates as specified in \cref{obj:def:design-class}, and the coefficient draw is independent of that experiment. Together with the sign-minimum bound, this is the prior comparison used in \cref{obj:thm:full-design-lower}:
\[
\sup_{m\in\mathcal H_{d,s}}
\frac4n\mathbb E_{X,Z}
\left(\sum_{i=1}^n Z_i m(X_i)\right)^2
\geq
\mathbb E_\xi\!\left[
\frac4n\mathbb E_{X,Z}
\left(\sum_{i=1}^n Z_i m_\xi(X_i)\right)^2
\right]
\geq \frac{2}{C_0^2L^{2s}}.
\]
Each function in this finite average is fixed before sampling and belongs to the class by \cref{obj:lem:cosine-prior}. Thus the comparison preserves the fixed-function supremum outside the sample expectation. The chosen cutoff has order \(\max\{1,\sqrt{n/B}\}\), yielding the saturated scale \(\min\{1,(B/n)^s\}\). Taking the infimum over admissible designs gives the full-design lower comparison in \cref{obj:thm:full-design-lower}.


\section{Proof assembly and verification note}

The preceding auxiliary results supply the regularity, assignment, and converse ingredients of the capacity characterization. This section collects their implications for the minimax criterion, the dimension threshold, and causal precision, as stated in \cref{obj:thm:capacity-answer,obj:thm:log-free-bivariate-capacity}. The following complete characterization assembles the common assignment procedure, the independent-prior certificate, and the exact causal interpretation.

\begin{theoremv}{thm:capacity-answer}[Complete capacity and excess characterization]*
For every pair of integers \(n,d\geq2\), let \(B\), the coordinate-pair count, and \(a_s(n,d)\), the comparison scale, be
\[
B=\binom d2,\qquad
a_s(n,d)=\min\{1,(B/n)^s\}.
\]
There is a single fair, outcome-independent Borel assignment kernel \(\pi^\star_{n,d}\) in the design class of \cref{obj:def:design-class}, using the original \(n\) units and independent of \(s\). For \(n\leq B\), this kernel assigns independent fair signs. For \(B<n\), it is the law of the prescribed spectral signing procedure, using independent reflection coins, an independent common Haar shift, and independent uniform rounding seeds. The rounding procedure starts at the zero fractional assignment and uses the first \(K=\lfloor n/4\rfloor\) prescribed real spectral rows evaluated at the reflected and commonly shifted sample. For every realization of the reflection coins and common shift, and every sequence of rounding seeds in \([0,1]\), its fractional assignment after \(n\) iterations equals the resulting assignment sign vector:
\[
u_i=Z_i,\qquad i=1,\ldots,n.
\]

For every \(0<s\leq1\), the lower comparison scale satisfies \(b_s(n,d)=a_s(n,d)\). With
\[
c_s=\frac{2}{(48+32\pi^2)16384^s},
\qquad
c_1=\frac{2}{(48+32\pi^2)16384},
\]
the capacity and signing loss of \cref{obj:def:criterion}, over the function class of \cref{obj:def:sobolev-class}, satisfy
\[
0<c_1\leq c_s,\qquad
c_sa_s(n,d)
\leq R_s(n,d)
\leq
\sup_{m\in\mathcal H_{d,s}}
\mathcal L_{n,d,s}(\pi^\star_{n,d},m)
\leq3072a_s(n,d).
\]
These inequalities hold simultaneously for the same kernel \(\pi^\star_{n,d}\).

Every coefficient-sign realization \(m_\xi\) of the finite pair-cosine prior in \cref{obj:def:cosine-prior}, with the cutoff and normalization displayed there, is centered and belongs to \(\mathcal H_{d,s}\). Drawing its coefficient signs independently of the covariate sample gives, for every admissible design \(\pi\),
\[
c_sa_s(n,d)
\leq
\mathbb E_\xi\!\left[\mathcal L_{n,d,s}(\pi,m_\xi)\right].
\]

For every admissible \(\pi\) and every \(m\in\mathcal H_{d,s}\), the signing loss is finite, and the actual Horvitz--Thompson estimator under the Gaussian completion satisfies
\[
n\operatorname{Var}(\widehat\theta)-4.01
=
\mathcal L_{n,d,s}(\pi,m).
\]
Consequently, taking the infimum over admissible designs and the supremum over Gaussian completions indexed by \(m\in\mathcal H_{d,s}\) gives exactly \(R_s(n,d)\); taking only the supremum for \(\pi^\star_{n,d}\) gives exactly its worst signing loss.

Let \(\mathcal Q_{d,s}\) be the frozen uniform outcome-law class: the covariate is uniform on \([0,1]^d\), both potential outcomes have finite second moments, and their conditional midpoint mean is a centered member of \(\mathcal H_{d,s}\). For a law \(P\) in this class, write
\[
m_P(u)=\mathbb E_P[(O(1)+O(0))/2\mid U=u],
\qquad
\widehat T_P=\frac2n\sum_{i=1}^n Z_iO_i((Z_i+1)/2),
\]
and
\[
\mathsf V(P)
=
\operatorname{Var}_P\!\left(
\mathbb E_P[O(1)-O(0)\mid U]\right)
+
2\mathbb E_P\operatorname{Var}_P(O(1)\mid U)
+
2\mathbb E_P\operatorname{Var}_P(O(0)\mid U).
\]
Here the estimator uses independent original-unit copies under \(P\), with assignment generated by \(\pi\). The exact published-law minimax excess is
\[
\inf_{\pi\in\mathcal D_{n,d,s}}
\sup_{P\in\mathcal Q_{d,s}}
\max\!\left\{
n\operatorname{Var}_{P,\pi}(\widehat T_P)-\mathsf V(P),\,0
\right\}
=
R_s(n,d).
\]

Finally, for every fixed \(0<s\leq1\) and every integer dimension sequence \((d_n)\) satisfying \(d_n\geq2\) for every \(n\),
\[
R_s(n,d_n)\longrightarrow0
\quad\Longleftrightarrow\quad
\frac{\binom{d_n}{2}}{n}\longrightarrow0
\quad\Longleftrightarrow\quad
d_n=o(\sqrt n),
\qquad n\longrightarrow\infty.
\]
\end{theoremv}
% CAUSALSMITH-CITED-SCOPE-BEGIN thm:capacity-answer
\verificationfootnotetext{\textbf{Formalization scope.} This result uses the cited conclusion \citep{BartlettMendelson2002Complexities} (Theorem 12(4), with Rademacher complexity as in Definition 2) and \citep{Cytrynbaum2026Limits} (Proposition 2.3 (Excess Variance), equation (2.4), with Assumption 2.1 and Definition 2.2). The cited source proof is not formalized here: Lean verifies its use and all remaining steps. If that cited conclusion is false, the portions of this result that depend on it are not certified.}
% CAUSALSMITH-CITED-SCOPE-END thm:capacity-answer

The common scale expresses the number of pairwise interactions relative to the available units. When that ratio is small, the spectral assignment controls prognostic imbalance at a rate determined by smoothness; the independent coefficient prior gives a matching comparison over every admissible design. The exact variance identities then translate this balance criterion into causal precision. The same finite assignment procedure supplies the upper comparison throughout the stated smoothness range.

\subsection{The common comparison scale}

The spectral assignment bound combines the pooled reflection budget in \cref{obj:lem:exact-reflection-budget}, the imbalance identity in \cref{obj:lem:reflection-transfer}, and the assignment guarantees in \cref{obj:lem:ordered-prefix-rounding}. The independent coefficient prior supplies the full-design converse in \cref{obj:thm:full-design-lower}. Both comparisons use the same saturated scale identified in \cref{obj:thm:capacity-answer}. Here the function and design classes are those of \cref{obj:def:sobolev-class,obj:def:design-class}, and the loss is specified in \cref{obj:thm:log-free-bivariate-capacity}. The assignment procedure in \cref{obj:def:capacity-handle} attains this comparison simultaneously across the stated smoothness range.

For every fixed \(s\in(0,1]\), the dimension conclusion of \cref{obj:thm:capacity-answer} applies to every integer sequence \(d_n\geq2\): vanishing minimax scaled excess is equivalent to a vanishing interaction-to-sample ratio and to \(d_n=o(\sqrt n)\). Thus smoothness determines the rate within the subcritical regime, while the threshold for vanishing minimax scaled excess is expressed by the interaction count relative to sample size. The equivalence accommodates arbitrary dimension sequences under fixed smoothness.

\subsection{Outcome-law capacity}

The causal interpretation rests on the exact transfer in \cref{obj:lem:published-prognostic-capacity}. For every admissible design, the uniform outcome-law class \(\mathcal Q_{d,s}\) defined there satisfies
\[
\sup_{P\in\mathcal Q_{d,s}}\mathcal V_n(\pi,P)
=
\sup_{m\in\mathcal H_{d,s}}\mathcal L_{n,d,s}(\pi,m).
\]
The outcome-law criterion therefore has the same minimax value \(R_s(n,d)\), and its worst-case value for \(\pi^\star_{n,d}\) equals that procedure's worst signing loss. These equalities retain finite potential-outcome second moments, uniform covariates, centered prognosis, and the assignment conditions of \cref{obj:def:design-class}. The transfer uses the excess-variance identity of \citet[Proposition 2.3, equation (2.4), under Assumption 2.1 and Definition 2.2]{Cytrynbaum2026Limits}.

The Gaussian completions in \cref{obj:synth_4} realize the entire prognostic class. As recorded in \cref{obj:thm:capacity-answer}, restricting the outcome-law supremum to those completions preserves both the minimax capacity and the worst-case loss of the specified procedure. For the larger bounded-density class defined in \cref{obj:lem:published-prognostic-capacity}, the class transfer instead supplies a lower comparison. With fixed density bounds \(0<\lambda\leq1\leq\Lambda<\infty\) and fixed smoothness, vanishing minimax scaled excess over that class requires \(d_n=o(\sqrt n)\).

\subsection{Relative-excess precision requirements}

Under the Gaussian completion, \cref{obj:lem:ht-transfer} identifies signing loss with actual scaled HT variance above \(V^*=4.01\). Consequently, the upper comparison in \cref{obj:thm:capacity-answer} gives the following sufficient precision requirement, for any \(\varepsilon>0\):
\[
n\geq B\max\left\{
1,\left(\frac{3072}{4.01\varepsilon}\right)^{1/s}
\right\}.
\]
At this sample size, the procedure \(\pi^\star_{n,d}\) satisfies
\[
\sup_{m\in\mathcal H_{d,s}}
n\operatorname{Var}(\widehat\theta)
\leq4.01(1+\varepsilon),
\]
with each variance evaluated under the corresponding Gaussian completion.

The lower comparison in \cref{obj:thm:capacity-answer} gives the necessary requirement for \(0<\varepsilon<c_s/4.01\):
\[
R_s(n,d)\leq4.01\varepsilon
\quad\Longrightarrow\quad
n\geq B\left(\frac{c_s}{4.01\varepsilon}\right)^{1/s}.
\]
It therefore applies to any admissible design attaining the stated relative-excess target uniformly over the prognostic class. Within this target range, the sufficient and necessary requirements both scale with the coordinate-pair count and with \(\varepsilon^{-1/s}\), up to the displayed constants.

\paragraph{Verification note.}
The theorem statements and proofs are machine-checked in Lean 4 under their stated hypotheses. Uniform sampling, exact order-two prognosis, the pooled Sobolev budget, conditional fairness, and outcome-independent assignment are stipulated model conditions in \cref{obj:ass:uniform-draw,obj:ass:order-two,obj:ass:pooled-budget,obj:ass:fair,obj:ass:assignment-independence}; the capacity bounds, transfer equalities, assignment guarantees, and precision implications are proved conclusions. The theorem-local verification disclosures record the cited Rademacher contraction conclusion \citep[Theorem 12(4), with Rademacher complexity as in Definition 2]{BartlettMendelson2002Complexities} and the cited excess-variance identity \citep[Proposition 2.3, equation (2.4), under Assumption 2.1 and Definition 2.2]{Cytrynbaum2026Limits}, wherever each is used. Those published conclusions enter as premises whose source proofs were not formalized: Lean checks their applications and the remaining deductions, with certification of dependent conclusions conditional on the cited premises.


\section{Proofs of the main results}\label{sec:deferred-proofs}

\begin{proof}[Proof of \cref{obj:prop:published-class-inclusion}]
Fix \(0<s\leq1\). We use the published excess-variance identity
\citep[Proposition 2.3, equation (2.4), under Assumption 2.1 and Definition 2.2]{Cytrynbaum2026Limits}
and the expected absolute-supremum contraction inequality
\citep[Theorem 12(4), with Definition 2]{BartlettMendelson2002Complexities}
through \cref{obj:lem:published-prognostic-capacity,obj:thm:full-design-lower}.

\begin{enumerate}
\item For every integer \(d\geq2\), \cref{obj:lem:published-prognostic-capacity} supplies the proper inclusion
\[
\left\{
P:\ P\text{ is the Gaussian completion of some }
m\in\mathcal H_{d,s}\text{ specified in }\cref{obj:synth_4}
\right\}
\subsetneq\mathcal Q_{d,s}.
\]
This is the first assertion.
% lean: CausalSmith.Experimentation.BivariateSobolevDesignCapacity.published_prognostic_capacity

\item Fix integers \(n,d\geq2\) and a Borel assignment probability kernel \(\pi\) based on the original covariate array. The admissibility equivalence in
\cref{obj:lem:published-prognostic-capacity} states that
\(\pi\in\mathcal D_{n,d,s}\) precisely when
\[
\mathbb E_\pi[Z_i\mid X]=0
\quad\text{almost surely},\qquad i=1,\ldots,n,
\]
under independent uniform covariates, and
\[
Z\ \perp\!\!\!\perp\
\bigl((O_i(0),O_i(1))\bigr)_{i=1}^n
\ \bigm|\ X
\]
in the experiment formed from \(n\) independent copies of every single-unit law \(P\) having bounded strictly positive covariate density and finite potential-outcome second moments, followed by assignment from \(\pi(X)\). These are exactly the stated specialization of the published admissibility conditions.
% lean: CausalSmith.Experimentation.BivariateSobolevDesignCapacity.published_prognostic_capacity

\item Fix an integer \(n\geq2\). With the lower constant and comparison scale defined by
\[
c_s=\frac{2}{(48+32\pi^2)\,16384^s},
\qquad
b_s(n,d)=\min\left\{1,\left(\frac{\binom d2}{n}\right)^s\right\},
\]
\cref{obj:thm:full-design-lower} gives
\[
R_s(n,2)\geq c_s b_s(n,2).
\]
Since \(\binom22=1\), \(n\geq1\), and \(-s<0\),
\[
n^{-s}\leq1,
\qquad
b_s(n,2)=\min\{1,n^{-s}\}=n^{-s}.
\]
Substitution yields
\[
R_s(n,2)\geq
\frac{2}{(48+32\pi^2)\,16384^s}\,n^{-s}.
\]
% lean: CausalSmith.Experimentation.BivariateSobolevDesignCapacity.full_design_lower
% lean: CausalSmith.Experimentation.BivariateSobolevDesignCapacity.bScale_two

\item Fix integers \(n,d\geq2\) and \(m\in\mathcal H_{d,s}\). Define the constant assignment kernel on \(([0,1]^d)^n\) by
\[
\pi_{\mathrm{fair}}(\,\cdot\,)
=
\bigotimes_{i=1}^n
\left(\frac12\delta_{-1}+\frac12\delta_{+1}\right).
\]
Its signs are independent and fair, independently of the original covariate sample. Consequently, for \(i,i'\in\{1,\ldots,n\}\),
\[
\mathbb E_{\pi_{\mathrm{fair}}}[Z_iZ_{i'}\mid X]
=
\begin{cases}
1,&i=i',\\
0,&i\ne i',
\end{cases}
\]
and expansion of the finite square gives
\[
\mathbb E_{\pi_{\mathrm{fair}}}
\left[
\left(\sum_{i=1}^n Z_i m(X_i)\right)^2
\middle|X
\right]
=
\sum_{i=1}^n m(X_i)^2.
\]
Square-integrability of \(m\) justifies integrating this identity. Each covariate marginal is uniform, so
\[
\begin{aligned}
\frac4n\,\mathbb E
\left[
\left(\sum_{i=1}^n Z_i m(X_i)\right)^2
\right]
&=\frac4n\sum_{i=1}^n\mathbb E[m(X_i)^2]\\
&=4\int_{[0,1]^d}m(x)^2\,dx\\
&=4\mathbb E[m(X_1)^2].
\end{aligned}
\]
Finally, \cref{obj:lem:exact-reflection-budget} supplies
\[
\int_{[0,1]^d}|m(x)|^2\,dx\leq1.
\]
Since \(m\) is real-valued, multiplying this inequality by \(4\) proves the claimed upper bound.
% lean: CausalSmith.Experimentation.BivariateSobolevDesignCapacity.loss_fairSigns_eq
% lean: CausalSmith.Experimentation.BivariateSobolevDesignCapacity.sobolev_cube_second_moment_le_one
\end{enumerate}
\end{proof}

\begin{proof}[Proof of \cref{obj:thm:log-free-bivariate-capacity}]
Fix integers \(n,d\geq2\), and set
\[
B:=\binom d2,\qquad K:=\lfloor n/4\rfloor.
\]
In particular, \(B>0\). Expectations of squares and sums of nonnegative terms are initially interpreted in the extended nonnegative reals.

\begin{enumerate}
\item
We use the expected absolute-supremum contraction inequality of
\citep[Theorem 12(4), with the complexity convention of Definition 2]{BartlettMendelson2002Complexities}
through \cref{obj:thm:full-design-lower,obj:lem:ht-transfer}. The outcome-law comparison uses the excess-variance identity of
\citep[Proposition 2.3, equation (2.4), under Assumption 2.1 and Definition 2.2]{Cytrynbaum2026Limits}
through \cref{obj:lem:published-prognostic-capacity}.

For every \(0<s\leq1\) and every admissible \(\pi\), \cref{obj:thm:full-design-lower} supplies
\[
c_s b_s(n,d)
\leq
\mathbb E_{m\sim\nu^\star_{n,d,s}}
\mathcal L_{n,d,s}(\pi,m),
\qquad
c_s b_s(n,d)\leq R_s(n,d).
\]
Also, directly from \cref{obj:def:criterion}, selecting any admissible kernel in the design infimum gives
\[
R_s(n,d)
\leq
\sup_{m\in\mathcal H_{d,s}}\mathcal L_{n,d,s}(\pi,m).
\]
We will apply this inequality to the kernel constructed below.
% lean: CausalSmith.Experimentation.BivariateSobolevDesignCapacity.full_design_lower
% lean: CausalSmith.Experimentation.BivariateSobolevDesignCapacity.published_prognostic_capacity
% lean: CausalSmith.Experimentation.BivariateSobolevDesignCapacity.log_free_bivariate_capacity

\item
We first construct the kernel using \(n,d\). For clarity, the spectral notation of \cref{obj:synth_3} is
\[
\mathbb T_2^d:=(\mathbb R/2\mathbb Z)^d,
\qquad
\operatorname{supp}(k):=\{j\in\{1,\ldots,d\}:k_j\ne0\},
\]
\[
t(k):=
\begin{cases}
\|k\|_2^2/|\operatorname{supp}(k)|,&k\ne0,\\
0,&k=0,
\end{cases}
\qquad
e_k(y):=\exp(\mathrm i\pi k^{\mathsf T}y)
\quad(k\in\mathbb Z^d,\ y\in\mathbb T_2^d).
\]
Let
\[
\mathscr K_d^+
:=
\left\{
k\in\mathbb Z^d:
1\leq|\operatorname{supp}(k)|\leq2,\
\text{the first nonzero coordinate of }k\text{ is positive}
\right\}.
\]
For \(h_{\mathrm{freq}}\in\mathbb N_0\), define \(k^{[h_{\mathrm{freq}}]}\) to be the
\((h_{\mathrm{freq}}+1)\)-st member of \(\mathscr K_d^+\), ordered by increasing \(t(k)\), with lexicographic ties. The prescribed real features are
\[
\phi_{2h_{\mathrm{freq}}+1}(y)
:=\sqrt2\,\operatorname{Re}e_{k^{[h_{\mathrm{freq}}]}}(y),
\qquad
\phi_{2h_{\mathrm{freq}}+2}(y)
:=\sqrt2\,\operatorname{Im}e_{k^{[h_{\mathrm{freq}}]}}(y).
\]
Since \(|e_k(y)|=1\),
\[
|\phi_h(y)|\leq\sqrt2
\qquad(h\geq1,\ y\in\mathbb T_2^d).
\]
For \(w=(w_1,\ldots,w_n)\in(\mathbb T_2^d)^n\), define
\[
a_h(w):=(\phi_h(w_i))_{i=1}^n,
\qquad
a(w):=(\phi_h(w_i))_{\substack{1\leq h\leq K\\1\leq i\leq n}}.
\]
Let
\[
Z^{\mathrm{sp}}(w,\zeta)
:=
\text{the output of ordered boundary rounding on }(a(w),\zeta),
\qquad
\zeta\in[0,1]^{3n},
\]
where the procedure is precisely \cref{obj:def:ordered-boundary-rounding}.

Draw independently a fair reflection array
\(\eta\in\{-1,1\}^{n\times d}\), a normalized Haar shift
\(T\in\mathbb T_2^d\), and a seed vector \(\zeta\) with independent uniform entries.
For every deterministic original sample
\(\mathbf x=(x_1,\ldots,x_n)\in([0,1]^d)^n\), define, modulo two,
\[
[\operatorname{lift}(\mathbf x,\eta)_i]_j
:=
\begin{cases}
(x_i)_j,&\eta_{ij}=1,\\
2-(x_i)_j,&\eta_{ij}=-1,
\end{cases}
\qquad
W_i(\mathbf x,\eta,T)
:=
\operatorname{lift}(\mathbf x,\eta)_i+T.
\]
The construction in \cref{obj:def:capacity-handle} is therefore
\[
\pi^\star_{n,d}(\mathbf x)
:=
\begin{cases}
\operatorname{Unif}(\{-1,1\}^n),&n\leq B,\\
\operatorname{Law}_{\eta,T,\zeta}
\bigl(Z^{\mathrm{sp}}(W(\mathbf x,\eta,T),\zeta)\bigr),&B<n.
\end{cases}
\]

For every fixed \(\mathbf x,\eta,T\), \cref{obj:lem:ordered-prefix-rounding}, applied with row bound \(\sqrt2\), gives
\[
\mathbb E_\zeta
Z_i^{\mathrm{sp}}(W(\mathbf x,\eta,T),\zeta)=0.
\]
Integration over \(\eta,T\), together with the independent-sign branch, yields
\[
\int_{\{-1,1\}^n}z_i\,\pi^\star_{n,d}(\mathbf x,dz)=0
\qquad(1\leq i\leq n)
\]
for every original sample. The feature evaluations, reflection lift, torus addition, and rounding output are Borel measurable; integrating their output against the probability seed law gives a Borel probability kernel. Drawing assignments from this covariate kernel gives
\[
Z\perp Y\mid X
\]
for every pre-assignment law allowed in \cref{obj:ass:assignment-independence}. Consequently,
\[
\pi^\star_{n,d}\in\mathcal D_{n,d,s}.
\]
Its branch choice, feature ordering, and seed distribution are independent of \(s\).
% lean: CausalSmith.Experimentation.BivariateSobolevDesignCapacity.piStar_pointwise_fair
% lean: CausalSmith.Experimentation.BivariateSobolevDesignCapacity.piStar_designClass

\item
We prove the upper guarantee for an arbitrary \(0<s\leq1\). First suppose \(n\leq B\). Conditional on \(X\), independent fair signs cancel every off-diagonal term:
\[
\mathbb E_Z\left[
\left(\sum_{i=1}^n Z_i m(X_i)\right)^2\Bigm|X
\right]
=
\sum_{i=1}^n m(X_i)^2.
\]
Thus, for every \(m\in\mathcal H_{d,s}\),
\[
\mathcal L_{n,d,s}(\pi^\star_{n,d},m)
=
4\int_{[0,1]^d}m(x)^2\,dx
\leq4,
\]
where the last inequality is supplied by \cref{obj:lem:exact-reflection-budget}.
Moreover,
\[
B/n\geq1
\quad\Longrightarrow\quad
(B/n)^s\geq1
\quad\Longrightarrow\quad
a_s(n,d)=1.
\]
Hence
\[
\mathcal L_{n,d,s}(\pi^\star_{n,d},m)
\leq4\leq3072a_s(n,d).
\]
% lean: CausalSmith.Experimentation.BivariateSobolevDesignCapacity.fairSignKernel_loss_le_four

\item
Suppose now that \(B<n\). Fix \(w\in(\mathbb T_2^d)^n\).
The moment guarantee in \cref{obj:lem:ordered-prefix-rounding} gives, for every \(h\geq1\),
\[
\mathbb E_\zeta
\left[
\bigl(a_h(w)^{\mathsf T}Z^{\mathrm{sp}}(w,\zeta)\bigr)^2
\right]
\leq32(\sqrt2)^2\min\{h,n\}
=
64\min\{h,n\}.
\]
This includes rows beyond the supplied prefix. When \(K=0\), the input matrix has zero rows and the same guarantee applies.

The two real rows associated with \(k^{[h_{\mathrm{freq}}]}\) satisfy the pointwise identity
\[
\left|
\sum_{i=1}^n
Z_i^{\mathrm{sp}}(w,\zeta)
e_{k^{[h_{\mathrm{freq}}]}}(w_i)
\right|^2
=
\frac{
\bigl(a_{2h_{\mathrm{freq}}+1}(w)^{\mathsf T}Z^{\mathrm{sp}}(w,\zeta)\bigr)^2
+
\bigl(a_{2h_{\mathrm{freq}}+2}(w)^{\mathsf T}Z^{\mathrm{sp}}(w,\zeta)\bigr)^2
}{2}.
\]
Indeed, the row imbalances are \(\sqrt2\) times the real and imaginary parts of the character imbalance. Taking expectations therefore gives
\[
\mathbb E_\zeta
\left|
\sum_{i=1}^n
Z_i^{\mathrm{sp}}(w,\zeta)
e_{k^{[h_{\mathrm{freq}}]}}(w_i)
\right|^2
\leq64\min\{2h_{\mathrm{freq}}+2,n\}.
\]
% lean: CausalSmith.Experimentation.BivariateSobolevDesignCapacity.character_imbalance_eq_real_rows
% lean: CausalSmith.Experimentation.BivariateSobolevDesignCapacity.spectral_rounding_character_second_moment

\item
We convert the row index into frequency complexity. Define, for \(t_{\mathrm{cut}}\geq1\),
\[
N_{\mathrm{real}}(t_{\mathrm{cut}})
:=
2\#\{k\in\mathscr K_d^+:t(k)\leq t_{\mathrm{cut}}\}.
\]
A support-one representative has a positive integer amplitude at most
\(\lfloor\sqrt{t_{\mathrm{cut}}}\rfloor\). A support-two representative has an ordered coordinate pair, a positive amplitude on its first coordinate, either sign on its second coordinate, and both absolute amplitudes at most
\(\lfloor\sqrt{2t_{\mathrm{cut}}}\rfloor\). Consequently,
\[
N_{\mathrm{real}}(t_{\mathrm{cut}})
\leq
2d\lfloor\sqrt{t_{\mathrm{cut}}}\rfloor
+
4B\lfloor\sqrt{2t_{\mathrm{cut}}}\rfloor^2.
\]
For \(t_{\mathrm{cut}}\geq1\),
\[
\lfloor\sqrt{t_{\mathrm{cut}}}\rfloor
\leq\sqrt{t_{\mathrm{cut}}}\leq t_{\mathrm{cut}},
\qquad
\lfloor\sqrt{2t_{\mathrm{cut}}}\rfloor^2
\leq2t_{\mathrm{cut}}.
\]
Also,
\[
2B=d(d-1)\geq d.
\]
It follows that
\[
N_{\mathrm{real}}(t_{\mathrm{cut}})
\leq2dt_{\mathrm{cut}}+8Bt_{\mathrm{cut}}
\leq12Bt_{\mathrm{cut}}.
\]

Every nonzero integer frequency satisfies
\[
\|k\|_2^2\geq|\operatorname{supp}(k)|,
\qquad t(k)\geq1.
\]
The first \(h_{\mathrm{freq}}+1\) representatives all have complexity at most
\(t(k^{[h_{\mathrm{freq}}]})\), so
\[
2h_{\mathrm{freq}}+2
\leq
N_{\mathrm{real}}\bigl(t(k^{[h_{\mathrm{freq}}]})\bigr)
\leq12B\,t(k^{[h_{\mathrm{freq}}]}).
\]
If \(Bt(k^{[h_{\mathrm{freq}}]})\leq n\), the preceding row-index bound gives a character moment at most \(768Bt(k^{[h_{\mathrm{freq}}]})\). If
\(n<Bt(k^{[h_{\mathrm{freq}}]})\), the bound by \(64n\) gives a character moment at most \(768n\). Thus
\[
\mathbb E_\zeta
\left|
\sum_{i=1}^n
Z_i^{\mathrm{sp}}(w,\zeta)
e_{k^{[h_{\mathrm{freq}}]}}(w_i)
\right|^2
\leq
768\min\{n,Bt(k^{[h_{\mathrm{freq}}]})\}.
\]

Every support-one or support-two frequency is a listed representative or its negative: its first nonzero coordinate determines the orientation, and there are finitely many representatives below each fixed complexity. Since real signs give
\[
\sum_i Z_i^{\mathrm{sp}}(w,\zeta)e_{-k}(w_i)
=
\overline{\sum_i Z_i^{\mathrm{sp}}(w,\zeta)e_k(w_i)},
\qquad
t(-k)=t(k),
\]
we obtain, for every supported frequency and every deterministic \(w\),
\[
\mathbb E_\zeta
\left|
\sum_{i=1}^n Z_i^{\mathrm{sp}}(w,\zeta)e_k(w_i)
\right|^2
\leq768\min\{n,Bt(k)\},
\qquad
1\leq|\operatorname{supp}(k)|\leq2.
\]
% lean: CausalSmith.Experimentation.BivariateSobolevDesignCapacity.featureFrequency_real_rows_le_complexity
% lean: CausalSmith.Experimentation.BivariateSobolevDesignCapacity.spectral_rounding_supported_second_moment

\item
Fix \(m\in\mathcal H_{d,s}\). Define its reflected function and Fourier coefficients by
\[
F(y):=m(b_d(y)),
\qquad
[b_d(y)]_j:=\min\{y_j,2-y_j\}
\quad(y\in[0,2)^d),
\]
\[
\widehat F(k)
:=
2^{-d}\int_{[0,2)^d}
F(y)\exp(-\mathrm i\pi k^{\mathsf T}y)\,dy
\qquad(k\in\mathbb Z^d).
\]
Under the independent uniform original sample, \cref{obj:lem:reflection-transfer} gives
\[
\operatorname{Law}(W,T)
=
\operatorname{Haar}(\mathbb T_2^d)^{\otimes n}
\otimes\operatorname{Haar}(\mathbb T_2^d)
\]
with normalized Haar laws. Its spectral identity gives
\[
\mathcal L_{n,d,s}(\pi^\star_{n,d},m)
=
\frac4n
\sum_{k\in\mathbb Z^d}
|\widehat F(k)|^2
\mathbb E_{W,\zeta}
\left|
\sum_{i=1}^n Z_i^{\mathrm{sp}}(W,\zeta)e_k(W_i)
\right|^2.
\]
Here \(W\) has the product Haar law, independently of the rounding seeds.

By \cref{obj:lem:exact-reflection-budget},
\[
\widehat F(0)=0,
\qquad
\widehat F(k)=0
\quad\text{if }|\operatorname{supp}(k)|>2,
\]
and
\[
\sum_{\substack{k\in\mathbb Z^d\\
1\leq|\operatorname{supp}(k)|\leq2}}
(1+\pi^2t(k))^s|\widehat F(k)|^2
\leq1.
\]
Integrating the conditional character bound and using
\[
\frac4n\,768\min\{n,Bt(k)\}
=
3072\min\{1,(B/n)t(k)\}
\]
therefore yields
\[
\mathcal L_{n,d,s}(\pi^\star_{n,d},m)
\leq
3072
\sum_{\substack{k\in\mathbb Z^d\\
1\leq|\operatorname{supp}(k)|\leq2}}
|\widehat F(k)|^2\min\{1,(B/n)t(k)\}.
\]
The integration and summation are justified by nonnegativity.
% lean: CausalSmith.Experimentation.BivariateSobolevDesignCapacity.spectralRounding_loss_fourier
% lean: CausalSmith.Experimentation.BivariateSobolevDesignCapacity.spectralRounding_loss_le_truncated

\item
We bound this truncated sum by the comparison scale. For every
\(x_{\mathrm{scalar}}\geq0\),
\[
\min\{1,x_{\mathrm{scalar}}\}\leq x_{\mathrm{scalar}}^s.
\]
Indeed, when \(0\leq x_{\mathrm{scalar}}\leq1\),
\[
x_{\mathrm{scalar}}
=x_{\mathrm{scalar}}^1
\leq x_{\mathrm{scalar}}^s,
\]
and when \(x_{\mathrm{scalar}}>1\),
\[
1\leq x_{\mathrm{scalar}}^s.
\]
These inequalities include \(s=1\).

Put
\[
r_{\mathrm{sample}}:=B/n.
\]
For every \(t_{\mathrm{weight}}\geq0\), the inequalities
\[
t_{\mathrm{weight}}\leq1+\pi^2t_{\mathrm{weight}},
\qquad
1\leq1+\pi^2t_{\mathrm{weight}}
\]
give
\[
\begin{aligned}
\min\{1,r_{\mathrm{sample}}t_{\mathrm{weight}}\}
&\leq
(r_{\mathrm{sample}}t_{\mathrm{weight}})^s\\
&=
r_{\mathrm{sample}}^s t_{\mathrm{weight}}^s\\
&\leq
r_{\mathrm{sample}}^s(1+\pi^2t_{\mathrm{weight}})^s,
\end{aligned}
\qquad
\min\{1,r_{\mathrm{sample}}t_{\mathrm{weight}}\}
\leq1\leq(1+\pi^2t_{\mathrm{weight}})^s.
\]
Use the former bound if \(r_{\mathrm{sample}}^s\leq1\), and the latter if
\(r_{\mathrm{sample}}^s>1\). This proves
\[
\min\{1,r_{\mathrm{sample}}t_{\mathrm{weight}}\}
\leq
\min\{1,r_{\mathrm{sample}}^s\}
(1+\pi^2t_{\mathrm{weight}})^s
=
a_s(n,d)(1+\pi^2t_{\mathrm{weight}})^s.
\]
Multiplying by the squared Fourier coefficients and summing gives
\[
\begin{aligned}
&\sum_{\substack{k\in\mathbb Z^d\\
1\leq|\operatorname{supp}(k)|\leq2}}
|\widehat F(k)|^2\min\{1,(B/n)t(k)\}\\
&\qquad\leq
a_s(n,d)
\sum_{\substack{k\in\mathbb Z^d\\
1\leq|\operatorname{supp}(k)|\leq2}}
(1+\pi^2t(k))^s|\widehat F(k)|^2
\leq a_s(n,d).
\end{aligned}
\]
Combining both design branches and then taking the fixed-function supremum proves
\[
\sup_{m\in\mathcal H_{d,s}}
\mathcal L_{n,d,s}(\pi^\star_{n,d},m)
\leq3072a_s(n,d).
\]
The construction was fixed before choosing \(s\), so this guarantee holds for the same kernel for every \(0<s\leq1\).
% lean: CausalSmith.Experimentation.BivariateSobolevDesignCapacity.spectral_min_rate_weight
% lean: CausalSmith.Experimentation.BivariateSobolevDesignCapacity.sobolev_truncated_spectral_sum_le_scale
% lean: CausalSmith.Experimentation.BivariateSobolevDesignCapacity.spectral_design_upper

\item
We also record the pathwise terminal-state guarantee. For
\(a\in\mathbb R^{K\times n}\), \(\zeta\in[0,1]^{3n}\), and
\(h_{\mathrm{move}}\in\mathbb N_0\), write
\[
u^{[h_{\mathrm{move}}]}(a,\zeta)
:=
\text{the fractional state after }h_{\mathrm{move}}
\text{ iterations of ordered boundary rounding},
\qquad
u^{[0]}(a,\zeta)=0.
\]
Applying \cref{obj:lem:ordered-prefix-rounding} to the rows
\(a_h(w)\), whose entries are bounded by \(\sqrt2\), gives
\[
u_i^{[n]}(a(w),\zeta)=Z_i^{\mathrm{sp}}(w,\zeta)
\qquad(1\leq i\leq n)
\]
for every \(w\) and every seed vector in the closed cube. In particular,
\[
u_i^{[n]}
\bigl(a(W(\mathbf x,\eta,T)),\zeta\bigr)
=
Z_i^{\mathrm{sp}}(W(\mathbf x,\eta,T),\zeta)
\]
for every original sample, every reflection array, every common shift, and every permitted rounding seed vector. Together with the displayed definition of \(\pi^\star_{n,d}\), this proves the two branch descriptions and the asserted equality between fractional coordinates and output signs.
% lean: CausalSmith.Experimentation.BivariateSobolevDesignCapacity.ordered_prefix_rounding
% lean: CausalSmith.Experimentation.BivariateSobolevDesignCapacity.log_free_bivariate_capacity

\item
Now fix \(0<s\leq1\). The two scales agree by their definitions:
\[
a_s(n,d)=b_s(n,d)=\min\{1,(B/n)^s\}\geq0.
\]
If \(n<B\), then \(B/n>1\), so
\[
a_s(n,d)=1.
\]
If \(B\leq n\), then \(0<B/n\leq1\), so
\[
a_s(n,d)=(B/n)^s.
\]
In particular, both expressions equal one at \(B=n\).

The constant calibration follows from
\[
0<48+32\pi^2,
\qquad
1\leq16384^s\leq16384,
\]
which imply
\[
c_1=\frac{2}{(48+32\pi^2)16384}>0,
\qquad
c_1\leq c_s.
\]
Furthermore,
\[
(48+32\pi^2)16384^s\geq48,
\qquad
c_s\leq\frac2{48}<3072.
\]
Multiplication by the nonnegative scale, the lower certificate, the design-infimum inequality, and the upper guarantee now give
\[
c_1a_s(n,d)
\leq c_sa_s(n,d)
\leq R_s(n,d)
\leq
\sup_{m\in\mathcal H_{d,s}}
\mathcal L_{n,d,s}(\pi^\star_{n,d},m)
\leq3072a_s(n,d).
\]
% lean: CausalSmith.Experimentation.BivariateSobolevDesignCapacity.cLower_uniform_positive
% lean: CausalSmith.Experimentation.BivariateSobolevDesignCapacity.log_free_bivariate_capacity

\item
For the prior in \cref{obj:def:cosine-prior}, its cutoff, feature dimension, and normalization are
\[
L:=\max\left\{1,\left\lceil\sqrt{4096n/B}\right\rceil\right\},
\qquad
M:=BL^2,
\qquad
C_0:=\sqrt{48+32\pi^2}.
\]
Since \(L\geq1\), \cref{obj:lem:cosine-prior} applies to every coefficient-sign array and gives
\[
\int_{[0,1]^d}m_\xi(x)\,dx=0,
\qquad
m_\xi\in\mathcal H_{d,s}.
\]
The coefficients are drawn independently of \(X\), as specified in \cref{obj:def:cosine-prior}. Using \(a_s=b_s\), the prior inequality from \cref{obj:thm:full-design-lower} becomes
\[
c_sa_s(n,d)
\leq
\mathbb E_{m\sim\nu^\star_{n,d,s}}
\mathcal L_{n,d,s}(\pi,m)
\qquad(\pi\in\mathcal D_{n,d,s}).
\]
% lean: CausalSmith.Experimentation.BivariateSobolevDesignCapacity.cosinePrior_legal_and_isotropic
% lean: CausalSmith.Experimentation.BivariateSobolevDesignCapacity.full_design_lower

\item
For the frozen uniform outcome class, \cref{obj:lem:published-prognostic-capacity} identifies the admissible kernels with \(\mathcal D_{n,d,s}\), identifies the prognostic image with \(\mathcal H_{d,s}\) modulo null sets, and gives
\[
\mathcal V_n(\pi,P)
=
\mathcal L_{n,d,s}(\pi,m_P)
\qquad
(\pi\in\mathcal D_{n,d,s},\ P\in\mathcal Q_{d,s}).
\]
Consequently, for every admissible \(\pi\),
\[
\sup_{P\in\mathcal Q_{d,s}}\mathcal V_n(\pi,P)
=
\sup_{m\in\mathcal H_{d,s}}\mathcal L_{n,d,s}(\pi,m).
\]
Taking the common design infimum proves
\[
\inf_{\pi\in\mathcal D_{n,d,s}}
\sup_{P\in\mathcal Q_{d,s}}\mathcal V_n(\pi,P)
=
R_s(n,d).
\]
Selecting \(\pi=\pi^\star_{n,d}\) in the kernel-level equality gives
\[
\sup_{P\in\mathcal Q_{d,s}}
\mathcal V_n(\pi^\star_{n,d},P)
=
\sup_{m\in\mathcal H_{d,s}}
\mathcal L_{n,d,s}(\pi^\star_{n,d},m).
\]
% lean: CausalSmith.Experimentation.BivariateSobolevDesignCapacity.published_prognostic_capacity

\item
Define the Gaussian subclass by
\[
\mathcal G_{d,s}
:=
\left\{
P:
P\text{ is the single-unit Gaussian completion of some }
m\in\mathcal H_{d,s}
\right\},
\]
using the completion in \cref{obj:synth_4}. For the completion of \(m\), its conditional means satisfy
\[
m_P=m\quad\text{almost everywhere},
\qquad
\mathbb E_P[O(1)-O(0)\mid U]
=
0.2+0.1\sqrt2\sin(2\pi U_1).
\]
Its two conditional potential-outcome variances are one. Since \(U_1\) is uniform and
\[
\int_0^1\sqrt2\sin(2\pi v)\,dv=0,
\qquad
\int_0^1\bigl(\sqrt2\sin(2\pi v)\bigr)^2\,dv=1,
\]
the law-specific benchmark is
\[
\mathsf V(P)=0.01+2\cdot1+2\cdot1=4.01.
\]

For every admissible \(\pi\) and every \(m\in\mathcal H_{d,s}\),
\cref{obj:lem:ht-transfer} supplies the actual-estimator identity and finiteness:
\[
n\operatorname{Var}(\widehat\theta)-4.01
=
\mathcal L_{n,d,s}(\pi,m)<\infty.
\]
Under this completion, \(\widehat\theta=\widehat T_P\). Taking the supremum over the completions of all \(m\in\mathcal H_{d,s}\) therefore gives
\[
\sup_{P\in\mathcal G_{d,s}}\mathcal V_n(\pi,P)
=
\sup_{m\in\mathcal H_{d,s}}\mathcal L_{n,d,s}(\pi,m).
\]
Taking the admissible-design infimum, and separately selecting the constructed kernel, yields
\[
\inf_{\pi\in\mathcal D_{n,d,s}}
\sup_{P\in\mathcal G_{d,s}}\mathcal V_n(\pi,P)
=R_s(n,d),
\]
\[
\sup_{P\in\mathcal G_{d,s}}
\mathcal V_n(\pi^\star_{n,d},P)
=
\sup_{m\in\mathcal H_{d,s}}
\mathcal L_{n,d,s}(\pi^\star_{n,d},m).
\]
Finally, \cref{obj:lem:published-prognostic-capacity} supplies the proper inclusion
\[
\mathcal G_{d,s}\subsetneq\mathcal Q_{d,s}.
\]
% lean: CausalSmith.Experimentation.BivariateSobolevDesignCapacity.gaussianWorst_eq_worstLoss
% lean: CausalSmith.Experimentation.BivariateSobolevDesignCapacity.gaussianCapacity_eq_capacity
% lean: CausalSmith.Experimentation.BivariateSobolevDesignCapacity.ht_transfer
% lean: CausalSmith.Experimentation.BivariateSobolevDesignCapacity.published_prognostic_capacity

\item
For every pair of real bounds
\[
0<\mathrm{lower}\leq1\leq\mathrm{upper},
\]
the bounded-density comparison in \cref{obj:lem:published-prognostic-capacity} gives
\[
R_s(n,d)
\leq
\inf_{\pi\in\mathcal D_{n,d,s}}
\sup_{P\in\mathcal A_{d,s}}\mathcal V_n(\pi,P),
\]
where \(\mathcal A_{d,s}\) has the density, moment, prognostic-budget, and permitted-intercept conditions specified there. This proves the finite-sample comparison for the full displayed range of density bounds.
% lean: CausalSmith.Experimentation.BivariateSobolevDesignCapacity.published_prognostic_capacity

\item
Let \(\varepsilon>0\), and suppose
\[
n\geq
B\max\left\{1,
\left(\frac{3072}{4.01\varepsilon}\right)^{1/s}\right\}.
\]
Since \(B,n,4.01\varepsilon\) are positive,
\[
\left(\frac{3072}{4.01\varepsilon}\right)^{1/s}
\leq\frac nB.
\]
Raising both sides to the positive power \(s\) gives
\[
\frac{3072}{4.01\varepsilon}
\leq\left(\frac nB\right)^s.
\]
The reciprocal-power identity is
\[
\left(\frac Bn\right)^s
\left(\frac nB\right)^s=1.
\]
Multiplying the preceding inequality by \((B/n)^s\) and then by
\(4.01\varepsilon\) yields
\[
3072a_s(n,d)
\leq3072(B/n)^s
\leq4.01\varepsilon.
\]
Hence the upper guarantee proves
\[
\sup_{m\in\mathcal H_{d,s}}
\mathcal L_{n,d,s}(\pi^\star_{n,d},m)
\leq4.01\varepsilon.
\]
% lean: CausalSmith.Experimentation.BivariateSobolevDesignCapacity.scale_sample_size_sufficient

\item
Conversely, suppose
\[
0<\varepsilon<c_s/4.01,
\qquad
R_s(n,d)\leq4.01\varepsilon.
\]
The lower bound gives
\[
c_sa_s(n,d)\leq4.01\varepsilon<c_s.
\]
If \((B/n)^s\geq1\), then \(a_s(n,d)=1\), contradicting this strict inequality. Thus
\[
(B/n)^s<1,
\qquad
a_s(n,d)=(B/n)^s.
\]
Multiplying the lower inequality by \((n/B)^s\), using the reciprocal-power identity, and dividing by \(4.01\varepsilon>0\) gives
\[
\frac{c_s}{4.01\varepsilon}
\leq\left(\frac nB\right)^s.
\]
Taking the positive \(1/s\)-power and multiplying by \(B>0\) proves
\[
n\geq
B\left(\frac{c_s}{4.01\varepsilon}\right)^{1/s}.
\]
% lean: CausalSmith.Experimentation.BivariateSobolevDesignCapacity.scale_sample_size_necessary

\item
For the asymptotic claims, fix \(0<s\leq1\) and a sequence \(d_n\geq2\). Work on the tail \(n\geq2\), and define
\[
B_n:=\binom{d_n}{2}.
\]
First,
\[
a_s(n,d_n)\longrightarrow0
\quad\Longleftrightarrow\quad
B_n/n\longrightarrow0.
\]
For the forward direction, a vanishing scale is eventually less than one, so eventually
\[
a_s(n,d_n)=(B_n/n)^s,
\qquad
B_n/n=a_s(n,d_n)^{1/s}.
\]
Continuity of the positive \(1/s\)-power at zero gives \(B_n/n\to0\).
For the reverse direction, continuity of the positive \(s\)-power gives
\[
(B_n/n)^s\longrightarrow0,
\qquad
0\leq a_s(n,d_n)\leq(B_n/n)^s,
\]
and squeezing gives the asserted scale convergence.

The pair count satisfies
\[
B_n=\frac{d_n(d_n-1)}2,
\qquad
\frac{d_n^2}{4}\leq B_n\leq\frac{d_n^2}{2},
\]
because \(d_n-1\geq d_n/2\) and \(d_n-1\leq d_n\). Therefore
\[
0\leq\frac{d_n^2}{n}\leq4\frac{B_n}{n},
\qquad
0\leq\frac{B_n}{n}\leq\frac{d_n^2}{n},
\]
and
\[
B_n/n\longrightarrow0
\quad\Longleftrightarrow\quad
d_n^2/n\longrightarrow0.
\]
Since \(d_n\geq0\) and \(n>0\) on this tail,
\[
\frac{d_n}{\sqrt n}
=
\sqrt{\frac{d_n^2}{n}},
\qquad
\left(\frac{d_n}{\sqrt n}\right)^2
=
\frac{d_n^2}{n}.
\]
Continuity of square roots and squares, together with the defining ratio criterion for little \(o\), gives
\[
B_n/n\longrightarrow0
\quad\Longleftrightarrow\quad
d_n=o(\sqrt n).
\]
% lean: CausalSmith.Experimentation.BivariateSobolevDesignCapacity.pairCount_quadratic_bounds
% lean: CausalSmith.Experimentation.BivariateSobolevDesignCapacity.scale_vanishing

\item
Suppose \(R_s(n,d_n)\to0\). For every \(n\geq2\), the lower certificate gives
\[
0\leq c_sa_s(n,d_n)\leq R_s(n,d_n).
\]
Squeezing and division by the fixed constant \(c_s>0\) yield
\[
c_sa_s(n,d_n)\longrightarrow0,
\qquad
a_s(n,d_n)\longrightarrow0.
\]
The scale equivalences just proved imply
\[
d_n=o(\sqrt n),
\qquad
B_n/n\longrightarrow0.
\]

Conversely, if \(B_n/n\to0\), then \(a_s(n,d_n)\to0\). The constructed admissible kernel gives, for every \(n\geq2\),
\[
0\leq R_s(n,d_n)
\leq
\sup_{m\in\mathcal H_{d_n,s}}
\mathcal L_{n,d_n,s}(\pi^\star_{n,d_n},m)
\leq3072a_s(n,d_n).
\]
The right-hand side tends to zero, so squeezing gives \(R_s(n,d_n)\to0\). Thus
\[
R_s(n,d_n)\longrightarrow0
\quad\Longleftrightarrow\quad
\frac{\binom{d_n}{2}}n\longrightarrow0
\quad\Longleftrightarrow\quad
d_n=o(\sqrt n).
\]
% lean: CausalSmith.Experimentation.BivariateSobolevDesignCapacity.dimension_necessity_of_scale_lower
% lean: CausalSmith.Experimentation.BivariateSobolevDesignCapacity.log_free_bivariate_capacity

\item
Finally, fix density bounds satisfying
\(0<\mathrm{lower}\leq1\leq\mathrm{upper}\).
For these fixed bounds and every sequence \(d_n\geq2\),
\cref{obj:lem:published-prognostic-capacity} supplies
\[
\inf_{\pi\in\mathcal D_{n,d_n,s}}
\sup_{P\in\mathcal A_{d_n,s}}\mathcal V_n(\pi,P)
\longrightarrow0
\quad\Longrightarrow\quad
d_n=o(\sqrt n).
\]
This establishes the bounded-density dimension implication under the stated conditions.
% lean: CausalSmith.Experimentation.BivariateSobolevDesignCapacity.published_prognostic_capacity
% lean: CausalSmith.Experimentation.BivariateSobolevDesignCapacity.log_free_bivariate_capacity
\end{enumerate}
\end{proof}

\begin{proof}[Proof of \cref{obj:thm:full-design-lower}]
Fix \(s\in(0,1]\) and integers \(n,d\geq2\). We use the parameters and features of \cref{obj:def:cosine-prior}, written explicitly as
\[
B=\binom d2,\qquad
L=\max\left\{1,\left\lceil\sqrt{4096n/B}\right\rceil\right\},
\qquad
M=BL^2,\qquad
C_0=\sqrt{48+32\pi^2},
\]
\[
\mathcal I_{d,L}
=\{(j,h,k,\ell):1\leq j<h\leq d,\ 1\leq k,\ell\leq L\},
\qquad
V_\alpha(x)=2\cos(\pi kx_j)\cos(\pi\ell x_h),
\qquad
V(x)=(V_\alpha(x))_{\alpha\in\mathcal I_{d,L}}.
\]
Thus \(B>0\), \(L\geq1\), \(M>0\), and \(C_0>0\).

\begin{enumerate}
\item
The constant
\[
c_s=\frac{2}{(48+32\pi^2)\,16384^s}>0
\]
depends on \(s\) alone. We first compare it with the cutoff normalization. Define
\[
\eta_{\mathrm{cut}}=\frac{4096n}{B},
\qquad
\Gamma_{\mathrm{cut}}=\max\{1,n/B\},
\qquad
\Lambda_{\mathrm{cut}}=16384^s.
\]
Since \(\eta_{\mathrm{cut}}\leq4096\Gamma_{\mathrm{cut}}\) and
\(\Gamma_{\mathrm{cut}}\geq1\), the ceiling bound gives
\[
\sqrt{\eta_{\mathrm{cut}}}\leq64\sqrt{\Gamma_{\mathrm{cut}}},
\qquad
\left\lceil\sqrt{\eta_{\mathrm{cut}}}\right\rceil
\leq\sqrt{\eta_{\mathrm{cut}}}+1,
\]
and hence
\[
L\leq64\sqrt{\Gamma_{\mathrm{cut}}}+1
\leq128\sqrt{\Gamma_{\mathrm{cut}}},
\qquad
L^2\leq16384\Gamma_{\mathrm{cut}}.
\]
Raising the last inequality to the positive power \(s\) yields
\[
L^{2s}=(L^2)^s
\leq\Lambda_{\mathrm{cut}}\Gamma_{\mathrm{cut}}^s.
\]

We next verify the scale comparison in its two cases. If \(n/B\leq1\), then
\[
b_s(n,d)\Gamma_{\mathrm{cut}}^s=b_s(n,d)\leq1.
\]
If \(n/B>1\), then
\[
\begin{aligned}
b_s(n,d)\Gamma_{\mathrm{cut}}^s
&\leq (B/n)^s(n/B)^s\\
&=\bigl((B/n)(n/B)\bigr)^s=1.
\end{aligned}
\]
Because \(b_s(n,d)\geq0\) and \(\Lambda_{\mathrm{cut}}>0\), these inequalities imply
\[
L^{2s}b_s(n,d)
\leq\Lambda_{\mathrm{cut}}
       b_s(n,d)\Gamma_{\mathrm{cut}}^s
\leq\Lambda_{\mathrm{cut}}.
\]
Using \(C_0^2=48+32\pi^2\) and positive denominators, we obtain
\[
c_sb_s(n,d)
=\frac{2b_s(n,d)}{C_0^2\Lambda_{\mathrm{cut}}}
\leq\frac{2}{C_0^2L^{2s}}.
\]
% lean: CausalSmith.Experimentation.BivariateSobolevDesignCapacity.priorCutoff_sq_upper
% lean: CausalSmith.Experimentation.BivariateSobolevDesignCapacity.priorCutoff_normalized_lower

\item
Fix an admissible design \(\pi\in\mathcal D_{n,d,s}\). With \(X\) drawn as in
\cref{obj:ass:uniform-draw} and \(Z\) drawn conditionally from \(\pi(X)\), write the scaled signing loss appearing in \cref{obj:def:criterion} as
\[
\mathcal L_{n,d,s}(\pi,m)
=\frac4n\mathbb E_{X,Z}
 \left[\left(\sum_{i=1}^n Z_i m(X_i)\right)^2\right].
\]
The coefficient signs have the independent uniform law of
\cref{obj:def:cosine-prior}, independently of this experiment. Their second moments satisfy
\[
\mathbb E_\xi[\xi_\alpha\xi_\beta]
=
\begin{cases}
1,&\alpha=\beta,\\
0,&\alpha\ne\beta.
\end{cases}
\]
For the first case, \(\xi_\alpha^2=1\). For the second, flipping the sign at
\(\alpha\) preserves the uniform coefficient law and reverses the product
\(\xi_\alpha\xi_\beta\), so its mean is zero.

For every deterministic sample
\(x=(x_1,\ldots,x_n)\in([0,1]^d)^n\) and signing
\(z\in\{-1,1\}^n\), the prior formula gives
\[
\sum_{i=1}^n z_i m_\xi(x_i)
=\frac{1}{C_0L^s\sqrt M}
 \sum_{\alpha\in\mathcal I_{d,L}}
 \xi_\alpha\sum_{i=1}^n z_iV_\alpha(x_i).
\]
Expanding the square and applying the second-moment identity therefore gives
\[
\begin{aligned}
\mathbb E_\xi
 \left[\left(\sum_{i=1}^n z_i m_\xi(x_i)\right)^2\right]
&=\frac{1}{C_0^2L^{2s}M}
 \sum_{\alpha,\beta\in\mathcal I_{d,L}}
 \mathbb E_\xi[\xi_\alpha\xi_\beta]
 \left(\sum_i z_iV_\alpha(x_i)\right)
 \left(\sum_i z_iV_\beta(x_i)\right)\\
&=\frac{1}{C_0^2L^{2s}M}
 \sum_{\alpha\in\mathcal I_{d,L}}
 \left(\sum_i z_iV_\alpha(x_i)\right)^2\\
&=\frac{1}{C_0^2L^{2s}M}
 \left\|\sum_i z_iV(x_i)\right\|_2^2.
\end{aligned}
\]
Finite coefficient averaging commutes with the nonnegative sample and assignment integrals. Consequently,
\[
\mathbb E_\xi\mathcal L_{n,d,s}(\pi,m_\xi)
=\frac{4}{nC_0^2L^{2s}M}
 \mathbb E_{X,Z}
 \left\|\sum_i Z_iV(X_i)\right\|_2^2.
\]
For each fixed sample \(x\), the probability kernel satisfies
\[
\int_{\{-1,1\}^n}
 \left\|\sum_i z_iV(x_i)\right\|_2^2\,\pi(x,dz)
\geq
\min_{z\in\{-1,1\}^n}
 \left\|\sum_i z_iV(x_i)\right\|_2^2.
\]
Integrating this pointwise comparison yields
\[
\mathbb E_\xi\mathcal L_{n,d,s}(\pi,m_\xi)
\geq
\frac{4}{nC_0^2L^{2s}M}
\mathbb E_X
 \min_{z\in\{-1,1\}^n}
 \left\|\sum_i z_iV(X_i)\right\|_2^2.
\]
% lean: CausalSmith.Experimentation.BivariateSobolevDesignCapacity.coefficientSigns_product_integral
% lean: CausalSmith.Experimentation.BivariateSobolevDesignCapacity.coefficientSigns_sum_second_moment
% lean: CausalSmith.Experimentation.BivariateSobolevDesignCapacity.mXi_signed_sum_prior_second_moment
% lean: CausalSmith.Experimentation.BivariateSobolevDesignCapacity.priorRisk_eq_feature_energy
% lean: CausalSmith.Experimentation.BivariateSobolevDesignCapacity.priorRisk_ge_min_feature_energy

\item
The cutoff also supplies the feature-dimension threshold. Indeed,
\[
\sqrt{4096n/B}
\leq\left\lceil\sqrt{4096n/B}\right\rceil
\leq L.
\]
Squaring these nonnegative quantities and multiplying by \(B>0\) gives
\[
\frac{4096n}{B}\leq L^2,
\qquad
4096n\leq BL^2=M.
\]
Under \cref{obj:ass:uniform-draw}, the feature rows \(V(X_i)\) are independent copies. Their coordinates are measurable and square-integrable, and
\cref{obj:lem:cosine-prior} supplies
\[
\mathbb E\!\left[V(X_1)V(X_1)^{\mathsf T}\right]=I_M.
\]
The expected absolute-supremum contraction inequality of
\citep[Theorem~12(4), with Rademacher complexity as in Definition~2]{BartlettMendelson2002Complexities}
supplies the contraction step underlying \cref{obj:lem:isotropic-row-signing}, for its finite classes and finite-dimensional Euclidean unit-ball linear class. Applying that signing bound with feature dimension \(M\geq4096n\) gives
\[
\mathbb E_X
 \min_{z\in\{-1,1\}^n}
 \left\|\sum_i z_iV(X_i)\right\|_2^2
\geq\frac{nM}{2}.
\]
The finite minimum is measurable, so this is the same sample integral that appears in the preceding step.
% lean: CausalSmith.Experimentation.BivariateSobolevDesignCapacity.priorCutoff_dimension_lower
% lean: CausalSmith.Experimentation.BivariateSobolevDesignCapacity.prior_min_feature_energy_lower

\item
Substituting the feature-energy bound into the prior-risk comparison and cancelling the positive factors gives
\[
\begin{aligned}
\mathbb E_\xi\mathcal L_{n,d,s}(\pi,m_\xi)
&\geq
\frac4n\left(\frac{1}{C_0L^s\sqrt M}\right)^2
\frac{nM}{2}\\
&=\frac{2}{C_0^2L^{2s}}\\
&\geq c_sb_s(n,d).
\end{aligned}
\]
The final inequality uses the cutoff comparison \(c_sb_s(n,d)\leq 2/(C_0^2L^{2s})\). This proves the asserted prior lower bound for every admissible design.
% lean: CausalSmith.Experimentation.BivariateSobolevDesignCapacity.priorRisk_ge_normalized_cutoff
% lean: CausalSmith.Experimentation.BivariateSobolevDesignCapacity.full_design_lower

\item
Since \(d\geq2\), \(L\geq1\), and \(0<s\leq1\),
\cref{obj:lem:cosine-prior} ensures that every prior realization belongs to
\(\mathcal H_{d,s}\). The coefficient index set has cardinality \(M\), so its uniform law has total mass
\[
\sum_{\xi\in\{-1,1\}^{\mathcal I_{d,L}}}2^{-M}=1.
\]
For each fixed admissible \(\pi\), every prior atom satisfies
\[
\mathcal L_{n,d,s}(\pi,m_\xi)
\leq
\sup_{m\in\mathcal H_{d,s}}\mathcal L_{n,d,s}(\pi,m).
\]
Multiplying by the prior weights and summing therefore yields
\[
\begin{aligned}
\mathbb E_\xi\mathcal L_{n,d,s}(\pi,m_\xi)
&=\sum_{\xi\in\{-1,1\}^{\mathcal I_{d,L}}}
 2^{-M}\mathcal L_{n,d,s}(\pi,m_\xi)\\
&\leq
\left(\sum_{\xi\in\{-1,1\}^{\mathcal I_{d,L}}}2^{-M}\right)
\sup_{m\in\mathcal H_{d,s}}\mathcal L_{n,d,s}(\pi,m)\\
&=\sup_{m\in\mathcal H_{d,s}}\mathcal L_{n,d,s}(\pi,m).
\end{aligned}
\]
Thus every admissible design has worst-case loss at least \(c_sb_s(n,d)\). Taking the infimum in \cref{obj:def:criterion} gives
\[
R_s(n,d)
=\inf_{\pi\in\mathcal D_{n,d,s}}
 \sup_{m\in\mathcal H_{d,s}}\mathcal L_{n,d,s}(\pi,m)
\geq c_sb_s(n,d).
\]
% lean: CausalSmith.Experimentation.BivariateSobolevDesignCapacity.priorRisk_le_worstLoss_of_legal
% lean: CausalSmith.Experimentation.BivariateSobolevDesignCapacity.capacity_lower_of_priorRisk
\end{enumerate}
\end{proof}

\begin{proof}[Proof of \cref{obj:thm:capacity-answer}]
We apply \cref{obj:thm:log-free-bivariate-capacity}, using the contraction inequality of \citep[Theorem 12(4), with the complexity convention of Definition 2]{BartlettMendelson2002Complexities} and the excess-variance identity of \citep[Proposition 2.3, equation (2.4), under Assumption 2.1 and Definition 2.2]{Cytrynbaum2026Limits}. These published conclusions supply, respectively, the stated measurable-class contraction inequality and the HT excess identity for fair, conditionally outcome-independent assignment under bounded strictly positive covariate density and finite outcome second moments.

\begin{enumerate}
\item Fix integers \(n,d\geq2\). The construction conclusion of \cref{obj:thm:log-free-bivariate-capacity} supplies a single kernel
\[
\pi^\star_{n,d}\in\mathcal D_{n,d,s}
\]
independent of \(s\). For \(n\leq B\), it assigns independent fair signs; for \(B<n\), it is the prescribed spectral-rounding law with independent reflection coins, an independent common Haar shift, and independent uniform rounding seeds. The procedure starts from zero and uses the first \(\lfloor n/4\rfloor\) prescribed real spectral rows evaluated at the reflected and shifted sample. Its deterministic termination conclusion gives
\[
u_i=Z_i,\qquad i=1,\ldots,n,
\]
after \(n\) iterations, for every original sample, every realization of the reflection coins and common shift, and every rounding-seed sequence in \([0,1]\). This supplies the asserted construction and admissibility simultaneously for all smoothness parameters.
% lean: CausalSmith.Experimentation.BivariateSobolevDesignCapacity.log_free_bivariate_capacity

\item Fix \(0<s\leq1\). The comparison conclusions of the same result give
\[
b_s(n,d)=a_s(n,d)=\min\{1,(B/n)^s\},
\qquad
0<c_1\leq c_s,
\]
and
\[
c_sa_s(n,d)
\leq R_s(n,d)
\leq
\sup_{m\in\mathcal H_{d,s}}
\mathcal L_{n,d,s}(\pi^\star_{n,d},m)
\leq3072a_s(n,d),
\]
with exactly the constants in the statement. Its prior conclusions, with the cutoff and normalization of \cref{obj:def:cosine-prior}, give, for every coefficient-sign realization,
\[
\int_{[0,1]^d}m_\xi(x)\,dx=0,
\qquad
m_\xi\in\mathcal H_{d,s}.
\]
For coefficient signs drawn independently of the covariate sample, they also give
\[
c_sa_s(n,d)
\leq
\mathbb E_\xi\!\left[\mathcal L_{n,d,s}(\pi,m_\xi)\right],
\qquad
\pi\in\mathcal D_{n,d,s}.
\]
The kernel selected in the first step remains the same throughout these conclusions.
% lean: CausalSmith.Experimentation.BivariateSobolevDesignCapacity.log_free_bivariate_capacity

\item The Gaussian variance-transfer conclusion of \cref{obj:thm:log-free-bivariate-capacity} gives, for every admissible \(\pi\) and every \(m\in\mathcal H_{d,s}\),
\[
\mathcal L_{n,d,s}(\pi,m)<\infty,
\qquad
n\operatorname{Var}(\widehat\theta)-4.01
=\mathcal L_{n,d,s}(\pi,m).
\]
Taking the supremum over Gaussian completions therefore gives the worst signing loss for each such design; taking the infimum over admissible designs gives \(R_s(n,d)\).

For the frozen uniform outcome-law class, \cref{obj:lem:published-prognostic-capacity} supplies
\[
n\operatorname{Var}_{P,\pi}(\widehat T_P)-\mathsf V(P)
=\mathcal L_{n,d,s}(\pi,m_P)\geq0,
\qquad P\in\mathcal Q_{d,s}.
\]
Thus the positive part equals the excess itself. The published-law capacity conclusion of \cref{obj:thm:log-free-bivariate-capacity} consequently yields
\[
\inf_{\pi\in\mathcal D_{n,d,s}}
\sup_{P\in\mathcal Q_{d,s}}
\max\!\left\{
n\operatorname{Var}_{P,\pi}(\widehat T_P)-\mathsf V(P),0
\right\}
=R_s(n,d).
\]
% lean: CausalSmith.Experimentation.BivariateSobolevDesignCapacity.capacity_answer

\item Finally, the dimension-sequence conclusion of \cref{obj:thm:log-free-bivariate-capacity} applies to every fixed \(0<s\leq1\) and every integer sequence \(d_n\geq2\), and gives
\[
R_s(n,d_n)\longrightarrow0
\quad\Longleftrightarrow\quad
\frac{\binom{d_n}{2}}n\longrightarrow0
\quad\Longleftrightarrow\quad
d_n=o(\sqrt n).
\]
This is the asserted asymptotic characterization.
% lean: CausalSmith.Experimentation.BivariateSobolevDesignCapacity.log_free_bivariate_capacity
\end{enumerate}
\end{proof}

\bibliographystyle{plainnat}

\bibliography{references}

\end{document}
